% Options for packages loaded elsewhere
\PassOptionsToPackage{unicode}{hyperref}
\PassOptionsToPackage{hyphens}{url}
\PassOptionsToPackage{dvipsnames,svgnames,x11names}{xcolor}
\documentclass[
  english,
  11pt,
  reqno]{amsart}
\usepackage{xcolor}
\usepackage{amsmath,amssymb}
\setcounter{secnumdepth}{-\maxdimen} % remove section numbering
\usepackage{iftex}
\ifPDFTeX
  \usepackage[T1]{fontenc}
  \usepackage[utf8]{inputenc}
  \usepackage{textcomp} % provide euro and other symbols
\else % if luatex or xetex
  \usepackage{unicode-math} % this also loads fontspec
  \defaultfontfeatures{Scale=MatchLowercase}
  \defaultfontfeatures[\rmfamily]{Ligatures=TeX,Scale=1}
\fi
\usepackage{lmodern}
\ifPDFTeX\else
  % xetex/luatex font selection
\fi
% Use upquote if available, for straight quotes in verbatim environments
\IfFileExists{upquote.sty}{\usepackage{upquote}}{}
\IfFileExists{microtype.sty}{% use microtype if available
  \usepackage[]{microtype}
  \UseMicrotypeSet[protrusion]{basicmath} % disable protrusion for tt fonts
}{}
\usepackage{longtable,booktabs,array}
\newcounter{none} % for unnumbered tables
\usepackage{calc} % for calculating minipage widths
% Correct order of tables after \paragraph or \subparagraph
\usepackage{etoolbox}
\makeatletter
\patchcmd\longtable{\par}{\if@noskipsec\mbox{}\fi\par}{}{}
\makeatother
% Allow footnotes in longtable head/foot
\IfFileExists{footnotehyper.sty}{\usepackage{footnotehyper}}{\usepackage{footnote}}
\makesavenoteenv{longtable}
\ifLuaTeX
\usepackage[bidi=basic,shorthands=off]{babel}
\else
\usepackage[bidi=default,shorthands=off]{babel}
\fi
\ifLuaTeX
  \usepackage{selnolig} % disable illegal ligatures
\fi
\setlength{\emergencystretch}{3em} % prevent overfull lines
\providecommand{\tightlist}{%
  \setlength{\itemsep}{0pt}\setlength{\parskip}{0pt}}
% Portable additions for the arXiv-oriented profile.
% Keep packages within the standard TeX Live distribution.
\usepackage{mathtools}
\usepackage{amssymb}
\usepackage{amsthm}
\usepackage{booktabs}
\usepackage{listings}
\usepackage{microtype}
\usepackage{centernot}
\usepackage{newunicodechar}

\ifPDFTeX
  % unicode-math supplies mathscr on Xe/LuaTeX; pdfTeX needs its own font.
  \usepackage{mathrsfs}
  % Explicit, meaning-preserving mappings only. Other unsupported input errors
  % are fatal; never discard characters or guess a different math expression.
  \newunicodechar{α}{\ensuremath{\alpha}}
  \newunicodechar{β}{\ensuremath{\beta}}
  \newunicodechar{γ}{\ensuremath{\gamma}}
  \newunicodechar{δ}{\ensuremath{\delta}}
  \newunicodechar{ε}{\ensuremath{\varepsilon}}
  \newunicodechar{θ}{\ensuremath{\theta}}
  \newunicodechar{λ}{\ensuremath{\lambda}}
  \newunicodechar{μ}{\ensuremath{\mu}}
  \newunicodechar{π}{\ensuremath{\pi}}
  \newunicodechar{σ}{\ensuremath{\sigma}}
  \newunicodechar{ω}{\ensuremath{\omega}}
  \newunicodechar{ℂ}{\ensuremath{\mathbb{C}}}
  \newunicodechar{ℕ}{\ensuremath{\mathbb{N}}}
  \newunicodechar{ℚ}{\ensuremath{\mathbb{Q}}}
  \newunicodechar{ℝ}{\ensuremath{\mathbb{R}}}
  \newunicodechar{ℤ}{\ensuremath{\mathbb{Z}}}
  \newunicodechar{∞}{\ensuremath{\infty}}
  \newunicodechar{∈}{\ensuremath{\in}}
  \newunicodechar{∅}{\ensuremath{\emptyset}}
  \newunicodechar{⊂}{\ensuremath{\subset}}
  \newunicodechar{⊆}{\ensuremath{\subseteq}}
  \newunicodechar{∩}{\ensuremath{\cap}}
  \newunicodechar{∪}{\ensuremath{\cup}}
  \newunicodechar{≤}{\ensuremath{\leq}}
  \newunicodechar{≥}{\ensuremath{\geq}}
  \newunicodechar{≠}{\ensuremath{\neq}}
  \newunicodechar{→}{\ensuremath{\rightarrow}}
  \newunicodechar{↦}{\ensuremath{\mapsto}}
  \newunicodechar{⇒}{\ensuremath{\Rightarrow}}
  \newunicodechar{⇔}{\ensuremath{\Leftrightarrow}}
  \newunicodechar{−}{\ensuremath{-}}
\fi
\newunicodechar{□}{\ensuremath{\square}}

% Portable negated hooked arrow used by several manuscripts.
\providecommand{\nhookrightarrow}{\centernot\hookrightarrow}

% U+220E END OF PROOF: render a small hollow square at the right margin.
\newcommand{\rightqedsquare}{%
  \ifhmode
    \unskip\nobreak\hfill\ensuremath{\square}%
  \else
    \noindent\hfill\ensuremath{\square}%
  \fi
}
\newunicodechar{∎}{\rightqedsquare}

% Reduce avoidable overfull lines without changing the source manuscript.
\setlength{\emergencystretch}{2em}
% Match amsart's theorem preskip while keeping Markdown statement bodies roman.
\makeatletter
\newcommand{\pandocstatementheading}[1]{%
  \par\addpenalty{\@beginparpenalty}%
  \addvspace{.5\baselineskip plus .2\baselineskip minus .2\baselineskip}%
  {\normalfont\bfseries #1\par}%
  \nobreak\@afterindenttrue\@afterheading
}
\makeatother
% House style for standard/arxiv-ams: match the References section heading.
% The native abstract box, margins and body font remain owned by amsart.
\usepackage{etoolbox}
\newcommand{\pandocabstractheading}{%
  \par\begingroup
    \normalsize\section*{\abstractname}%
  \endgroup
  \nobreak\noindent\ignorespaces
}
\makeatletter
\patchcmd{\abstract}
  {\item[\hskip\labelsep\scshape\abstractname.]}
  {\item[]\pandocabstractheading}
  {}%
  {\PackageError{studio-ams-abstract}
     {Unsupported amsart abstract definition}
     {Check the installed class before updating the abstract-heading patch.}}
\makeatother
% Keep this minimal and portable across engines.

% Math shortcuts
\newcommand{\N}{\mathbb{N}}
\newcommand{\Z}{\mathbb{Z}}
\newcommand{\Q}{\mathbb{Q}}
\newcommand{\R}{\mathbb{R}}
\newcommand{\C}{\mathbb{C}}
\newcommand{\eps}{\varepsilon}

% Theorem-like environments (lightweight)
% If you need full theoremstyles, consider amsthm or thmtools here.
% \usepackage{amsthm}   % already available via TeX Live if needed
% \newtheorem{theorem}{Theorem}
% \newtheorem{lemma}{Lemma}

% Code listings default (avoid minted to keep shell-escape off)
\lstset{basicstyle=\ttfamily\small,breaklines=true,columns=fullflexible}

\usepackage{bookmark}
\IfFileExists{xurl.sty}{\usepackage{xurl}}{} % add URL line breaks if available
\urlstyle{same}
\providecommand{\xmpquote}[1]{#1}
\hypersetup{
  pdftitle={Global twistor sources for analytic Einstein spacetimes},
  pdfauthor={Lior Isthmus},
  pdflang={en},
  pdfkeywords={\xmpquote{gravitational googly
problem}, \xmpquote{twistor correspondence}, \xmpquote{analytic Einstein
metrics}, \xmpquote{null monodromy}, \xmpquote{microdifferential
algebroids}, \xmpquote{Cauchy development}},
  colorlinks=true,
  linkcolor={blue},
  filecolor={Maroon},
  citecolor={blue},
  urlcolor={blue},
  pdfcreator={LaTeX via pandoc}}

\title{Global twistor sources for analytic Einstein spacetimes}
\author{Lior Isthmus\textsuperscript{1}}
\date{2026-11-13}

\begin{document}
\maketitle

\markboth{\MakeUppercase{Lior Isthmus}}{\MakeUppercase{Global twistor sources for analytic Einstein spacetimes}}

\protect\phantomsection\label{sec:title}

\label{global-twistor-sources-for-analytic-einstein-spacetimes}

\protect\phantomsection\label{sec:abstract}

\pandocabstractheading\label{abstract}

We globalize the local analytic projective twistor correspondence for
the gravitational googly problem, retaining its convergent nonlinear
treatment of both gravitational chiralities. For each fixed Einstein
coefficient, the correspondence identifies independently defined atlases
of two analytic Čech classes of weights \(2\) and \(-6\), structured
operator sources on null monodromy presentations, and holomorphic
oriented spin Einstein four-manifolds. Its inverse comparisons preserve
the original source operations, Einstein scale, spin structure and
complete point family, together with calibrated modules on parametrized
incidence. For analytic Lorentz developments with a compact marked
Cauchy surface, germs along the entire real development give the
corresponding equivalence, and the maximal analytic Cauchy development
determines a terminal source germ. A Nariai example admits no Hausdorff
contact carrier containing the full complex skies over all points of the
real development. In its product complexification and every allowed
point restriction containing the entire real development, every real
point sky has nonclosed saturation, so no ordinary locally finitely
generated module has exactly that saturated support. These obstructions
explain the use of null monodromy and parametrized incidence in the
global correspondence, with compact point objects recovered by central
specialization.

\footnotetext[1]{Email: \href{mailto:lioristhmus@proton.me}{\nolinkurl{lioristhmus@proton.me}}}
\setcounter{footnote}{1}

\clearpage
\tableofcontents

\protect\phantomsection\label{sec:introduction}

\section{1. Introduction and main
results}\label{introduction-and-main-results}

Penrose's gravitational googly problem asks how to incorporate the
interacting chirality absent from the nonlinear-graviton construction
within the same twistor conventions (Penrose 1976; 2015, secs. 4--5).
The local analytic correspondence in Isthmus (2026a), hereafter the
\emph{local paper}, treats the nonlinear coupling of both gravitational
chiralities through two analytic Čech classes of weights \(2\) and
\(-6\) on a fixed flat twistor germ. It constructs every marked local
analytic Einstein four-germ with a prescribed Einstein coefficient,
proves convergence of its prescribed nonlinear construction, and
establishes inverse correspondences on independently admitted geometric
and operator domains. The present paper globalizes this local analytic
solution over an entire point manifold and, with Lorentz reality, over
an entire Cauchy development.

Several earlier constructions address the geometry beyond self-duality.
LeBrun reconstructs complex conformal geometry locally from compact
quadrics in a space of complex null geodesics and gives an ambitwistor
recognition of conformally Einstein metrics (LeBrun 1983, 1985). The
thickening approach of Baston and Mason has a fully curved extension due
to LeBrun, with a conformal Einstein interpretation when the Weyl
curvature is algebraically general (Baston and Mason 1987; LeBrun 1991).
Sharma obtains a twistor action for Euclidean general relativity with
zero cosmological constant, using partially integrable almost complex
structures, and shows that this action reduces to the Plebański action
(Sharma 2021). More recent approaches include Penrose's bi-twistors,
which replace the role of \(\alpha\)-planes by null geodesics, and the
construction of a metric conformal to Schwarzschild from holomorphic
data in the twistor space of self-dual Taub--NUT (Penrose 2026; Adamo et
al. 2026). Our starting point is the specific convergent two-class
correspondence in Isthmus (2026a); the task here is to continue its
source data, operations and inverse comparisons globally.

The operator formulation is also related to Penrose's palatial proposal,
which discusses local twistor algebras, their patching and point
recovery in connection with analytic Lorentz globally hyperbolic
geometry (Penrose 2015, secs. 5--6). Marcolli and Penrose extend
Donaldson--Friedman gluing in the Riemannian self-dual setting to
several quantized twistor constructions by deformation theory (Marcolli
and Penrose 2021). Here the global objects are independently admitted
source atlases and their operator realizations on the null monodromy
presentation of an entire point manifold. Their inverse correspondence
preserves the original source classes and physical transition maps, as
well as the supplied operator associators and framed point modules.

Global twistor geometry has established precedents in other signatures
and settings. LeBrun and Mason construct a global correspondence for
self-dual conformal structures of split signature near the standard
structure on \(S^2\times S^2\), using holomorphic disks (LeBrun and
Mason 2007). In Lorentz geometry, reconstruction from real light rays
and skies has its own causal and sky-separation hypotheses (Bautista et
al. 2015). Non-Hausdorff twistor spaces also occur in the Geroch-group
construction of Woodhouse and Mason (Woodhouse and Mason 1988). The
obstruction proved below concerns a specific stronger requirement: a
Hausdorff complex carrier for the full complex null quadrics over every
point of an entire real Einstein development.

A global source is an atlas of the original input classes with their
nonlinear physical comparison law. Its point domains and comparison maps
specify the topology, continuation domain and complete point family.
They determine the complex null phase and its monodromy groupoid, which
remembers null paths modulo endpoint-fixed homotopy. Structured operator
algebroids descend on contact transversals of this groupoid. The point
modules descend on Hausdorff parametrized incidence, retaining the
ambient relative operator germs and their central point parameters; each
point object is recovered by central specialization. This construction
also accommodates the nonclosed point-sky saturations in the Nariai
example.

The external local results are collected in
\hyperref[sec:local-input-v08]{§2}, with precise references to the
posted version v1 of the local paper. The new contributions are the
global construction and independent inverse correspondence, their
compatibility with spin structures and Lorentz reality, and the
obstruction to stronger ordinary realizations. The null-geodesic
incidence construction and structured local contact quantization are
used with their established geometric and analytic conventions (Isthmus
2026a, Propositions 4.2--4.3; Kashiwara 1996, Theorem 2; Polesello 2005,
Theorem 3.3).

\pandocstatementheading{\protect\phantomsection\label{thm:global-source-correspondence-v08}Theorem 1.1 (Global source correspondence).}

Fix \(\Lambda\in\mathbb C\). Use the local correspondence and analytic
conventions of \hyperref[sec:local-input-v08]{§2}. There are
equivalences
\begin{equation}\protect\phantomsection\label{eq:focused-holomorphic-correspondence-v08}{
\mathsf{SrcAtl}^{\mathrm{hol}}_\Lambda
\simeq
\mathsf{OpMon}^{\mathrm{hol}}_\Lambda
\simeq
\mathsf{MonRay}^{\mathrm{hol}}_\Lambda
\simeq
\mathsf{Ein}^{\mathrm{hol}}_\Lambda.
}\end{equation} Here the source atlas consists of original two-class
inputs with physical comparison germs as in
\hyperref[def:independent-original-source-atlas]{Definition 5.1}; the
geometric and operator monodromy domains are independently admitted by
their incidence, transport, optical and operator conditions. Metric
objects are connected holomorphic oriented spin Einstein four-manifolds
with coefficient \(\Lambda\). Morphisms retain the specified spin lifts
of point embeddings; operator morphisms are physical Hom classes modulo
the stated normalized comparison isomorphisms.

Both inverse comparisons retain the entire chosen point manifold, its
proper full quadric family, spin structure and Einstein scale. They
preserve the original local source operations and prescribed analytic
representatives, the projective spin, optical and infinity diagram, and
the calibrated module family on parametrized incidence. The source unit
recovers the original input classes and overlap instructions; the
operator counit preserves the supplied associators and framed module
comparisons. For the
\hyperref[def:conditioned-holomorphic-source-families]{coherent
holomorphic families specified in §5}, these comparisons are holomorphic
on conditioned local parents and preserve finite parameter jets on each
retained finite diagram. The equivalence commutes with point restriction
when monodromy is recomputed inside the restricted point domain.

\emph{Proof.}
\hyperref[thm:original-source-atlas-operator-equivalence]{Theorem 5.6}
constructs the first equivalence and its inverse comparisons.
\hyperref[thm:operator-monodromy-equivalence]{Theorem 4.6} and
\hyperref[thm:independent-monodromy-equivalence]{Theorem 3.11} give the
remaining equivalences.
\hyperref[prop:source-atlas-complete-diagram]{Proposition 5.7},
\hyperref[thm:global-parametrized-point-module-family]{Theorem 4.8} and
\hyperref[sec:spin-reality-v08]{§6} identify the preserved structures.
The conditioned family assertions follow from
\hyperref[prop:source-atlas-local-analytic-execution]{Proposition 5.4}
and the finite-jet argument of
\hyperref[prop:source-atlas-complete-diagram]{Proposition 5.7}.
\hyperref[lem:source-atlas-phase-completion]{Lemma 5.5} and the
restriction comparisons of
\hyperref[sec:global-null-monodromy-v08]{§§3--4}, in particular
\hyperref[cor:null-monodromy-first-unit-v08]{Corollary 4.3}, prove
compatibility with holomorphic point restrictions.
\hyperref[prop:recomputed-monodromy-restriction]{Proposition 7.1} gives
its form for germs along the entire real development. ∎

Hausdorff contact carriers do not suffice in general for entire real
Einstein developments. The following result concerns the complete real
Nariai spacetime
\begin{equation}\protect\phantomsection\label{eq:focused-nariai-metric-v08}{
M_a=\mathbb R\times S^1\times S^2,
\qquad
g_a=\mathrm dt^2-a^2\cosh^2(t/a)\,\mathrm d\theta^2-a^2g_{S^2},
\qquad a>0,
}\end{equation} where \(g_{S^2}\) is the unit round metric and
\(\theta\) has period \(2\pi\). Its Einstein coefficient is \(-a^{-2}\),
and its Weyl tensor is nonzero.

\pandocstatementheading{\protect\phantomsection\label{thm:global-realization-obstructions-v09}Theorem 1.2 (Obstructions to ordinary realizations).}

For the entire Nariai development \((M_a,g_a)\), the following
assertions hold.

\begin{enumerate}
\tightlist
\item
  In any holomorphic Einstein complexification, no neighborhood of the
  full complex null quadrics over every real point admits a holomorphic
  map to a Hausdorff complex manifold whose differential has exactly the
  null-geodesic line as kernel. In particular, a global Hausdorff
  contact carrier with those whole skies does not exist.
\item
  In the explicit product complexification used in
  \hyperref[sec:nariai-obstructions-v08]{§8}, every real point sky has
  nonclosed saturation after every allowed point restriction (defined in
  \hyperref[sec:germ-correspondence-v08]{§7}) retaining all of \(M_a\).
  No ordinary locally finitely generated operator module on the
  resulting null monodromy presentation has exactly this saturated
  support. Every other compatible holomorphic Einstein complexification
  has an allowed point neighborhood on which this support obstruction
  holds for every real point and persists under every further allowed
  restriction.
\end{enumerate}

The Hausdorff-carrier obstruction holds on every complexification. The
support obstruction holds on every allowed product neighborhood and, for
any other complexification, on a cofinal collection of representatives
of its germ along the entire real development.

\emph{Proof.}
\hyperref[thm:nariai-no-hausdorff-whole-sky-carrier]{Theorem 8.2} proves
the first assertion.
\hyperref[thm:nariai-all-real-nonclosed-point-skies-v09]{Theorem 8.6},
\hyperref[lem:finite-module-closed-support]{Lemma 8.7} and
\hyperref[cor:nariai-cofinal-point-module-obstruction-v09]{Corollary
8.8} prove the second. ∎

The module in Theorem 1.1 is a module over the inverse-image algebroid
of ambient relative operator germs on incidence. Its support there is
the full incidence manifold; its compact point fibers are central
specializations. Thus it has the module action required by the local
correspondence while respecting the support obstruction above. The
theorem concerns the complete point families supplied in source
admission, not a classification of all compact objects of a bare
operator algebroid as spacetime points.

\pandocstatementheading{\protect\phantomsection\label{cor:focused-maximal-cauchy-source-v08}Corollary 1.3 (Maximal analytic Cauchy sources).}

Fix real \(\Lambda\) and analytic Einstein constraint data
\((\Sigma,h,K,P_\Sigma)\) with compact connected \(\Sigma\). Impose the
compatible Lorentz reality, exact real point locus and initial markings
defined in \hyperref[sec:germ-correspondence-v08]{§7}. After passing to
germs on point neighborhoods of the entire real development, the
correspondence gives
\begin{equation}\protect\phantomsection\label{eq:focused-cauchy-correspondence-v08}{
\mathsf{SrcAtl}^{\mathrm{germ}}_\Lambda(\Sigma,h,K)
\simeq
\mathsf{OpMon}^{\mathrm{germ}}_\Lambda(\Sigma,h,K)
\simeq
\mathsf{Ein}^{\omega}_\Lambda(\Sigma,h,K).
}\end{equation} The maximal analytic spin Cauchy development determines
a terminal source germ. This correspondence is natural under
initial-data isomorphisms with their specified spin lifts and retains
all the structures in Theorem 1.1.

\emph{Proof.} Apply
\hyperref[thm:source-atlas-cauchy-equivalence-v08]{Theorem 7.4} and
\hyperref[cor:maximal-source-atlas-germ-v08]{Corollary 7.5}, using the
analytic Cauchy construction of
\hyperref[sec:analytic-cauchy-appendix-v08]{Appendix A}. ∎

The real point manifold is the entire retained development; the
neighborhoods defining its complex germ may have widths that vary from
point to point. \hyperref[sec:germ-correspondence-v08]{Section 7}
defines point restriction by recomputing null monodromy. The maximal
analytic development is constructed as an open subdevelopment of the
smooth maximal globally hyperbolic development, with the scope stated in
\hyperref[rem:analytic-maximality-scope]{Remark A.9}.

Sections \hyperref[sec:global-null-monodromy-v08]{3} and
\hyperref[sec:operator-incidence-v08]{4} construct the geometric and
operator global sources.
\hyperref[sec:original-source-atlases-v08]{Section 5} proves that they
are obtained from, and recover, atlases of the original input classes.
\hyperref[sec:spin-reality-v08]{Section 6} treats the projective spin
diagram and Lorentz reality; \hyperref[sec:germ-correspondence-v08]{§7}
establishes the Cauchy correspondence.
\hyperref[sec:nariai-obstructions-v08]{Section 8} proves the
obstructions, and \hyperref[sec:analytic-cauchy-appendix-v08]{Appendix
A} supplies the analytic Cauchy facts needed for the maximality
statement.

The computational supplement (Isthmus 2026b) provides universal
metric-jet checks of the Gaussian differential identities, with explicit
metric cross-checks from a separate implementation, as well as exact
symbolic checks of selected Nariai formulas and the tractor and
infinity-pair normalizations.

\protect\phantomsection\label{sec:local-input-v08}

\section{2. Local correspondence and
conventions}\label{local-correspondence-and-conventions}

We use the posted version v1 of the local paper, \emph{A projective
twistor correspondence for local analytic Einstein germs}, Hexagon
2611.00123v1, dated November 5, 2026 (Isthmus 2026a). Propositions
2.1--2.3 recall the forms of its original-source, holomorphic geometric
and operator results needed below; their proofs belong to the local
paper. Lemma 2.4 proves conjugation naturality for the independent
geometric reconstruction. The common analytic convention for these local
inputs is a bounded holomorphic family over a finite-dimensional complex
manifold germ, with positive coordinate and conic buffers, bounded maps
and inverses, and positive lower bounds for the finitely many
coefficients inverted on a specified finite comparison diagram. Finite
parameter jets are restrictions of such holomorphic families to actual
finite Taylor neighborhoods. The global constructions use these bounds
locally on compact subsystems.

Our curvature convention is
\(R(a,b)=[\nabla_a,\nabla_b]-\nabla_{[a,b]}\) and
\(\operatorname{Ric}(b,c)=\operatorname{tr}(a\mapsto R(a,b)c)\). The
Einstein coefficient \(\Lambda\) is defined by
\(\operatorname{Ric}(g)=\Lambda g\). Holomorphic point manifolds have
complex dimension four. Real spacetimes are oriented, time-oriented and
four-dimensional, with signature \((+,-,-,-)\). Thus the de Sitter
metric of radius \(a\) has \(\Lambda=-3a^{-2}\) in these conventions.
The Schouten tensor is
\(P=\frac12(\operatorname{Ric}-\frac16\operatorname{Scal}\,g)\).

All point manifolds and real spacetimes are Hausdorff and second
countable. Groupoid arrow manifolds and explicitly identified contact
leaf spaces may be non-Hausdorff. Point restrictions are taken
componentwise when their domains are disconnected. Spin structures and
specified full spin lifts are retained as data: the nontrivial central
spin transformation remains distinct even when it acts trivially on
tensor records or projective skies. An Einstein embedding is an open
isometric embedding with its specified spin lift. The contact carrier
has complex dimension five. A full complex sky is a compact quadric
\(\mathbb P^1\times\mathbb P^1\) with ordered rulings. Completeness of
the point family means the local four-dimensional completeness specified
in \hyperref[prop:independent-local-ray-recognition-v09]{Proposition
2.2}, together with retention of the whole quadric over every point of
the chosen domain.

\pandocstatementheading{\protect\phantomsection\label{prop:original-two-class-local-input}Proposition 2.1 (Original source correspondence).}

The following local statements are drawn from Definition 2.1, Theorems
3.1, 6.3, 8.3 and 9.10--9.12, and Proposition 8.4 of the local paper.
Their proofs are in the local paper.

Let \[
T_0=\operatorname{Tot}(\mathcal O(1)^{\oplus2})
\longrightarrow\mathbb P(S')
\] be the marked flat seed as a germ along its entire zero line. For
integer \(m\), let \(\mathscr H_m(T_0)\) be the two-chart analytic Čech
group in the direct limit over neighborhoods of that line. With north
coordinates \((w,\zeta)\) and south coordinates \((w/\zeta,1/\zeta)\),
the fixed differential is
\begin{equation}\protect\phantomsection\label{eq:original-source-cech-differential}{
(\delta a)_{NS}(w,\zeta)
=\zeta^m a_S(w/\zeta,1/\zeta)-a_N(w,\zeta).
}\end{equation} At a fixed Einstein coefficient \(\Lambda\in\mathbb C\),
the original input domain is
\begin{equation}\protect\phantomsection\label{eq:original-source-local-domain}{
\mathcal A_\Lambda
=\mathscr H_2(T_0)\oplus\mathscr H_{-6}(T_0).
}\end{equation} The fixed residue and middle-Laurent projections
identify these classes with their unique canonical analytic
representatives and with independent complete curvature records
\(U=(U^+,U^-)\). They give an analytic quotient correspondence,
including analytic coboundaries and common smaller representatives for
bounded holomorphic families. Adding a source coboundary leaves the
canonical pair and its prescribed normalized output unchanged (Isthmus
2026a, Theorem 3.1 and §8.1).

Let \(\mathfrak O_\Lambda\) be the independently admitted local operator
domain recalled in
\hyperref[prop:independent-local-operator-input]{Proposition 2.3}. The
original source constructor \(\mathcal S_{\mathrm{op}}\) and its decoder
\(\mathcal D_{\mathrm{op}}\) satisfy
\begin{equation}\protect\phantomsection\label{eq:original-source-local-operator-inverses}{
\mathcal D_{\mathrm{op}}\mathcal S_{\mathrm{op}}
=1_{\mathcal A_\Lambda},
\qquad
\mathcal S_{\mathrm{op}}\mathcal D_{\mathrm{op}}
\simeq1_{\mathfrak O_\Lambda}.
}\end{equation} The first equality is on the original marked input
classes and their Einstein coefficient. The second is a comparison of
operator presentations, with the supplied and reconstructed associators,
framed module maps, torsor, optical root and point incidence retained.
These comparisons are natural on physical point/full-spin Hom classes.
Both domains contain arbitrary independent chiral records, every
algebraic Weyl type and \(\Lambda=0\) (Isthmus 2026a, Theorems
9.10--9.11).

The physical source action is defined directly on the complete records
by
\begin{equation}\protect\phantomsection\label{eq:original-source-physical-record-action}{
R_{a,A}U
=A^*\!\left(\operatorname{Fl}_{H_a}^{\,1}U\right),
}\end{equation} where \(H_a\) is the original constant-direction Cartan
record field, \(a\) is an admissible short marked displacement and \(A\)
is an actual full spin element. Admissibility requires that the
indicated flow and marked maps lie in their retained parents. If
\(V=R_{a,A}U\), the associated normal geometric map is
\begin{equation}\protect\phantomsection\label{eq:original-source-physical-map-direction}{
\Xi_{a,A}:g[V]\longrightarrow g[U],
\qquad
v\longmapsto\operatorname{Exp}_{x_a}^{g[U]}(\nu_aAv),
}\end{equation} where \(x_a=\operatorname{Exp}_o(a)\) and \(\nu_a\) is
initialized radial spin transport. The covariant source instruction
\(U\to V\) uses \(\Xi_{a,A}^{-1}\). The geometry in this formula is
derived from the source; it is not an additional admission datum.

For successive admissible instructions the law is
\begin{equation}\protect\phantomsection\label{eq:original-source-physical-composition}{
\begin{aligned}
(b,B)\circ(a,A)
&=\bigl(h_U(a,Ab),K_U(a,Ab)AB\bigr),\\
(a,A)^{-1}
&=\bigl(-A^{-1}a,A^{-1}\bigr)
\quad\text{at }V.
\end{aligned}
}\end{equation} Here \(h_U\) is the endpoint displacement read in the
initial marked normal coordinates, and \(K_U\) is the initialized spin
rotation comparing direct radial transport with the corresponding broken
transport. These laws include identities and associativity wherever the
finite diagram is admitted. Every physical local isometry with its
specified full spin lift has a unique such short instruction in the
marked representatives. The contact-affine comparisons obey the
corresponding composition laws, and no independent vertical comparison
remains after the physical map and full spin lift are specified (Isthmus
2026a, Proposition 8.1 and Theorem 8.3). Full spin labels remain
distinct even when their action on tensor records is identical.

The original nonlinear constructor is the fixed finite source execution
and prescribed normalizer specified in Proposition 5.9 and §6 of the
local paper. Its outputs preserve the two original input classes, their
nonlinear source operations and the six scale normalizations through the
natural comparisons in Theorem 9.12 of the local paper. The negative
class determines its dual lift; that lift is not an additional free
input. The six scalar choices fix one value, four gradient components
and the Einstein coefficient of the selected section. Their components
are recomputed in each auxiliary scale and frame, with the Cauchy-state
transformation; preserving them means preserving this selected-section
comparison, rather than equality of raw chart components (Isthmus 2026a,
Theorem 5.5 and Proposition 6.2).

For bounded holomorphic input families with finite-dimensional complex
parameter bases and each fixed finite conditioned comparison diagram,
the constructor, decoder, normalized representatives, their inverses and
comparisons have common smaller analytic parents and the stated
finite-jet dependence. The conditions retain positive coordinate and
conic buffers, bounds for maps and inverses, positive unit margins, and
parents containing every specified finite composite. No uniform estimate
on all global charts or unbounded comparison histories is part of this
input.

The Lorentz operation conjugates the complete records and exchanges
their spin factors and chiralities, with \(\Lambda\) conjugated. The
exact residue dictionaries transfer this operation to the original
classes while retaining their weights \(2\) and \(-6\). At real
\(\Lambda\), its fixed records and full real spin instructions give the
Lorentz restriction. The operator comparison carries the specified weak
square and module coherence of
\hyperref[prop:operator-monodromy-lorentz-reality]{Proposition 6.5}.

The source category retains original cohomology classes, their
prescribed representatives and physical instructions. Its
raw-representative and atlas conventions are specified in
\hyperref[def:independent-original-source-atlas]{Definition 5.1}; the
Hamiltonian and parameter conventions are stated after
\hyperref[prop:source-atlas-local-analytic-execution]{Proposition 5.4}.
Local comparisons use common refinements near the whole retained compact
sky.

\pandocstatementheading{\protect\phantomsection\label{prop:independent-local-ray-recognition-v09}Proposition 2.2 (Independent geometric recognition).}

The independently admitted local objects used below are those of
Definition 2.2 of the local paper. Such an object has a holomorphic
contact fivefold germ with a complete regular four-dimensional family of
embedded Legendrian quadrics \(Q\simeq\mathbb P^1\times\mathbb P^1\),
ordered rulings, and contact line of restriction \(\mathcal O(1,1)\).
Its incidence is \(p:F\to U\), \(q:F\to Z\), with one-dimensional
\(q\)-fibers projecting nontrivially to \(U\). It also has a quadratic
rank-six holomorphic bundle \((\mathcal V,H)\), a section \(\mathsf I\)
with \(H(\mathsf I,\mathsf I)=-\Lambda/3\), and an affine line subbundle
\(\mathscr A\subset\mathcal V\). The bundle is trivial on the central
sky, the translation line restricts to \(\mathcal O(-1,-1)\), and
\begin{equation}\protect\phantomsection\label{eq:independent-ray-affine-calibration}{
H(Y,Y)=0,
\qquad H(\mathsf I,Y)=1
\qquad(Y\in\mathscr A).
}\end{equation} Complete-sky evaluation reconstructs a rank-six bundle
\(\mathcal T\) on \(U\), a connection \(D\), and a section
\(\mathsf X\), with an evaluation isomorphism
\(p^*\mathcal T\simeq q^*\mathcal V\). Put
\(W=(\mathsf I,\mathsf X)^\perp\) and \(\Theta(v)=D_v\mathsf X\).
Admission requires \(\Theta:TU\to W\) to be invertible,
\(F^D\mathsf X=0\), and the Ricci contraction of
\begin{equation}\protect\phantomsection\label{eq:independent-ray-middle-curvature}{
(v,w,u)\longmapsto
\Theta^{-1}\operatorname{pr}_{W}F^D(v,w)\Theta(u)
}\end{equation} to vanish. The connection in these conditions is
reconstructed by inverse null transport; it is not freely supplied.

The local recognition theorem constructs
\(g(v,w)=H(\Theta(v),\Theta(w))\), proves
\(\operatorname{Ric}(g)=\Lambda g\), and identifies the supplied object
with the canonical ray geometry of \(g\). Evaluation intertwines \(D\)
with the constant partial connection along \(q\)-fibers, takes
\(\mathsf X\) to the supplied affine position, and preserves
\(H,\mathsf I\) and the spin marking. Conversely, every holomorphic
oriented spin Einstein four-germ has such a holomorphic whole-sky model,
without restrictions on chirality or algebraic Weyl type. For the
compatible Lorentz real form, Lemma 2.4 below proves reality of the
reconstructed metric; its retained tangent calibration gives signature
\((+,-,-,-)\) and real skies \(S^2\). The converse retains this reality.
Enriched comparison germs are in bijection with holomorphic local
isometries carrying a specified spin lift, compatibly with restriction
and composition; their compatible real restrictions give the
corresponding analytic Lorentz comparisons.

The holomorphic assertions are the results of Theorem 2.4, Propositions
4.1--4.3 and Theorem 8.3 of the local paper, in the forms used here. The
Lorentz converse and real-sky description use the constructions in §8.4
and Proposition I.3 of the local paper. In particular, Proposition 4.3
of the local paper proves recognition on independently supplied objects
and both comparison directions. The hypotheses concern the supplied
incidence, bundle, affine calibration and reconstructed normality
conditions, rather than membership in a generated image.

\pandocstatementheading{\protect\phantomsection\label{prop:independent-local-operator-input}Proposition 2.3 (Independent operator recognition).}

The following local inputs are drawn from Definition H.1, Lemma H.2,
Proposition H.3, Theorem H.4 and Theorem G.5 of the local paper.

An independently admitted operator object has analytic half-density
algebras with their order filtrations, full graded-symbol
identifications and complex-linear formal adjoints, including the
specified square and degree-zero duality coherences. On presentation
overlaps its algebra maps and normalized units obey
\begin{equation}\protect\phantomsection\label{eq:local-operator-algebroid-input}{
\begin{aligned}
\varphi_{ab}\varphi_{bc}
&=\operatorname{Ad}(C_{abc})\varphi_{ac},\\
C_{abc}C_{acd}
&=\varphi_{ab}(C_{bcd})C_{abd},\\
C_{abc}^*C_{abc}&=1,
\qquad \sigma_0(C_{abc})=1.
\end{aligned}
}\end{equation} Its symbols determine a separated contact germ \(Y\)
along a compact characteristic quadric \(Q\), with ordered rulings and
contact line \(\mathcal L|_Q=\mathcal O(1,1)\). The intrinsic first-unit
sheaf is identified with \(\mathcal L^{-1}\), and a torsor for this
sheaf is retained as substantive affine data. The calibrated affine
object satisfies the optical root, projective-extension and
Einstein-scale recognition conditions in Definition 5.6 of the local
paper. A compact locally standard order-zero simple module with reduced
multiplicity-one characteristic \(Q\) determines a complete
four-dimensional nearby point family. Its half-canonical root has its
actual square identification and any retained base factor; its framing
is a nonzero local-system vector along the calibrated cone for the whole
parameter family. It is distinct from the optical density root.

Taking full symbols, the torsor and the characteristic gives an
equivalence with the independently admitted calibrated affine whole-sky
objects. Both inverse comparisons are comparisons of presentations.
Physical Hom sets are taken modulo normalized comparison isomorphisms
and retain the specified full spin lift. Actual structured comparisons
retain both algebra associators and the semilinear framed-module maps;
their defining equations are
(\ref{eq:operator-source-comparison-descent}).

Both associators occur in these equations. Normalization of algebra
comparison units is separate from module framing: the module maps are
semilinear isomorphisms carrying the whole local-system vector, and are
not additionally restricted to principal scalar one. Over a fixed
admitted physical map and actual full spin label, the structured
comparison groupoid is contractible. Its uniqueness includes the framed
module comparison (Isthmus 2026a, Lemma H.2, Theorem H.4, Proposition
G.4 and Theorem G.5). These statements commute with restriction and
common refinement and hold on the conditioned finite diagrams and
bounded local holomorphic families of the cited domain.

These assertions are external inputs, not new quantization or module
existence results. The admissible point modules are the locally standard
framed modules specified above; arbitrary coherent or flat relative
modules are not included. In particular, Theorem H.4 of the local paper
does not turn an arbitrary contact germ into a physical map of complete
point objects.

\pandocstatementheading{\protect\phantomsection\label{lem:independent-ray-lorentz-recognition-v092}Lemma 2.4 (Conjugation and Lorentz recognition).}

Independent whole-sky recognition commutes with compatible
antiholomorphic comparison germs of the incidence and quadratic-affine
data, including their actual spin identifications. In particular, if
such a comparison is a reality involution over \(\sigma:U\to U\) at real
\(\Lambda\), the reconstructed metric satisfies
\begin{equation}\protect\phantomsection\label{eq:independent-ray-metric-reality-v092}{
\sigma^*\overline g=g.
}\end{equation} If the reality has the prescribed Lorentz spin/tangent
calibration at a fixed mark, its real restriction has signature
\((+,-,-,-)\) near that mark and its fixed sky is \(S^2\). Every marked
real analytic Lorentz Einstein germ has a compatible whole-sky model
with this calibration.

\emph{Proof.} Apply the reconstruction in Isthmus (2026a, Proposition
4.3) to a model and its conjugate. On each complete compact sky,
conjugation carries holomorphic sections of \(q^*\mathcal V\) to
holomorphic sections on the conjugate sky. Hence it induces the
comparison of \(\mathcal T=p_*q^*\mathcal V\), and evaluation
intertwines this comparison with the supplied one on \(q^*\mathcal V\).
The descended \(H\) and \(\mathsf I\) are therefore preserved. The
affine torsor on a sky has translation line \(\mathcal O(-1,-1)\), whose
\(H^0\) and \(H^1\) vanish. Its unique section is taken to the unique
section on the conjugate sky, so the reconstructed \(\mathsf X\) is
preserved as well.

The constant partial connection along the \(q\)-fibers commutes with the
comparison, since it is pulled back from the ray bundle. In an
evaluation frame, inverse null transport reconstructs its coefficients
as the unique sections of
\(\operatorname{End}(\mathcal T)\otimes\mathcal O(1,1)\) representing
linear functions of the tangent direction. Conjugating the reconstructed
connection gives a connection with these same prescribed restrictions on
every null direction. Uniqueness of that reconstruction therefore
intertwines \(D\) with the conjugate connection. This also shows
naturality for antiholomorphic comparisons between distinct models and
for their composites.

For an involution, use these induced bundle identifications to write \[
\sigma^*\overline H=H,\qquad
\sigma^*\overline{\mathsf I}=\mathsf I,\qquad
\sigma^*\overline{\mathsf X}=\mathsf X,\qquad
\sigma^*\overline D=D.
\] Consequently \(\Theta=D\mathsf X\) is equivariant and
\(g=H(\Theta,\Theta)\) gives the asserted identity. Thus \(g\) restricts
to a real nondegenerate metric on the fixed point locus. At a fixed
mark, the prescribed Lorentz spin/tangent calibration identifies its
real tangent form with \(\eta=\operatorname{diag}(1,-1,-1,-1)\).
Nondegeneracy makes this signature constant after restriction.
Recognition identifies the complete sky with the projective null
quadric; its fixed locus is the projectivized real Lorentz null cone,
namely \(S^2\). The spin identifications, including their central signs,
are retained throughout.

Conversely, the construction in §8.4 of the local paper realizes every
marked real analytic Lorentz Einstein germ with this tangent
calibration. Its canonical whole-sky model in Proposition 4.2 of the
local paper is constructed by null geodesic flow and quadratic-affine
transport. The initialized holomorphic ODEs commute with conjugation, so
their uniqueness gives the required reality on that model. The fixed-sky
description agrees with the Lorentz spinor calculation in Proposition
I.3 of the local paper. ∎

For the Lorentz Cauchy application, fix a connected oriented compact
analytic three-manifold \(\Sigma\) without boundary, an analytic
positive-definite metric \(h\), a symmetric analytic tensor \(K\), and a
specified spatial spin structure \(P_\Sigma\). With
\(K(X,Y)=-g(\nabla_Xn,Y)\) for the future unit normal, the constraints
are
\begin{equation}\protect\phantomsection\label{eq:ordinary-analytic-constraints-v08}{
\operatorname{Scal}(h)+(\operatorname{tr}_hK)^2-|K|_h^2=-2\Lambda,
\qquad
\operatorname{div}_hK-\mathrm d\operatorname{tr}_hK=0.
}\end{equation} A Cauchy hypersurface meets every inextendible
nonconstant causal curve exactly once. The category
\(\mathsf{Ein}^{\omega}_\Lambda(\Sigma,h,K)\) consists of analytic spin
globally hyperbolic developments with the prescribed initial
identification; its morphisms preserve that identification. The spin
structure \(P_\Sigma\) is suppressed only in the category notation.
\hyperref[sec:analytic-cauchy-appendix-v08]{Appendix A} defines this
category precisely, constructs an analytic collar, and proves the
maximal analytic development statement used in §7. The local
correspondence above is applied to the analytic Einstein germs of these
developments.

The principal notation is as follows.

{\def\LTcaptype{none} % do not increment counter
\begin{longtable}[]{@{}
  >{\raggedright\arraybackslash}p{(\linewidth - 2\tabcolsep) * \real{0.5000}}
  >{\raggedright\arraybackslash}p{(\linewidth - 2\tabcolsep) * \real{0.5000}}@{}}
\toprule\noalign{}
\begin{minipage}[b]{\linewidth}\raggedright
Symbol
\end{minipage} & \begin{minipage}[b]{\linewidth}\raggedright
Meaning
\end{minipage} \\
\midrule\noalign{}
\endhead
\bottomrule\noalign{}
\endlastfoot
\(T_0\), \(\mathcal A_\Lambda\) & The fixed local flat germ and its two
original input classes \\
\(R_{a,A}\) & The nonlinear physical change of marked point and full
spin frame \\
\(B\), \(M\) & A holomorphic point manifold and its designated Lorentz
real locus \\
\(p:F\to B\), \(K\), \(E\) & The full compact sky family, its null
characteristic line and its even-contact hyperplane \\
\(\operatorname{Mon}(K)\) & Null paths modulo homotopy with endpoints
fixed \\
\(\mathcal T\), \(H\), \(\mathsf I\), \(\mathsf X\) & The rank-six
quadratic bundle, its form, Einstein scale and affine position \\
\(\mathcal T_{\mathrm{LT}}\) & The rank-four normal local-twistor bundle
with spin lift \\
\(\mathscr A\) & The calibrated affine line of positions along a ray \\
\(\mathscr E_{\mathfrak X}\), \(\mathscr E_F^{\mathrm{inc}}\),
\(\mathscr M_F\) & The contact operator algebroid, its relative ambient
germs on incidence, and their point-module family \\
\(\mathsf{SrcAtl}^{\mathrm{hol}}_\Lambda\) & Original-source atlases
with their physical comparison germs \\
\(\mathsf{MonRay}^{\mathrm{hol}}_\Lambda\),
\(\mathsf{OpMon}^{\mathrm{hol}}_\Lambda\) & Independently admitted
geometric and operator monodromy sources \\
\(M_\omega\) & The maximal analytic Cauchy subdevelopment of the smooth
maximal development \\
\end{longtable}
}

\protect\phantomsection\label{tab:local-dependencies-v095}

\textbf{Local inputs to the global results.} The table records the main
dependencies on the local paper; the cited propositions above and the
proofs below specify their hypotheses and uses. Entries for the global
theorems include the local inputs used by their intermediate lemmas.

{\def\LTcaptype{none} % do not increment counter
\begin{longtable}[]{@{}
  >{\raggedright\arraybackslash}p{(\linewidth - 4\tabcolsep) * \real{0.3333}}
  >{\raggedright\arraybackslash}p{(\linewidth - 4\tabcolsep) * \real{0.3333}}
  >{\raggedright\arraybackslash}p{(\linewidth - 4\tabcolsep) * \real{0.3333}}@{}}
\toprule\noalign{}
\begin{minipage}[b]{\linewidth}\raggedright
Global result
\end{minipage} & \begin{minipage}[b]{\linewidth}\raggedright
Results in the local correspondence
\end{minipage} & \begin{minipage}[b]{\linewidth}\raggedright
Role
\end{minipage} \\
\midrule\noalign{}
\endhead
\bottomrule\noalign{}
\endlastfoot
Theorem 3.11 & Theorem 2.4, Propositions 4.1--4.3, Theorem 8.3 &
Whole-sky recognition and the local spin-isometric Hom correspondence \\
Theorems 4.2 and 4.6 & Lemmas F.8 and H.2, Proposition H.3, Theorem H.4
& Analytic structured quantization comparisons and independent operator
recognition \\
Theorem 4.8 & Lemmas G.1--G.2, Proposition G.4, Theorem G.5, Proposition
H.3 & Relative simple modules, framed comparisons and central
specialization \\
Lemma 5.3 and Proposition 5.4 & Theorems 3.1 and 3.3; Theorem 5.5 and
Proposition 5.9; Propositions 6.1--6.2 and Theorem 6.3; Proposition 8.1
and Theorem 8.3; Theorem G.5, Proposition H.5 and Lemma K.16 & Source
dictionaries, nonlinear execution, finite-diagram convergence and
rebasing; Lemma K.16 supplies the zero-scalar amplitude model \\
Theorem 5.6 and Proposition 5.7 & Theorems 2.4, 6.3 and 8.3; Proposition
9.9 and Theorems 9.10--9.12 & Original-source construction, decoding and
the two presentation-level comparisons, using the preceding analytic
lemmas \\
Theorem 6.3 and Propositions 6.4--6.5 & Proposition I.3, Theorem J.2,
Propositions J.3--J.4 and J.7, Theorem J.6, Theorem G.5 & Spin, optical,
infinity and Lorentz comparisons, including the module square \\
Theorem 7.4 and Corollary 7.5 & The preceding local inputs, including
§8.4 and Proposition 8.4 & Lorentz admission and the real restrictions
of the source and operator comparisons \\
\end{longtable}
}

The global integration and descent arguments, including the
correspondence of germs along the entire real development, are carried
out in Sections 3--7. Section 8 proves the Nariai obstructions directly;
Appendix A constructs the analytic Cauchy category and its maximal
development, using the smooth Cauchy theorem stated there.

\protect\phantomsection\label{sec:global-null-monodromy-v08}

\section{3. Global null incidence and
monodromy}\label{global-null-incidence-and-monodromy}

We construct the geometric object on which the operator and
original-input atlases will be defined. The point manifold and its full
quadric incidence remain Hausdorff and second countable. The arrow space
of a Lie groupoid may be non-Hausdorff; quotient notation denotes that
presentation with its equivariant data. The local recognition input is
\hyperref[prop:independent-local-ray-recognition-v09]{Proposition 2.2},
in the form proved in Propositions 4.2--4.3 of the local paper.

\pandocstatementheading{\protect\phantomsection\label{prop:einstein-complexification-phase}Proposition 3.1.}

Let \((M,g)\) be a Hausdorff second-countable real analytic oriented
spin Lorentz four-manifold with \(\operatorname{Ric}(g)=\Lambda g\). It
has a Hausdorff second-countable complexification \(B\) with an
antiholomorphic involution whose fixed locus is exactly \(M\), and a
nondegenerate holomorphic Einstein metric \(g_B\) extending \(g\). The
orientation, specified spin structure and Einstein scale extend. These
data are unique as germs along all of \(M\).

The full projective null bundle
\begin{equation}\protect\phantomsection\label{eq:full-complex-null-phase-v08}{
F=\{(x,[k])\in\mathbb P(TB):g_B(k,k)=0\},
\qquad p:F\longrightarrow B,
}\end{equation} is a proper holomorphic submersion with fiber
\(Q\simeq\mathbb P^1\times\mathbb P^1\). Its nonsingular null-geodesic
line field \(K\subset TF\) projects to the tautological null direction.
Its real fixed locus \(F_{\mathbb R}\) is the real projective null phase
of \(M\).

\emph{Proof.} The complexification theorem of Whitney and Bruhat (1959,
sec. 1, Proposition 1 and the following remark) supplies a
complexification, unique as a germ, with the stated real involution.
Choose a countable representative and shrink around the entire closed
real locus. Analytic tensor coefficients extend in complexified
coordinate charts, and their transition identities extend by the
identity theorem on the maximal totally real slice. A common shrinking
gives a holomorphic metric. Its determinant is nonzero on \(M\), so it
remains nonzero on a smaller neighborhood of all of \(M\). The
holomorphic Einstein tensor vanishes there by the same identity theorem.

For the spin extension, choose local analytic lifts of the oriented
orthonormal frame transitions. Their identities extend holomorphically,
including the prescribed discrete central signs. They glue a spin double
covering of the extended frame bundle. Conjugation extends with its
Lorentz action on the two chiral factors. Local uniqueness followed by a
common shrinking proves uniqueness of the tensor, spin and reality data
as germs along \(M\).

Nondegeneracy of the metric gives the locally trivial smooth quadric
bundle in the statement; compactness of its fiber makes \(p\) proper.
The Levi-Civita connection supplies the projective geodesic spray. Its
projection is a nonzero representative of the null direction, so \(K\)
is nonsingular. Conjugation preserves it, and a fixed projective
direction has a real representative. All shrinkings concern
neighborhoods of the entire real manifold; no uniform complex radius or
geodesic completeness is required. ∎

For an object of the \hyperref[def:analytic-ghd-v2]{analytic Cauchy
category of Definition A.1}, with its specified compact Cauchy
hypersurface \(\Sigma\), let \(n\) be the future unit normal and
\(h=-g|_{T\Sigma}\). A real ray is a maximal future-oriented null
geodesic modulo orientation-preserving affine reparametrization. Write
\(\mathcal N\) for their set. The smooth contact identification with a
Cauchy cosphere bundle is classical (Chernov 2018, sec. 1); the affine
domain will also be needed.

\pandocstatementheading{\protect\phantomsection\label{prop:real-ray-cauchy-coordinates}Proposition 3.2.}

The ray space, with its quotient topology from \(F_{\mathbb R}\), is a
compact analytic contact five-manifold canonically identified with
\(S_h^*\Sigma\). Normalize each ray by
\begin{equation}\protect\phantomsection\label{eq:real-ray-affine-normalization}{
\gamma(0)\in\Sigma,
\qquad g(\dot\gamma(0),n)=1.
}\end{equation} If \(I_\gamma\) is its maximal affine interval, then
\begin{equation}\protect\phantomsection\label{eq:real-ray-affine-domain}{
D_M=\{(\gamma,s)\in\mathcal N\times\mathbb R:s\in I_\gamma\}
}\end{equation} is open and the geodesic flow is an analytic
diffeomorphism
\begin{equation}\protect\phantomsection\label{eq:real-ray-flow-identification}{
\Phi:D_M\longrightarrow F_{\mathbb R},
\qquad (\gamma,s)\longmapsto(\gamma(s),[\dot\gamma(s)]).
}\end{equation} Each sky \(Q_x=\{\gamma:x\in\gamma\}\) is an embedded
analytic Legendrian two-sphere.

\emph{Proof.} A maximal null geodesic is inextendible as a causal curve.
Indeed, near a putative causal endpoint choose a coordinate with
timelike time differential. Parametrized by this time, its directions
satisfy a regular projective geodesic equation on a compact direction
sphere. Bounded derivatives give a limiting direction, and the logarithm
of its affine speed has bounded time derivative. Its affine tangent thus
has a finite nonzero limit, extending the geodesic. The Cauchy property
now gives one intersection with \(\Sigma\).

The normalized initial tangent is \(n+v\) with \(h(v,v)=1\); associate
the covector \(h(v,\cdot)\). Analytic dependence of the geodesic
equation gives the open maximal domain \(D_M\). Unique Cauchy
intersection and geodesic uniqueness make \(\Phi\) bijective; a return
to the same phase point would give a closed causal curve. Its inverse is
analytic locally: the finite geodesic segment to \(\Sigma\) is covered
by analytic flow maps, and its hit on \(\Sigma\) is transverse because
\(g(k,n)>0\) for future null \(k\). The normalization and affine travel
time are analytic. These local inverses agree. Projection
\(D_M\to S_h^*\Sigma\) is an open submersion with interval fibers,
proving the quotient-topology assertion and compactness.

The normalized tangent defines \(\alpha=g(k,\mathrm d p\,\cdot)\) on the
real phase. For a variation \(J\) through normalized affine null
geodesics, \[
\frac{\mathrm d}{\mathrm ds}g(\dot\gamma,J)
=g(\dot\gamma,\nabla_J\dot\gamma)
=\tfrac12Jg(\dot\gamma,\dot\gamma)=0.
\] Consequently
\(\Phi^*\alpha=-\operatorname{pr}_{\mathcal N}^*\lambda\), with
\(\lambda\) the tautological contact form on \(S_h^*\Sigma\). The sign
follows from the Lorentz signature and the chosen covector. The ray
projection on \(F_{\mathbb R,x}\simeq S^2\) is injective, since a ray
cannot return to \(x\). Its differential is injective because its kernel
projects nontrivially to \(TM\), whereas \(TF_{\mathbb R,x}\) is
vertical. Compactness gives an embedding, and \(\alpha\) vanishes on
this vertical tangent space, proving the Legendrian assertion. ∎

In the specified Einstein scale, the canonical rank-six bundle is
\(\mathcal T=\mathbb R\oplus TM\oplus\mathbb R\), with
\begin{equation}\protect\phantomsection\label{eq:real-ray-canonical-tractor}{\begin{aligned}
H((\sigma,\mu,\rho),(\sigma',\mu',\rho'))
 &=\sigma\rho'+\rho\sigma'+g(\mu,\mu'), \\
D_v(\sigma,\mu,\rho)
 &=\bigl(v\sigma-g(v,\mu),\\
 &\quad \,
       \nabla_v\mu+((\Lambda/6)\sigma+\rho)v,\\
 &\quad \,
       v\rho-(\Lambda/6)g(v,\mu)\bigr), \\
\mathsf X&=(0,0,1),\qquad
\mathsf I=(1,0,-\Lambda/6).
\end{aligned}}\end{equation} These are the normal standard-tractor formulas of the
local correspondence. They also define the holomorphic data after
complexification.

\pandocstatementheading{\protect\phantomsection\label{lem:real-ray-affine-tractor}Lemma 3.3.}

The connection \(D\) preserves \(H\) and \(\mathsf I\). With
\(\Theta(v)=D_v\mathsf X=(0,v,0)\),
\begin{equation}\protect\phantomsection\label{eq:real-ray-tractor-identities}{\begin{aligned}
H(\mathsf I,\mathsf I)&=-\Lambda/3, \\
H(\mathsf I,\mathsf X)&=1, \\
H(\mathsf X,\mathsf X)&=0, \\
D_v\Theta(w)
 &=\Theta(\nabla_vw)-g(v,w)\mathsf I
                 -(\Lambda/3)g(v,w)\mathsf X.
\end{aligned}}\end{equation} Along every affine null geodesic, \(D_k^2\mathsf X=0\)
and \(D_k\mathsf X\ne0\).

\emph{Proof.} Substitution in the derivative of \(H\) cancels the scalar
and middle-component terms, proving metric compatibility. Substitution
of \(\mathsf I\) gives \(D\mathsf I=0\). The pairings follow directly.
Differentiating \((0,w,0)\) gives
\((-g(v,w),\nabla_vw,-\Lambda g(v,w)/6)\), which is the stated identity.
For an affine null tangent, both \(\nabla_kk\) and \(g(k,k)\) vanish,
while \(\Theta\) is injective. ∎

\pandocstatementheading{\protect\phantomsection\label{thm:global-real-affine-ray-realization}Theorem 3.4.}

There exist an analytic quadratic rank-six bundle \((\mathcal V,H)\) on
\(\mathcal N\), a scale section \(\mathsf I\), and an affine line
subbundle \(\mathscr A\subset\mathcal V\) satisfying \(H(Y,Y)=0\) and
\(H(\mathsf I,Y)=1\). Affine position gives an open embedding
\begin{equation}\protect\phantomsection\label{eq:global-real-affine-incidence}{
j:F_{\mathbb R}\longrightarrow\mathscr A
}\end{equation} whose image \(\mathscr D_M\) retains exactly the
intervals \(I_\gamma\). The lifted skies \(j(F_{\mathbb R,x})\) are
disjoint embedded two-spheres with normal bundle
\(T_xM\otimes\mathbb R_{S^2}\).

\emph{Proof.} Pull \(\mathcal T|_\Sigma\) back by the
initial-intersection map \(e_\Sigma:\mathcal N\to\Sigma\). Let
\(P_\gamma(s)\) be tractor transport from zero to \(s\). Finite analytic
transport equations along a compact parameter segment show that it is
analytic jointly on \(D_M\). By
\hyperref[lem:real-ray-affine-tractor]{Lemma 3.3},
\begin{equation}\protect\phantomsection\label{eq:global-real-affine-tractor-line}{
P_\gamma(s)^{-1}\mathsf X(\gamma(s))
=\mathsf X_0(\gamma)+sV_\gamma,
\qquad V_\gamma=\Theta(\dot\gamma(0))\ne0.
}\end{equation} Define
\(\mathscr A_\gamma=\mathsf X_0(\gamma)+\mathbb RV_\gamma\). This gives
an analytic trivialization \(a:\mathcal N\times\mathbb R\to\mathscr A\).
The tractor pairings imply its calibration, since \(V_\gamma\) is null
and orthogonal to \(\mathsf X_0\) and \(\mathsf I\). Set
\(j=a\Phi^{-1}\). Its image is \(a(D_M)\), so \(j\) is an open
embedding. The lifted skies are the images of the distinct fibers of
\(F_{\mathbb R}\to M\). Their normal bundles follow from the
differential of that submersion. The complete affine fiber is larger
than its realized part whenever \(I_\gamma\ne\mathbb R\). ∎

\pandocstatementheading{\protect\phantomsection\label{prop:global-real-ray-naturality}Proposition 3.5.}

A spin isometric Cauchy embedding \(f:M\to M'\) preserving orientation
and time orientation and identifying the specified initial hypersurfaces
induces a contact diffeomorphism of their real ray spaces and a
quadratic affine-bundle isomorphism preserving \(\mathsf I\). Under the
Cauchy normalization,
\begin{equation}\protect\phantomsection\label{eq:global-real-ray-naturality}{
(f_{\mathcal N}\times\operatorname{id})(D_M)\subset D_{M'},
\qquad f_{\mathscr A}\circ j=j'\circ F(f).
}\end{equation} These maps preserve lifted skies and composition, with
the specified spin lift retained.

\emph{Proof.} The Cauchy isometry identifies cosphere bundles, future
normals and normalized initial tangents. Geodesic uniqueness gives the
interval inclusion. The map
\((\sigma,\mu,\rho)\mapsto(\sigma,\mathrm df(\mu),\rho)\) intertwines
the canonical tractor connections and the initial affine generators.
Transport uniqueness gives the second identity and composition. The full
spin lift remains separate morphism data, including its central sign. ∎

For a holomorphic Einstein fourfold \(B\), retain the entire phase
\(F\), not only its real locus. The tautological even-contact hyperplane
distribution \(E\subset TF\) has characteristic line \(K\), and its Levi
form is nondegenerate on \(E/K\). This follows from the characteristic
construction on the null hypersurface in \(T^*B\).

\pandocstatementheading{\protect\phantomsection\label{def:null-monodromy-groupoid}Definition 3.6.}

The null monodromy groupoid
\(\mathcal G=\operatorname{Mon}(K)\rightrightarrows F\) has as arrows
the endpoint-preserving homotopy classes of piecewise smooth paths in a
single complex null leaf. Composition concatenates paths and inversion
reverses them. Its analytic charts are the path-continuation charts of
the foliation groupoid. Distinct path classes remain distinct even if
their transverse holonomy germs coincide.

\pandocstatementheading{\protect\phantomsection\label{prop:holomorphic-null-monodromy}Proposition 3.7.}

The groupoid \(\mathcal G\) is holomorphic of complex dimension seven,
with source-connected, source-simply-connected fibers. A countable
disjoint union \(T\) of local holomorphic transversals meeting every
leaf gives an étale groupoid \[
\mathcal G_T=T\times_{F,t}\mathcal G\times_{F,s}T\rightrightarrows T
\] with an invariant contact structure on its five-dimensional object
space. Complete transversal systems give equivalent presentations
through their common union.

\emph{Proof.} The source fiber at \(u\) is the universal cover of its
leaf. This is the foliation case of source-simply-connected integration
(Crainic and Fernandes 2003, sec. 0, pp.~577--579). The line algebroid
\(K\) has injective anchor \(K\hookrightarrow TF\). Its integration is
holomorphic by the integration theorem (Laurent-Gengoux et al. 2009,
Theorem 3.17), and has dimension \(\dim F+\operatorname{rank}K=7\).
Explicitly, a finite chain of holomorphic flow boxes along a path gives
holomorphic continuation charts, multiplication and inversion, without
requiring the arrow space to be Hausdorff.

Choose a countable cover by flow boxes and one transversal in each.
Transversality makes the source and target maps of \(\mathcal G_T\)
local biholomorphisms. The distribution induced by \(E\) on each
transversal is contact because its Levi form is the nondegenerate form
on \(E/K\); leafwise continuation preserves it. In a common union of two
systems, each inclusion is fully faithful and meets every orbit. Flow
boxes supply the local submersions required for the equivalence of
presentations. ∎

\pandocstatementheading{\protect\phantomsection\label{thm:monodromy-affine-forward-construction}Theorem 3.8.}

Every holomorphic oriented spin Einstein fourfold has the contact
monodromy presentation above, an equivariant quadratic rank-six bundle
with scale section, and an invariant calibrated affine line whose
translation line is the inverse contact line. Its proper full-quadric
incidence has an affine-position morphism
\begin{equation}\protect\phantomsection\label{eq:monodromy-affine-position-morphism}{
j:F\longrightarrow[A_F/\mathcal G]\simeq[A_T/\mathcal G_T]
}\end{equation} that is étale in local groupoid charts. The prescribed
point spin structure and local whole-sky models, with full-spin
comparisons, are retained. The construction is natural for holomorphic
spin isometries. No separated leaf quotient, trivial leafwise holonomy
or injective affine position is assumed.

\emph{Proof.} Put \(\mathcal E=p^*\mathcal T\) with its partial
connection along \(K\). Its curvature vanishes because \(K\) has complex
rank one. Parallel transport therefore depends only on the
endpoint-preserving leafwise homotopy class (Crainic and Fernandes 2003,
sec. 1.4, Proposition 1.6), and gives a representation \[
P_\gamma:\mathcal E_{s(\gamma)}\longrightarrow\mathcal E_{t(\gamma)}.
\] It is holomorphic by solving the linear transport equation in a
finite chain of holomorphic flow boxes, and preserves \(H\) and
\(\mathsf I\).

For \(u=(x,[k])\), set
\begin{equation}\protect\phantomsection\label{eq:monodromy-affine-line}{
(A_F)_u=\mathsf X_x+\mathbb C D_k\mathsf X_x.
}\end{equation} The line is independent of the representative \(k\), and
its nonzero translation vector and calibration follow from
\hyperref[lem:real-ray-affine-tractor]{Lemma 3.3}. Since
\(D_k^2\mathsf X=0\) along affine null geodesics, transport preserves
the entire affine line. Thus it is \(\mathcal G\)-invariant.

The pairing \(H(D_k\mathsf X,D_v\mathsf X)=g(k,v)\) identifies the
translation line with the inverse contact line. It is invariant under
null transport: for a commuting variation \(v\) through affine null
geodesics, its derivative is \(g(k,\nabla_vk)=\tfrac12v g(k,k)=0\). This
is the tautological contact pairing, so it supplies a nonzero duality on
each transverse chart. Restriction to \(T\) gives the asserted
equivariant bundles.

The section \(\mathsf X:F\to A_F\) defines \(j\). Its base change
against the atlas \(A_F\) is represented by
\begin{equation}\protect\phantomsection\label{eq:monodromy-affine-etale-base-change}{
\Psi:\mathcal G\longrightarrow A_F,
\qquad \Psi(\gamma)=P_\gamma\mathsf X_{p(s(\gamma))}
\in(A_F)_{t(\gamma)}.
}\end{equation} This map covers \(t\). The one-dimensional kernel of
\(\mathrm dt\) is identified by \(\mathrm ds\) with \(K\) at the source,
since the anchor is injective. Varying the source while fixing the
target differentiates \(\Psi\) to the transport of \(D_k\mathsf X\),
which is nonzero. Together with the identity on the six-dimensional base
this makes \(\mathrm d\Psi\) invertible. Hence \(\Psi\) is locally
biholomorphic and \(j\) is étale. This argument gives neither
injectivity nor the monomorphism property.

Near a point of \(B\), the original local correspondence supplies a
whole-sky short-ray model. Inclusion of its paths gives a functor from
its short-leaf monodromy groupoid to \(\mathcal G\). Its affine
comparison is the same parallel transport, so its recognition data agree
with the retained global data. Extra arrows can come from paths leaving
the chart. The point spin bundle and its comparison cocycle are retained
on \(B\). A spin isometry sends null paths and homotopies to their
images and intertwines tractor transport; therefore all these
constructions, including the full spin lift, are natural. ∎

\pandocstatementheading{\protect\phantomsection\label{prop:real-path-groupoid-v08}Proposition 3.9.}

For a holomorphic extension of an analytic globally hyperbolic
development, the designated real path groupoid is \[
\mathcal G_{\mathbb R}=\operatorname{Mon}(K_{\mathbb R})
\simeq F_{\mathbb R}\times_{\mathcal N}F_{\mathbb R}.
\] Inclusion of real paths defines a conjugation-compatible functor
\(\mathcal G_{\mathbb R}\to\mathcal G\). Its quadratic and affine
descent is exactly
\hyperref[thm:global-real-affine-ray-realization]{Theorem 3.4}.

\emph{Proof.} The real leaves are intervals by
\hyperref[prop:real-ray-cauchy-coordinates]{Proposition 3.2}, so there
is one real path class between any two points on a leaf. The real ray
submersion gives the displayed analytic groupoid. Restricting the
canonical transport to these paths is precisely the Cauchy-normalized
transport used in Theorem 3.4. ∎

The designated groupoid uses paths in the real spacetime. The full
restriction of \(\mathcal G\) to real objects can also contain nonreal
paths and additional isotropy, and need not be this real path groupoid.

We now define the independent geometric source. Let
\(\mathsf{Ein}^{\mathrm{hol}}_\Lambda\) denote holomorphic oriented spin
Einstein fourfolds with connected Hausdorff second-countable point
manifolds and holomorphic spin isometric embeddings as morphisms.

\pandocstatementheading{\protect\phantomsection\label{def:independent-monodromy-source-v092}Definition 3.10.}

An independent monodromy source consists of the following data.

\begin{enumerate}
\item
  A connected Hausdorff second-countable complex fourfold \(B\), a
  proper holomorphic submersion \(p:F\to B\) with quadric fibers \(Q\)
  and ordered rulings, a line subbundle \(K\subset TF\), and a
  hyperplane distribution \(E\subset TF\). The map \(\mathrm dp|_K\) is
  nonzero; \(E\) contains \(K\) and \(\ker\mathrm dp\), has
  characteristic line \(K\), and has nondegenerate Levi form on \(E/K\).
\item
  A holomorphic Lie groupoid \(\mathcal G\rightrightarrows F\) with
  Hausdorff connected simply connected source fibers and Lie algebroid
  identified by its anchor with \(K\). It has a rank-six representation
  \(\mathcal V_F\), an invariant nondegenerate quadratic form \(H\), an
  invariant section \(\mathsf I\) of norm \(-\Lambda/3\), and an
  invariant affine line subbundle \(A_F\). Its calibration is
  \(H(Y,Y)=0\) and \(H(\mathsf I,Y)=1\). A holomorphic position section
  \(\mathsf X:F\to A_F\) is supplied; it need not be invariant.
\item
  An atlas of independently admitted whole-sky models from
  \hyperref[prop:independent-local-ray-recognition-v09]{Proposition
  2.2}. Over \(U_i\subset B\), each model identifies its full incidence
  with \(F_i=p^{-1}U_i\), its distributions with \(K,E\), and its
  pulled-back quadratic, affine and position data with the supplied data
  on \(F_i\). Writing \(q_i:F_i\to Z_i\) for its ray projection and
  \(C_i\) for its contact distribution, these identifications give
  \(K|_{F_i}=\ker\mathrm dq_i\) and \(E|_{F_i}=(\mathrm dq_i)^{-1}C_i\).
  Its point family is complete and four-dimensional near each retained
  quadric. Actual enriched whole-sky comparison germs on \(U_i\cap U_j\)
  cover its identity and are compatible with the supplied
  identifications of incidence, distributions, quadratic-affine data and
  position. They obey identity, inverse and triple laws, including their
  actual full spin lifts. Restriction and common refinement identify
  presentations.
\item
  The representation agrees with each local short-path action.
  Precisely, the canonical integrated functor \[
  \eta_i:\operatorname{Mon}(K|_{F_i})\longrightarrow\mathcal G
  \] takes a path in \(F_i\) to an action which, under
  \(\mathcal V_F|_{F_i}\simeq q_i^*\mathcal V_i\), is the identity on
  the common fiber over its local ray. On common subdomains these
  comparisons agree under the retained whole-sky transition.
  Equivalently, the representation's infinitesimal partial connection is
  the constant relative connection of \(q_i^*\mathcal V_i\) along
  \(\ker\mathrm dq_i\).
\end{enumerate}

The functor \(\eta_i\) exists by integration of the line-algebroid
inclusion; the equivalence in clause 4 follows from source
connectedness. The clause does not identify the local groupoid with the
full restriction of \(\mathcal G\): arrows whose paths leave \(F_i\) are
retained.

A morphism consists of a point embedding \(f:B\to B'\), a holomorphic
incidence map \(\widehat f:F\to F'\) over \(f\) that is an isomorphism
on each quadric and preserves its rulings and \(K,E\), and a holomorphic
groupoid functor \(\Phi:\mathcal G\to\mathcal G'\) over \(\widehat f\).
A \(\Phi\)-equivariant quadratic isomorphism
\(b:\mathcal V_F\to\widehat f^*\mathcal V'_{F'}\) preserves
\(\mathsf I\), affine lines and position. These maps have admitted local
whole-sky representatives with specified full spin lifts, compatible
with the source and target comparison cocycles; the groupoid functor
commutes with the local integrated functors. Equality includes the full
spin lifts. Write \(\mathsf{MonRay}^{\mathrm{hol}}_\Lambda\) for this
category.

No metric is supplied. The optical and affine local recognition
conditions are tested on the given models, and source connectedness and
source simple connectedness impose global completeness of the arrows.

\pandocstatementheading{\protect\phantomsection\label{thm:independent-monodromy-equivalence}Theorem 3.11.}

The local independent correspondence induces an equivalence
\begin{equation}\protect\phantomsection\label{eq:independent-monodromy-equivalence}{
\mathsf{MonRay}^{\mathrm{hol}}_\Lambda
\simeq\mathsf{Ein}^{\mathrm{hol}}_\Lambda.
}\end{equation} It retains the chosen complex point manifold, complete
quadric family, quadratic and affine data, and specified spin structure.
It agrees locally with the original correspondence and is fully faithful
for the stated embeddings with their full spin lifts.

\emph{Proof.} Recognition reconstructs a holomorphic Einstein metric on
each \(U_i\). The admitted comparisons cover the identity on each point
overlap and respect the supplied global identifications. They therefore
identify these metrics on the fixed overlap, so the metrics glue on
\(B\); nondegeneracy and the Einstein equation are local properties. The
spin comparison cocycle glues the full spin frame bundles and
orientation.

Recognition identifies every supplied quadric with the whole projective
null quadric of the resulting metric \(g\). Thus
\begin{equation}\protect\phantomsection\label{eq:monodromy-source-phase-identification}{
\varphi:F\longrightarrow F_g,
\qquad u\longmapsto\bigl(p(u),[\mathrm dp_u(K_u)]\bigr)
}\end{equation} is a global holomorphic isomorphism, with inverse given
on each whole-sky chart. It identifies \(K,E\) with their canonical
counterparts. An injective anchor forces the algebroid bracket to be the
vector-field bracket on \(K\). The uniqueness of source-connected,
source-simply-connected integration (Crainic and Fernandes 2003, sec. 0,
pp.~578--579) therefore supplies the unique isomorphism
\begin{equation}\protect\phantomsection\label{eq:monodromy-source-arrow-identification}{
\Xi:\mathcal G\xrightarrow{\ \sim\ }\operatorname{Mon}(K_g)
}\end{equation} inducing \(\varphi\) and its algebroid map. It is
holomorphic by the flow-box construction, or by the holomorphic
integration result (Laurent-Gengoux et al. 2009, Proposition 3.20). This
identifies all arrows, including those leaving any individual point
chart.

Complete-sky evaluation gives local isomorphisms from the canonical
\(p^*\mathcal T_g\) to \(\mathcal V_F\). Their naturality for the
comparisons glues them to
\begin{equation}\protect\phantomsection\label{eq:monodromy-source-evaluation-counit}{
\beta:p^*\mathcal T_g\xrightarrow{\ \sim\ }\mathcal V_F,
}\end{equation} preserving the quadratic form, scale and position, and
taking \(\mathsf X_g+\mathbb C D_k\mathsf X_g\) to \(A_F\). Local
recognition and clause 4 intertwine the infinitesimal partial
connections. This also identifies the full groupoid actions: join a unit
to any arrow by a path in its connected source fiber. Under its target
map, both actions solve the same leafwise transport equation with the
same initial value. Uniqueness gives equality at that arrow. Hence
\(\beta\) is \(\Xi\)-equivariant on all of \(\mathcal G\), and together
with \(\varphi\) and the local full-spin comparisons identifies the
entire supplied source with the canonical one.

Conversely, \hyperref[thm:monodromy-affine-forward-construction]{Theorem
3.8} gives the canonical data for any object of
\(\mathsf{Ein}^{\mathrm{hol}}_\Lambda\). Proposition 3.7 gives the
source-fiber and anchor conditions. Original whole-sky models provide
the independent local admission, and identity metric comparisons with
the specified spin structure give their cocycles by local full
faithfulness. Evaluation identifies canonical short-path transport with
the constant relative connection, proving clause 4. Recognition recovers
the given metric.

For morphisms, a spin isometric embedding \(f\) gives its differential
map on phase and the functor \([\gamma]\mapsto[\widehat f\circ\gamma]\).
Image homotopies prove well-definedness. This functor need not be
faithful: a path can become null-homotopic in the larger point domain.
Its tractor map intertwines transport and affine position. Local full
faithfulness gives the unique admitted comparisons with the specified
spin lift, so these maps define a source morphism and respect
composition.

An admitted source morphism decodes locally to a spin isometry; its
point and spin maps glue to its supplied point embedding and full spin
lift. To prove uniqueness, fix this decoded morphism. The phase
identification forces \(\widehat f\) to be its differential on null
directions. Evaluation and local full faithfulness force its quadratic
map and comparison germs. Its differentiated groupoid functor is then
fixed. Lie II gives a unique integration because the source groupoid is
source connected and source simply connected. Thus decoding is faithful
and the preceding construction makes it full. Uniqueness of evaluation
and integration also makes both inverse comparisons natural. ∎

Changing complete transversals changes only a presentation. Shrinking
the point manifold to \(B_0\subset B\) instead requires recomputing
\(\operatorname{Mon}(K|_{p^{-1}B_0})\); it need not be the restriction
of the old arrows to those endpoints. This second operation is used for
germs along the entire real development in
\hyperref[sec:germ-correspondence-v08]{§7}.

For the obstruction argument we record when an ordinary quotient can be
used. Here a possibly non-Hausdorff complex manifold means a
second-countable space with the usual local complex manifold charts,
without the global Hausdorff axiom. A leaf relation
\(R\subset F\times F\) is \emph{embedded} when it is a complex
submanifold with the subspace topology, so its inclusion is an injective
holomorphic immersion and a homeomorphism onto its image.

\pandocstatementheading{\protect\phantomsection\label{lem:ordinary-null-quotient-v08}Lemma 3.12.}

Let \(R\subset F\times F\) be the maximal complex null-leaf relation. If
\(R\) is embedded, \(Z=F/R\) is a possibly non-Hausdorff complex contact
fivefold and \(q:F\to Z\) is an open holomorphic submersion with
connected fibers. If \(R\) is also closed, \(Z\) is Hausdorff. The
canonical quadratic bundle descends with its prescribed null transport
exactly when transport around every leafwise loop is the identity. In
that case the calibrated affine lines descend, their translation line is
the inverse contact line, and affine position \(j:F\to A\) is locally
biholomorphic. It is an open embedding exactly when it is injective on
every leaf.

\emph{Proof.} At a diagonal pair the projections of \(R\) are
submersions, since the diagonal is a section of each. Its fiber is one
leaf and its tangent relation is \(\{(v,w):w-v\in K\}\). Finite
compositions of local leafwise flows in one factor carry this conclusion
to every related pair. A small transversal \(T\) has
\(R\cap(T\times T)\) equal to its diagonal after shrinking, because
\(R\) is embedded. Thus it meets each global leaf at most once.
Saturations of open sets are open, being unions of their images under
local flows. The transversals consequently give quotient charts with
holomorphic holonomy transitions and a countable base. The distribution
\(E/K\) supplies their contact structure. If \(R\) is closed, unrelated
points have product neighborhoods disjoint from \(R\); their saturations
are disjoint, proving Hausdorffness.

Necessity of trivial loop transport follows for a pulled-back bundle
with constant relative connection. For sufficiency choose local
holomorphic sections of \(q\). Transport between any two chosen
representatives of one fiber is unique, because the fiber is connected
and loop transport is trivial. A finite chain of flow boxes along a
connecting path proves holomorphic dependence. Concatenation proves the
cocycle identities, yielding the bundle and its unique
transport-compatible comparison. The preserved form and scale descend.
Lemma 3.3 makes the affine lines invariant, and the contact pairing used
in Theorem 3.8 identifies their translation line. The derivative of
affine position is nonzero on \(K\), and it covers \(q\), so it is
locally biholomorphic in all six dimensions. Equal affine images lie
over the same ray; hence global injectivity is equivalent to leafwise
injectivity. An injective local biholomorphism is an open embedding. ∎

\protect\phantomsection\label{sec:operator-incidence-v08}

\section{4. Operator algebroids and parametrized incidence
modules}\label{operator-algebroids-and-parametrized-incidence-modules}

We now retain the operator objects of the local correspondence on the
global monodromy presentation. There are two descent problems. First,
contact transformations along arbitrary null paths act on the structured
microdifferential algebroid. Second, the calibrated simple modules near
complete compact skies are compared along point overlaps and then
descend on the parametrized incidence. Keeping these constructions
separate allows the point family to remain complete even when its
projected skies fail to be closed in a ray quotient.

\protect\phantomsection\label{sec:operator-contact-quantization-v08}

\subsection{4.1. Quantization of the contact
groupoid}\label{quantization-of-the-contact-groupoid}

Let \(\mathcal H\rightrightarrows T\) be an étale holomorphic Lie
groupoid. Its object space is a Hausdorff second-countable complex
contact fivefold, with contact distribution \(C\) and contact line
\(\mathcal L=TT/C\). Its arrow manifold may be non-Hausdorff, and every
local bisection induces a contact local isomorphism. Write
\(\mathcal H_n\) for its manifold of composable \(n\)-tuples. Sheaves
and stacks on the arrow spaces are understood on their ordinary open-set
sites; their local analytic models are defined on Hausdorff coordinate
neighborhoods.

\pandocstatementheading{\protect\phantomsection\label{def:structured-contact-groupoid-quantization}Definition 4.1.}

A structured quantization of \(\mathcal H\) is a microdifferential
algebroid \(\mathscr E_T\) with a weak \(\mathcal H\)-action. It retains
its order filtration, the full graded-symbol identification \[
\operatorname{gr}\mathscr E_T
\simeq\left(\bigoplus_{m\in\mathbb Z}\mathcal L^m\right)^+,
\] and the complex-linear formal-adjoint anti-involution, including its
square and degree-zero duality coherences. The superscript \(+\) denotes
the algebroid associated with a sheaf of algebras. Locally this is the
structured algebroid of analytic half-density microdifferential
operators on a projective cotangent chart.

The weak action consists of a structured equivalence \[
\Phi:s^{-1}\mathscr E_T\longrightarrow t^{-1}\mathscr E_T
\] on \(\mathcal H_1\), multiplication comparisons
\(\mu_{g,h}:\Phi_g\Phi_h\Longrightarrow\Phi_{gh}\) on \(\mathcal H_2\),
and unit comparisons. They satisfy the associativity and unit equations
on the groupoid nerve. Comparisons preserve all the displayed structures
and coherences. Local algebra representatives and their comparisons may
use common refinements of object and arrow charts.

The quantization inputs are Kashiwara's existence theorem (Kashiwara
1996, Theorem 2) and the structured uniqueness theorem (Polesello 2005,
Theorem 3.3). With the filtration, all graded symbols and the compatible
adjoint coherences retained, a structured equivalence is unique up to a
unique structured isomorphism. The algebra-level normalization is
\begin{equation}\protect\phantomsection\label{eq:normalized-inner-uniqueness}{
\{P:\operatorname{ord}P=0,\ \sigma_0(P)=1,\ P^*P=1\}
\xrightarrow{\ \operatorname{Ad}\ }
\operatorname{Aut}_{\mathrm{gr},*}(\mathcal E),
}\end{equation} an isomorphism of sheaves, where the automorphisms
induce the identity on all graded symbols (Polesello 2005, Lemma 1.5).
Lemmas F.8 and H.2 of the local paper supply the analytic relative
representatives and their conditioned parameter estimates. These are
local statements; refinements permit their use without requiring one
quantizer on an entire previously chosen overlap.

\protect\phantomsection\label{def:effective-structured-descent-v09}

We use effective \(2\)-descent for algebroid stacks, their equivalences
and modules on the open-set sites of these spaces. Local algebra maps
and their inner associators give the algebroid gluing data; their
compatible semilinear module maps give its modules. The bimodule
formulation of this construction applies without a paracompactness
hypothesis (Kashiwara and Schapira 2012, sec. 2.1, Proposition 2.1.13).
For structured quantization comparisons, fix the underlying contact
data; for physical source comparisons, fix the physical map and spin
lift, as in \hyperref[prop:independent-local-operator-input]{Proposition
2.3}. In each application the relevant comparison groupoid is
contractible: it is nonempty and any two comparisons admit a unique
structured isomorphism. Hence any selected overlap comparisons have
unique coherence isomorphisms on triples, and the two coherence
composites on a quadruple overlap agree. Effective \(2\)-descent then
produces the global object and comparison, uniquely up to the prescribed
structured isomorphisms. The same argument applies to equivalences of
two fixed stacks, whose functors and natural transformations satisfy
descent.

\pandocstatementheading{\protect\phantomsection\label{thm:contact-groupoid-structured-quantization}Theorem 4.2.}

Every étale holomorphic contact groupoid as above has a structured
quantization. It is unique up to a structured equivariant equivalence,
and that equivalence is unique up to a unique structured equivariant
isomorphism. The construction is natural for étale contact functors and
invariant under common complete-transversal refinement of a holomorphic
contact foliation. Neither properness of the groupoid nor a Hausdorff
coarse quotient is required.

\emph{Proof.} Choose the canonical structured algebroid on each
component of \(T\). On every Hausdorff arrow chart, the contact action
identifies the contact structures and contact lines pulled back by \(s\)
and \(t\). The two pulled-back algebroids therefore quantize the same
contact data. Local quantized contact transformations give structured
equivalences on sufficiently small arrow charts. On intersections,
structured uniqueness supplies their unique comparison. The two
comparison composites on a triple intersection agree by the same
uniqueness. The effective descent convention above consequently glues
these local equivalences to \(\Phi\) on \(\mathcal H_1\).

On \(\mathcal H_2\), the equivalences \(\Phi_g\Phi_h\) and \(\Phi_{gh}\)
have the same prescribed contact symbols and adjoint structure. Their
unique structured comparison defines \(\mu_{g,h}\) and agrees across
arrow-chart refinements. On \(\mathcal H_3\), both composites in
\begin{equation}\protect\phantomsection\label{eq:contact-groupoid-operator-associativity}{
\mu_{gh,k}\circ(\mu_{g,h}*1_{\Phi_k})
=\mu_{g,hk}\circ(1_{\Phi_g}*\mu_{h,k})
}\end{equation} are structured comparisons of the same equivalences, so
they are equal. Here \(*\) denotes horizontal composition; the formal
adjoint is a separate operation. Unit and inverse comparisons follow in
the same way. The argument uses ordinary sheaf and stack descent and
does not require a locally finite cover of a non-Hausdorff arrow space.

For a second quantization, structured uniqueness first gives an
equivalence on \(T\). Its two composites with the actions have the same
contact symbols, and hence a unique comparison on \(\mathcal H_1\).
Uniqueness on \(\mathcal H_2\) gives its equivariance, and uniqueness
again compares any two such equivariant equivalences. This proves the
assertion for weak actions without replacing them by strict actions on
chosen algebra sheaves.

An étale contact functor pulls back all the structures. Uniqueness gives
the comparison with the quantization constructed on its source, and also
its identity and composition laws. For two complete transversal systems,
take their common union as in
\hyperref[prop:holomorphic-null-monodromy]{Proposition 3.7}. Restriction
from that groupoid gives both quantizations. The complete-transversal
inclusions are fully faithful and essentially surjective; their local
arrows transport objects and morphisms between the presentations, and
the action coherence makes this transport independent of its local
choices. Descent therefore identifies the represented algebroids and
module stacks.

Analytic estimates are needed only on the chosen local representatives.
Each specified compact finite diagram admits the common conditioned
refinements of Lemmas F.8 and H.2 of the local paper. The construction
imposes no uniform positive radius or uniform symbol bound on all arrows
of a noncompact development. ∎

On local algebra representatives, the action and its associator have the
form
\begin{equation}\protect\phantomsection\label{eq:groupoid-algebroid-local-coherence}{
\begin{aligned}
\varphi_g\varphi_h
&=\operatorname{Ad}(C(g,h))\varphi_{gh},\\
C(g,h)C(gh,k)
&=\varphi_g(C(h,k))C(g,hk).
\end{aligned}
}\end{equation} The units \(C(g,h)\) have order zero, symbol one and
unitary formal adjoint. These are analytic equations on common
refinements of the nerve. They remain part of the object.

\pandocstatementheading{\protect\phantomsection\label{cor:null-monodromy-first-unit-v08}Corollary 4.3.}

The contact transversal groupoid \(\mathcal G_T\) of an independent
monodromy source has a canonical structured operator algebroid. Its
intrinsic first-unit sheaf is \(\mathcal L^{-1}\), with a strict contact
action. The source's affine line bundle on \(T\) is an equivariant
torsor for this sheaf. These constructions are natural for source
embeddings, complete-transversal changes and the functors induced by
restricting the point domain and recomputing its monodromy.

\emph{Proof.} Apply Theorem 4.2. Locally the first-unit group is
\begin{equation}\protect\phantomsection\label{eq:contact-groupoid-first-unit-sheaf}{
\mathscr U_1
=\{1+\bar A:\bar A\in F_{-1}\mathcal E/F_{-2}\mathcal E\}
\simeq\mathcal L^{-1}.
}\end{equation} Products of negative-order classes vanish in this
quotient. Normalized inner changes act trivially, since the commutator
of an order-zero operator with an order-minus-one operator has order at
most minus two. The weak action thus induces a strict action on
\(\mathscr U_1\), and its graded identification is precisely the
inverse-contact-line action. The affine line of
\hyperref[thm:monodromy-affine-forward-construction]{Theorem 3.8} has
this translation representation, so its supplied affine transport makes
it the required torsor.

A point embedding, or the inclusion of a restricted point domain,
induces the functor taking each null path to the same path in the
target. Locally its phase map identifies open phase domains. Images of
source transversals can therefore be adjoined to the target transversal
system, reducing the comparison to the étale contact functors of Theorem
4.2. Its uniqueness makes the resulting comparisons compatible with
composition. ∎

The algebroid action can identify the effects of distinct null-path
classes having the same transverse contact germ. The monodromy groupoid
and its quadratic-affine representation remain supplied data and retain
the corresponding geometric transport. Likewise the affine torsor is
additional structure, as in the independent local operator domain. Its
existence does not produce a neutralization of the algebroid. Such a
neutralization would require a global invertible object with the full
groupoid descent; a simple module supported on a Legendrian quadric is a
different kind of object (Kashiwara 1996, Theorem 2 and the following
paragraph; Polesello 2005, Corollary 3.4).

\protect\phantomsection\label{sec:operator-independent-sources-v08}

\subsection{4.2. Independent operator
sources}\label{independent-operator-sources}

Use the independent local operator domain and its inverse comparisons
from \hyperref[prop:independent-local-operator-input]{Proposition 2.3}.
In particular, the optical density root and the module half-canonical
root remain distinct. Whole-family module framing fixes a vector in the
rank-one local-system factor on the calibrated cone; it does not require
a nowhere-zero section of the negative root line on a compact quadric.
Algebra comparison units are normalized separately from these framed
semilinear module maps.

\pandocstatementheading{\protect\phantomsection\label{def:independent-operator-monodromy-source-v092}Definition 4.4.}

An independent operator monodromy source consists of the following data.

\begin{enumerate}
\tightlist
\item
  The point and phase data \(B\), \(p:F\to B\), \(K\), \(E\) and the
  source-connected, source-simply-connected groupoid
  \(\mathcal G\rightrightarrows F\) of clauses 1--2 of
  \hyperref[def:independent-monodromy-source-v092]{Definition 3.10}. Its
  source fibers are Hausdorff and its anchor identifies its Lie
  algebroid with \(K\). The quadratic-affine representation
  \((\mathcal V_F,H,\mathsf I,A_F)\) and position section \(\mathsf X\)
  are retained. No metric is supplied.
\item
  A structured analytic operator algebroid on a complete contact
  transversal presentation of \(\mathcal G\), with the first-unit
  identification of Corollary 4.3. The restriction \(A_F|_T\) is
  retained as its equivariant torsor, including its supplied affine
  translations. Changes of transversal are compared by Theorem 4.2.
\item
  A point atlas \(\{U_i\}\) of independently admitted local operator
  objects \(\mathcal O_i\). Their derived complete point families and
  full quadric incidences are identified with \(U_i\) and
  \(F_i=p^{-1}(U_i)\), preserving the ordered rulings, \(K\) and \(E\).
  Their calibrated affine and complete-sky evaluation data identify with
  \((\mathcal V_F,H,\mathsf I,A_F,\mathsf X)|_{F_i}\). The infinitesimal
  groupoid action becomes the constant relative connection of their
  short-ray projections \(q_i:F_i\to Y_i\). After shrinking the local
  point chart, these projections have the contractible short fibers of
  the local whole-sky construction.
\item
  Actual physical operator comparisons on \(U_i\cap U_j\) near each
  whole retained sky, covering its identity and compatible with the
  supplied identifications of incidence, distributions, quadratic-affine
  data and position. They retain the actual full spin lifts, root square
  identifications, torsor maps and whole-family module frames. Their
  classical incidence and full spin maps satisfy strict identity,
  inverse and triple laws. Operator representatives have normalized
  comparison isomorphisms between a composite and the representative of
  the composite, satisfying the quadruple coherence law.
\item
  Structured comparisons between each local chart algebroid and the
  pullback of the contact-groupoid algebroid along its short-ray
  presentation. They preserve full symbols and the first-unit torsor,
  and their overlap diagrams commute through normalized comparison
  isomorphisms.
\end{enumerate}

In clause 5, the short-path functor
\(\operatorname{Mon}(K|_{F_i})\to\mathcal G\) integrates the inclusion
of the line algebroid. Its short-ray quotient presents \(Y_i\). Thus the
compatibility is a comparison of these groupoid descent data; it does
not apply the physical whole-sky correspondence to an arbitrary long
null path. All operator maps and their inverse comparisons may use
common local refinements.

For a comparison between source and target operator presentations, write
\(k_a\) for its algebra maps, \(T_{ab}\) for its normalized units and
\(F_a\) for its framed semilinear module maps. The descent equations are
\begin{equation}\protect\phantomsection\label{eq:operator-source-comparison-descent}{
\begin{aligned}
k_a\varphi^s_{ab}
&=\operatorname{Ad}(T_{ab})\varphi^t_{ab}k_b,\\
T_{ab}\varphi^t_{ab}(T_{bc})C^t_{abc}
&=k_a(C^s_{abc})T_{ac},\\
F_a\tau^s_{ab}
&=T_{ab}\tau^t_{ab}F_b.
\end{aligned}
}\end{equation} Both source and target associators occur. A morphism
contains the point embedding, phase map, groupoid functor and
equivariant quadratic-affine comparison of an independent
monodromy-source morphism, together with compatible structured algebroid
comparisons and the above physical comparisons on whole-sky charts. Its
full spin lift is retained. Two representatives of a morphism are
identified only by normalized comparison isomorphisms over the same
underlying maps and spin label. Composition uses the retained
coherences. Write \(\mathsf{OpMon}^{\mathrm{hol}}_\Lambda\) for the
resulting category of physical Hom classes.

These admission conditions concern operator symbols, calibrated torsors,
intrinsic optical recognition, complete point families and specified
groupoid actions. They can be checked on the supplied source. Atlases
related by restriction and common refinement represent the same source
presentation, with all structures and comparison maps restricted
together. A globally supported module on a ray quotient is not an
admission condition.

\pandocstatementheading{\protect\phantomsection\label{lem:operator-atlas-quadratic-completion}Lemma 4.5.}

The quadratic-affine representation and position section in Definition
4.4 can equivalently be recovered from the local operator charts, their
symbolic comparison cocycle and the supplied source-complete groupoid.
The recovery uses whole-sky evaluation and inverse null transport,
without supplying a global Einstein metric.

\emph{Proof.} Apply the full-symbol and calibrated-affine functor of
\hyperref[prop:independent-local-operator-input]{Proposition 2.3}. The
independently admitted local ray presentation and complete-sky
evaluation give a rank-six bundle \(\mathcal T_i\), connection \(D_i\),
form \(H_i\) and sections \(\mathsf I_i,\mathsf X_i\) on \(U_i\). Their
connection is the one reconstructed by inverse null transport in
\hyperref[prop:independent-local-ray-recognition-v09]{Proposition 2.2}.

Actual symbolic comparisons cover the identity of each point overlap,
respect the supplied incidence identifications and preserve all these
objects. Normalized operator comparison isomorphisms leave the symbols
and first-unit torsor unchanged. Consequently the weak operator cocycle
induces a strict cocycle on the evaluation bundles, their connections
and the point spin bundles. These data glue on \(B\).

Put \(\mathcal V_F=p^*\mathcal T\) and restrict its pulled-back
connection to \(K\). This partial connection is flat because \(K\) has
complex rank one. On every short-ray chart, complete-sky evaluation
identifies it with the constant relative connection on the local ray
bundle. The affine lines and position sections agree on point overlaps
and therefore give \(A_F\) and \(\mathsf X\) globally. Transport
preserves \(H\), \(\mathsf I\) and these affine lines.

The anchor and source-completeness hypotheses identify the supplied
groupoid with \(\operatorname{Mon}(K)\) by
\hyperref[thm:independent-monodromy-equivalence]{Theorem 3.11}. Flat
partial transport hence integrates to a holomorphic representation of
that groupoid. Any representation with this infinitesimal action agrees
with it on each source fiber by initialized transport uniqueness. The
same argument proves uniqueness of the completion and naturality for
physical embeddings, spin comparisons and restrictions. ∎

\pandocstatementheading{\protect\phantomsection\label{thm:operator-monodromy-equivalence}Theorem 4.6.}

Full symbols and calibrated affine data give an equivalence
\begin{equation}\protect\phantomsection\label{eq:operator-monodromy-equivalence}{
\mathsf{OpMon}^{\mathrm{hol}}_\Lambda
\simeq
\mathsf{MonRay}^{\mathrm{hol}}_\Lambda.
}\end{equation} It preserves the point manifold, complete quadric
incidence, affine position, selected scale and spin structure. On each
whole-sky chart it is the local operator correspondence. It is fully
faithful for the embedding morphisms of
\hyperref[def:independent-monodromy-source-v092]{Definition 3.10}. Both
inverse comparisons retain the operator associators and framed
point-module comparisons.

\emph{Proof.} Taking full symbols and the calibrated torsor of each
\(\mathcal O_i\) gives an independently admitted local ray model. Clause
3 identifies its whole-sky incidence with the supplied phase data, and
clause 4 gives a symbolic and full-spin comparison cocycle over the
identity on point overlaps, compatible with those global
identifications. The infinitesimal-action condition is precisely the
short-path condition of the independent monodromy-source definition. The
global phase, groupoid and quadratic-affine data are already retained,
so the resulting object is an admitted ray source.

For morphisms, full symbols and first-unit torsor maps are the admitted
physical ray comparisons. Normalized comparison isomorphisms preserve
their symbolic data and spin label. The operation therefore descends to
physical Hom classes. The two-associator equations above show that it
preserves composition, giving a functor \(\mathcal F\).

Conversely, start with an independent ray source. Theorem 4.2 quantizes
its contact transversal groupoid. Retain the supplied affine
representation under the first-unit identification of Corollary 4.3. On
each whole-sky chart, apply the inverse local operator correspondence to
its independently admitted affine model. This gives an operator object
with its compact module family, root and frame data.

The short-ray presentation maps to the full monodromy presentation by
including paths. Its local chart algebroid and the pullback of the
groupoid algebroid quantize the same contact data. Structured uniqueness
compares them on refined contact charts, with unique normalized
comparisons on intersections. Their first-unit torsor is the same
supplied affine torsor, giving clause 5 of Definition 4.4.

Lift each physical ray comparison over the identity on a point overlap
by the local operator theorem, with its specified spin lift and the same
supplied global identifications. A composite of lifts and the lift of
the composite cover the same physical map and spin label. The local
structured comparison groupoid over those fixed data is contractible, so
there is a unique normalized comparison between them. On a quadruple
overlap the two comparison composites agree by uniqueness. The framed
module comparisons satisfy the same coherence through their semilinear
descent equations and whole-family framing. Uniqueness also makes these
comparisons compatible with the contact-groupoid algebroid. This
constructs the inverse functor \(\mathcal Q\), up to its specified
normalized comparisons.

For a noncompact point manifold, use relatively compact point charts
contained in larger analytic parents. Properness of \(p\) retains every
full compact sky. Each finite comparison diagram can then be placed in
finitely many buffered parents after refinement, where the conditioned
local estimates apply. This gives analytic representatives and their
inverses for every required local diagram. The unique comparisons glue
them over the entire cover and the local charts of the possibly
non-Hausdorff arrow space. No estimate uniform over all global charts is
needed.

The local inverse comparisons identify \(\mathcal F\mathcal Q\) with the
supplied ray source, preserving its global phase, groupoid and
transport. They also preserve its spin cocycle, so they glue globally.

For \(\mathcal Q\mathcal F\), apply the local counit to each supplied
operator chart on a common refinement of the supplied and reconstructed
presentations. Its algebra and module components satisfy the displayed
comparison equations with both associators. On a point overlap, the two
counit composites cover the same physical comparison and spin label;
unique normalized comparison identifies them. The same uniqueness gives
coherence on triples. Theorem 4.2 compares the global algebroids and
makes that comparison compatible with the chart counits. Thus the counit
is a comparison of the full operator presentations, rather than only of
their decoded geometric objects.

To prove fullness, fix a ray-source embedding. Its local whole-sky
comparisons lift by the local operator correspondence, while its
contact-groupoid map has the comparison constructed in Corollary 4.3.
Structured uniqueness identifies the two constructions on short-ray
charts, and the comparison equations glue the lifts on point overlaps.
This gives an operator-source morphism. The induced functor between
monodromy groupoids need not be faithful; its supplied phase and contact
action are sufficient.

For faithfulness, two lifts of the same ray-source morphism have the
same physical local maps and spin lifts. The unique normalized operator
and framed module comparisons identify them on every chart. These
comparisons agree on overlaps by uniqueness, and the global algebroid
comparisons agree as well. Hence the lifts have the same physical Hom
class. Distinct spin labels are never identified. The unit, counit and
composition diagrams commute through these same comparisons and
therefore as equalities of physical Hom classes. ∎

\protect\phantomsection\label{sec:operator-parametrized-incidence-v08}

\subsection{4.3. The module on parametrized
incidence}\label{the-module-on-parametrized-incidence}

Let \(\mathfrak X=[T/\mathcal G_T]\) denote the contact quotient
presentation, let \(\mathscr E_{\mathfrak X}\) be its structured
algebroid, and let \(q:F\to\mathfrak X\) be the canonical phase
morphism. The proper point map \(p:F\to B\) remains a map of Hausdorff
manifolds. Here \(B\) is the physical point parameter, denoted by \(S\)
in the relative-module construction in Appendix G of the local paper; it
is not that appendix's additional external parameter base.

The operators below retain ambient contact germs along incidence,
including operators transverse to a point sky. They are not intrinsic
differential operators on \(F\).

\pandocstatementheading{\protect\phantomsection\label{prop:relative-operator-incidence-germs}Proposition 4.7.}

There is an analytic relative structured algebroid
\(\mathscr E_{\mathfrak X\times B/B}\). Its local symbols are jointly
holomorphic in contact coordinates and the passive parameter \(b\in B\),
with the local factorial bounds of the analytic operator domain, and its
operators differentiate only in the contact model. The inverse-image
algebroid
\begin{equation}\protect\phantomsection\label{eq:relative-operator-incidence-algebroid}{
\mathscr E_F^{\mathrm{inc}}
=(q,p)^{-1}\mathscr E_{\mathfrak X\times B/B}
}\end{equation} is well defined on \(F\) and contains
\(p^{-1}\mathcal O_B\) centrally. It is independent, up to structured
comparison, of transversals and local representatives, and natural for
physical embeddings and recomputed point restrictions.

\emph{Proof.} On a projective cotangent chart, take analytic relative
half-density operators with passive parameter in a point chart. Their
composition, filtration and adjoint are the corresponding relative
operator formulas; the adjoint fixes the passive parameter. Jointly
holomorphic symbols and the factorial estimates on smaller product
parents give the analytic relative algebra constructed in §§F.1, G.1 and
H.1 of the local paper.

Every contact-groupoid comparison and associator extends with the point
parameter passive. The exact operator equations remain valid
coefficientwise. Changes of point coordinates act by substitution. These
data give the relative algebroid and its weak action on
\(\mathcal G_T\times B\). Theorem 4.2 and the relative comparisons of
\hyperref[prop:independent-local-operator-input]{Proposition 2.3}
establish independence of the presentations.

For an explicit construction on \(F\), choose phase neighborhoods
\(V_i\) with holomorphic short-path choices into contact transversals.
They give maps \(\chi_i:V_i\to T\). On an overlap, two choices are
related by the groupoid arrow obtained by composing one chosen path with
the inverse of the other. These arrows satisfy the groupoid cocycle.
Pull back the relative operator germs along \((\chi_i,p)\) and glue
their algebroids using the weak operator comparisons and their
associators. The groupoid-nerve coherence is exactly the required
descent coherence. A change of short-path choices gives the
corresponding comparison.

This construction takes inverse images of ambient ring germs. It takes
neither a quotient by functions vanishing on incidence nor a tensor
product with \(\mathcal O_F\) over ambient functions. Thus transverse
operator germs remain available. The central passive functions give
\(p^{-1}\mathcal O_B\). Physical embeddings and recomputed restrictions
carry the local paths to their images; the comparisons of Corollary 4.3
and the same descent equations give naturality. ∎

\pandocstatementheading{\protect\phantomsection\label{thm:global-parametrized-point-module-family}Theorem 4.8.}

The calibrated whole-sky module families of an independent operator
monodromy source descend effectively to a locally finitely generated
algebroid module
\begin{equation}\protect\phantomsection\label{eq:global-parametrized-point-module-family}{
\mathscr M_F\in\operatorname{Mod}(\mathscr E_F^{\mathrm{inc}}).
}\end{equation} Its support on \(F\) is all of \(F\). Its local
realizations retain the root \(R^2=K_{F/B}\), optical calibration,
whole-family frames and operator associators. It is unique up to the
unique compatible calibrated comparison and natural for physical source
embeddings and recomputed point restrictions. Its Lorentz comparison is
the weak comparison established in \hyperref[sec:spin-reality-v08]{§6}.

For \(x\in B\), let \(i_x:F_x\hookrightarrow F\) and let
\(\mathfrak m_x\subset\mathcal O_{B,x}\) be the maximal ideal. The point
fiber is the central specialization
\begin{equation}\protect\phantomsection\label{eq:incidence-point-module-specialization}{
\begin{aligned}
\mathscr E_x^{\mathrm{inc}}
&=i_x^{-1}\mathscr E_F^{\mathrm{inc}}
  /\mathfrak m_x i_x^{-1}\mathscr E_F^{\mathrm{inc}},\\
\mathscr M_x
&=i_x^{-1}\mathscr M_F
  /\mathfrak m_x i_x^{-1}\mathscr M_F.
\end{aligned}
}\end{equation} The compact quadric \(F_x\) with these germs is locally
the calibrated simple sky object of the original correspondence at
\(x\).

\emph{Proof.} On a whole-sky point chart \(U_i\), the relative
construction of Proposition H.3 of the local paper, recalled in
\hyperref[prop:independent-local-operator-input]{Proposition 2.3}, gives
a module family \(\mathcal N_i\) on \(Y_i\times U_i/U_i\). Its support
is the incidence embedded by \[
\jmath_i=(q_i,p):F_i\longrightarrow Y_i\times U_i.
\] The differential is injective because \(\ker\mathrm dq_i=K\) and
\(\mathrm dp\) is nonzero on \(K\). The map is injective since each sky
is embedded and its point parameter is retained. It is proper as a map
into the retained product neighborhood: the inverse image of a compact
set is closed in the compact inverse image under \(p\) of that set's
point projection. Hence \(\jmath_i\) is a proper holomorphic embedding.
This argument concerns the local whole-sky representative, independently
of the global ray quotient.

Clause 5 of Definition 4.4 identifies the inverse-image relative chart
algebroid with \(\mathscr E_F^{\mathrm{inc}}|_{F_i}\). Thus
\(\jmath_i^{-1}\mathcal N_i\) is a module over that restriction. Its
finite local presentations remain finite presentations after inverse
image, which is exact on sheaves and can be checked stalkwise. Each
stalk is nonzero by the locally standard simple incidence model.

On a point overlap, the physical operator comparison intertwines the two
relative families. After transport to the fixed incidence algebroid, its
algebra and module components satisfy precisely the comparison equations
(\ref{eq:operator-source-comparison-descent}), including the supplied
and reconstructed associators. On triple overlaps these give the module
descent equation through the specified algebroid associator; the unique
normalized algebra and framed module comparisons give quadruple
coherence.

More explicitly, refine by algebra and module charts. The local module
sheaves have semilinear transition maps \(\tau_{ab}\) satisfying
\(\tau_{ab}\tau_{bc}=C_{abc}\tau_{ac}\). This is the descent datum
defining a module for the algebroid represented by \(\varphi_{ab}\) and
\(C_{abc}\). Sheaf descent glues the local objects with these equations.
Finite generation and nonzero stalks are preserved by the transition
equivalences, giving the asserted support. Root square identifications,
any retained parameter-base factors and their frames descend through
their comparison maps. The framing remains on the calibrated cone, as
specified in
\hyperref[prop:independent-local-operator-input]{Proposition 2.3}.

Any two choices have unique local calibrated comparisons. Their
restrictions agree on overlaps by uniqueness and hence glue. The same
argument proves naturality for embeddings and recomputed restrictions.
The square equations in
\hyperref[prop:operator-monodromy-lorentz-reality]{Proposition 6.5}
descend through these same comparisons, giving the asserted Lorentz
compatibility.

It remains to identify the fiber. Inverse image to \(F_x\) retains
passive parameter germs, so the quotient by \(\mathfrak m_x\) is
necessary. In a product chart centered at \(x\), evaluation of relative
operators has kernel generated by the point coordinates. Indeed, for
each holomorphic coefficient symbol, \[
a(y,b)-a(y,x)
=\sum_\ell(b^\ell-x^\ell)
\int_0^1\partial_{b^\ell}a(y,x+t(b-x))\,\mathrm dt.
\] On a smaller convex parameter parent, one Cauchy gap bounds the
integral coefficients with the same factorial symbol order. Applying the
formula coefficientwise gives division inside the analytic relative
algebra, and therefore identifies its central specialization with the
ordinary contact operator germ. The calculation commutes with the
algebra comparisons, so it also applies to the algebroid.

The relative flattenings and finite defining-ideal presentations in
Lemmas G.1--G.2 and Proposition H.3 of the local paper specialize to the
fixed-parameter simple-module presentation. Hence ordinary central
specialization gives precisely the local sky module. This uses the
admitted locally standard presentation, not a flatness claim about
arbitrary relative modules. Roots, frames and comparison equations
specialize simultaneously, completing the identification of the compact
point object. ∎

\pandocstatementheading{\protect\phantomsection\label{cor:operator-incidence-enrichment-v08}Corollary 4.9.}

Theorem 4.6 retains the global incidence module \(\mathscr M_F\) and all
specialized compact point objects
\((F_x,\mathscr E_x^{\mathrm{inc}},\mathscr M_x)\) naturally.
Equivalently, the local point-module component of Definition 4.4 can be
presented by this incidence module together with its locally standard
realizations, root identifications and frame comparisons. This change of
presentation preserves physical Hom classes.

\emph{Proof.} Theorem 4.8 constructs the incidence module from the
admitted local families and their comparisons. Conversely, the specified
locally standard realizations recover those local families and their
comparisons. On a closed embedded local incidence, a supported sheaf is
determined as a germ by its inverse image: direct image by that
embedding recovers its local supported sheaf. This stalkwise
identification respects the ambient operator action. The specified local
realizations retain the required finite operator presentations, so the
statement does not enlarge the admissible module class.

The two constructions are inverse through the unique calibrated
comparisons of Theorem 4.8 and commute with physical maps. Their
compatibility with the functors and both counits of Theorem 4.6 proves
the assertion, including its naturality under recomputed point
restriction. ∎

The module constructed here lives on the Hausdorff parametrized
incidence, with its ambient relative operator germs. Its support there
is all of \(F\), and its compact point labels are the fibers of \(p\).
The construction supplies no descent of that module along every null ray
and no globally supported coherent module on a projected sky in the
coarse quotient. This distinction will be essential in the Nariai
example.

\protect\phantomsection\label{sec:original-source-atlases-v08}

\section{5. Atlases of the original source
classes}\label{atlases-of-the-original-source-classes}

We now globalize the original input itself. The chart data are the two
cohomology classes of
\hyperref[prop:original-two-class-local-input]{Proposition 2.1}, and the
transition data are their nonlinear physical comparisons. The definition
records the topology and comparison law; the analytic and descent
arguments then construct their global geometric and operator
realizations. We use the descent-data and common-refinement conventions
of the Stacks Project (The Stacks project authors 2026, Tag 02ZC).

Use the original source domain and analytic constructor recalled in
\hyperref[sec:local-input-v08]{§2}. An original marked input is a pair
\begin{equation}\protect\phantomsection\label{eq:atlas-original-source-pairs}{
\mathbf a_i=(a_i^+,a_i^-)
\in\mathscr H_2(T_{0,i})\oplus\mathscr H_{-6}(T_{0,i})
}\end{equation} at the common Einstein coefficient \(\Lambda\), with its
full spin mark. Each \(T_{0,i}\) is a marked copy of the germ of
\(\operatorname{Tot}(\mathcal O(1)^{\oplus2})\) along the entire zero
line. The fixed residue and middle-Laurent projection selects its
canonical representatives. Raw representatives differing by a Čech
coboundary represent the same input; an arbitrary choice of a coboundary
primitive is not an additional physical arrow.

The convergent original constructor supplies a complete local point
domain \(P_i\) and its full compact sky incidence. A \emph{retained
point domain} means a representative inside the constructor's analytic
parents, with the whole compact sky over every retained point. Further
restriction and recentering are allowed. The six normalization
requirements belong to that prescribed construction; they are not six
additional freely chosen source classes.

On a common point domain, a physical source comparison is specified
locally after recentering at a point of that domain. It is the original
short instruction \((a,A)\) and its source comparison, with the full
spin element \(A\) retained. Its point map uses the covariant convention
in Theorem 8.3 of the local paper: the instruction
\(\mathbf a\to R_{a,A}\mathbf a\) is represented by \(\Xi_{a,A}^{-1}\).
Comparisons between other retained representatives are those of Isthmus
(2026a, Proposition 8.1). Their composition uses the source-dependent
physical law, not addition of displacement vectors or multiplication of
tensor labels alone.

The action on the two source classes is computed by the exact residue
dictionaries, the original Cartan flow on whole records, and the inverse
dictionaries. It need not be the pullback of both classes by a map
between the old flat seeds. In particular, the global input below is
descent data for this nonlinear physical comparison law.

\pandocstatementheading{\protect\phantomsection\label{def:independent-original-source-atlas}Definition 5.1.}

An original-source atlas consists of the following data.

\begin{enumerate}
\item
  A connected Hausdorff second-countable complex four-manifold \(B\).
\item
  A covering \(\{U_i\}\) of \(B\), original marked inputs
  \(\mathbf a_i\) as above, retained derived point domains \(P_i\), and
  biholomorphic identifications \(e_i:P_i\to U_i\).
\item
  On every point overlap, physical source comparison germs \(c_{ji}\)
  from the retained \(i\)-presentation to the retained
  \(j\)-presentation, covering \(e_j^{-1}e_i\). They are germs near
  every whole compact sky in the overlap, with holomorphic
  representatives after the conditioned local restrictions of
  \hyperref[prop:source-atlas-local-analytic-execution]{Proposition
  5.4}. The central points of the old charts need not belong to the
  overlap: recenter both charts near a common point before specifying a
  short comparison.
\item
  These comparisons satisfy identity, inverse and triple laws as
  original physical arrows, including their full spin lifts:
  \begin{equation}\protect\phantomsection\label{eq:original-source-atlas-cocycle}{
  c_{ii}=1,\qquad
  c_{ij}=c_{ji}^{-1},\qquad
  c_{kj}\circ c_{ji}=c_{ki}.
  }\end{equation}
\end{enumerate}

Equality here is equality of the original physical class and spin label
on the common germ domain. The contact-affine representatives obey the
corresponding strict comparison laws in Proposition 8.1 of the local
paper. Operator lifts, when constructed, retain their normalized inner
comparison isomorphisms.

Restriction of charts, recentering and common refinement give equivalent
presentations through specified physical comparison maps. Every
refinement still covers the same retained point manifold \(B\). Its
comparison maps retain the original classes and their full spin lifts.
These maps are changes of presentation; distinct physical morphisms are
not identified merely because their underlying point maps agree. In
particular, a refinement does not add arbitrary raw Čech zero-cochains
as new arrows.

Restricting the retained point manifold itself to an open subset
\(U\subset B\) is a separate operation producing a source object over
\(U\). It is not a change of presentation of the original object over
\(B\).

A morphism consists of a holomorphic point embedding \(f:B\to B'\) and
physical source comparison germs on common source and target point
charts covering \(f\). They have full spin lifts and commute with the
source and target overlap comparisons. Morphisms are compared on a
common refinement using the same specified source and target refinement
identifications, and are identified only when their germs and spin lifts
agree there. Independent changes of spin identification cannot turn the
negative central spin morphism into the identity. Composition is the
original physical composition on representatives, followed by common
refinement. Denote the resulting category by
\(\mathsf{SrcAtl}^{\mathrm{hol}}_\Lambda\).

Admission in this definition is tested on the two original analytic
cohomology classes, their record-flow action and physical composition
law, and the identifications of their derived complete point domains.
These data are defined before the geometric decoder; the global metric
is constructed in
\hyperref[thm:original-source-atlas-operator-equivalence]{Theorem 5.6}.

Equivalently, the point topology may be specified by the generated
domains \(P_i\) and their physical transition germs: require the usual
full atlas identity, inverse and cocycle laws, including the spin
comparisons, and require the quotient of the disjoint union by these
identifications to be connected, Hausdorff and second countable, with
each \(P_i\) mapping by an open embedding. The quotient then has a
unique complex structure for which these maps are holomorphic charts,
since the supplied transition maps are biholomorphic and agree as germs
on overlaps. This recovers the manifold \(B\) and the identifications in
the definition. The topology and the full spin transition law remain
source-atlas data.

\pandocstatementheading{\protect\phantomsection\label{rem:global-source-continuation-v08}Remark 5.2.}

The point topology and continuation data in this definition cannot be
recovered from one original germ pair. For example, the flat
developments \(\mathbb R\times(\mathbb R^3/L\mathbb Z^3)\) with metric
\(\mathrm dt^2-\sum_j(\mathrm dx^j)^2\) have zero scalar and zero
curvature records, hence the same original pair \((0,0)\) at a chosen
full spin mark. Their Cauchy volumes are \(L^3\), so different \(L\)
give distinct marked global data. For fixed \(L\), a proper time collar
and the full flat development have the same local pair but different
realized point domains. An atlas retains precisely the additional
identifications and domains involved in these examples.

\pandocstatementheading{\protect\phantomsection\label{lem:compatible-source-amplitude-diagrams}Lemma 5.3.}

Let \(u=(U^+,U^-;\Lambda)\) be an original analytic input with its
marked full spin frame. Fix a compact amplitude disk
\(|\varepsilon|\le T_b\), where \(T_b>1\) lies strictly inside a
retained larger disk, and a finite local star of point and full-spin
instructions, together with its specified inverses and composites. After
choosing sufficiently small passive displacement neighborhoods and
smaller point parents, these instructions have compatible bounded
holomorphic representatives over the entire compact amplitude disk.
Every vertex is obtained by rebasing the same central physical family
\begin{equation}\protect\phantomsection\label{eq:source-compatible-central-amplitude}{
u_\varepsilon
=(\varepsilon U^+,\varepsilon U^-;\varepsilon\Lambda).
}\end{equation} For passive instructions \((a_i,A_i)\), the vertex input
is
\begin{equation}\protect\phantomsection\label{eq:source-compatible-vertex-input}{
v_i(\varepsilon)=R_{a_i,A_i}u_\varepsilon
=\bigl(V_i^+(\varepsilon),V_i^-(\varepsilon);
       \varepsilon\Lambda\bigr).
}\end{equation} Its prescribed source constructor uses a separate
internal amplitude \(s\) on
\begin{equation}\protect\phantomsection\label{eq:source-independent-internal-amplitude}{
\bigl(sV_i^+(\varepsilon),sV_i^-(\varepsilon);
      s\varepsilon\Lambda\bigr),
}\end{equation} and is evaluated at \(s=1\) before the physical vertex
comparisons are made. The resulting analytic contact, affine, operator
and point-module comparisons retain their full spin labels, both
calibrated roots, associators and whole-family frames. The assertion
includes \(\Lambda=0\).

\emph{Proof.} The local analytic inputs are the results of Theorem 3.3,
Propositions 6.1--6.2, Theorem 6.3, Theorem 8.3 and Theorem G.5 of the
local paper. The compatible-amplitude construction is also displayed for
zero scalar in Lemma K.16 of the local paper. We give the finite-star
argument with the Einstein coefficient retained.

Propositions 6.1--6.2 of the local paper construct the central family
\(g_\varepsilon\) and its marked normal coordinates, initialized spin
transports, prescribed ray charts and augmented stored tuples on common
smaller parents over the compact amplitude disk. The metric and inverse,
the finitely many required derivatives, and the relevant inverse maps
have common bounds there.

Let \(\nu_i(\varepsilon)\) be the initialized full-spin transport to
\(x_i(\varepsilon)=\operatorname{Exp}^{g_\varepsilon}_o(a_i)\). The
exact record action makes \(v_i(\varepsilon)\) the whole curvature
record of the same \(g_\varepsilon\) at that point in the frame
\(\nu_i(\varepsilon)A_i\). Theorem 3.3 of the local paper gives a common
factorial bound on a smaller record radius for sufficiently small
displacements. The fixed finite spin parameters and their inverses are
bounded after restriction to compact parameter subcores. The two residue
encoders therefore give bounded holomorphic representatives of these
vertex inputs on common smaller annular neighborhoods of the whole
reference line.

The initialized metric comparisons are
\begin{equation}\protect\phantomsection\label{eq:source-compatible-initialized-comparisons}{
\Xi_i^\varepsilon:g[v_i(\varepsilon)]\longrightarrow g_\varepsilon,
\qquad
y\longmapsto
\operatorname{Exp}^{g_\varepsilon}_{x_i(\varepsilon)}
                    \bigl(\nu_i(\varepsilon)A_i y\bigr).
}\end{equation} At zero displacement their point differentials are the
invertible vector representations of the retained spin maps. Finitely
many fixed spin frames can be used as reference values; their central
signs remain distinct spin labels. The analytic inverse theorem,
compactness of the amplitude disk and the finite number of specified
maps give smaller common domains for \(\Xi_i^\varepsilon\), their
inverses and the requested composites. The passive neighborhoods can be
chosen so that their images have positive margins inside the assigned
parents. The required determinants and scalar coefficients used as units
then have positive modulus minima. This condition concerns the inverted
coefficients and determinants, not every individual matrix entry.

On these domains put
\begin{equation}\protect\phantomsection\label{eq:source-compatible-star-comparisons}{
f_{ij}^\varepsilon
=(\Xi_i^\varepsilon)^{-1}\Xi_j^\varepsilon.
}\end{equation} Their inverse and triple laws are identities of maps,
with the same initialized spin comparisons. Expressing them as original
source instructions gives the composition law in Isthmus (2026a, Theorem
8.3), including its radial-versus-broken-path spin factor. This is a
coordinate comparison law; it does not declare spin parallel transport
path independent.

The rebased inputs are bounded holomorphic families in \(\varepsilon\)
and the passive parameters. Apply the prescribed source construction to
each such family using the independent variable \(s\) in the statement.
On a compact internal disk of radius greater than one, Propositions
6.1--6.2 and Theorem 6.3 of the local paper supply its analytic
representative and inverse components. At \(s=1\) it is the source
associated with \(v_i(\varepsilon)\). The physical comparison
\(f_{ij}^\varepsilon\) is then exactly a comparison to which Theorem G.5
of the local paper applies. Its proof already uses the nonlinear outer
family and evaluates each vertex's independent internal constructor at
one. It gives the analytic structured algebra maps, normalized
comparison units and framed module maps, with both associators retained.
Normalized-inner uniqueness and framed-module uniqueness give all the
stipulated finite comparison equations. The Einstein coefficient is
\(\varepsilon\Lambda\) throughout the outer family and
\(s\varepsilon\Lambda\) in the internal constructor. The cited formulas
have no division by this scalar.

Finally, the full sky remains present under these restrictions. The
central sky over the compact retained parameter set is compact. A finite
cover of it by buffered direction charts and short hit-map domains gives
a common smaller point neighborhood on which these charts still cover
the full proper quadric family. All the finite comparison data, root
maps and module charts are restricted together. Thus the finite minima
above are positive without intersecting an uncontrolled infinite
collection of neighborhoods. ∎

The two amplitudes have different roles. In general,
\begin{equation}\protect\phantomsection\label{eq:source-amplitude-rebase-noncommutation}{
R_{a,A}(\varepsilon U^+,\varepsilon U^-;\varepsilon\Lambda)
\ne
\bigl(\varepsilon(R_{a,A}u)^+,
      \varepsilon(R_{a,A}u)^-;\varepsilon\Lambda\bigr).
}\end{equation} Consequently the lemma supplies physical comparisons
after each internal constructor is evaluated at one. It does not compare
separately linearly scaled endpoint inputs at every internal amplitude.

The central physical family in this lemma belongs to one local
comparison star. Different stars need not share an outer amplitude
family. In particular, the lemma does not deform a whole global source
atlas, or an arbitrary global loop of point charts, to one flat source.

\pandocstatementheading{\protect\phantomsection\label{prop:source-atlas-local-analytic-execution}Proposition 5.4.}

Consider a collection of original local source presentations with
physical comparison germs, on a Hausdorff second-countable point
manifold. Assume the supplied holomorphic families and specified finite
local comparison diagrams have the conditioned parent, inverse and unit
margins of Definition 2.1 of the local paper. On every retained compact
point and passive-parameter subsystem, the collection may be refined to
finitely many representatives on which the original prescribed nonlinear
construction and its physical comparisons are analytic, including their
whole compact skies. Their Einstein coefficient and six scale
normalizations are preserved through the calibrated comparisons.

The construction may therefore be used on a locally finite source atlas
over a noncompact point manifold. Its estimates are local on compact
subsystems and finite comparison diagrams. It requires no common
positive radius or common norm bound over the whole point manifold.

\emph{Proof.} The two source representatives determine their complete
factorially bounded records by the exact residue dictionaries.
Conversely, the same dictionaries recover unique original cohomology
classes from those records, with representative and coboundary
comparisons on smaller buffered domains (Isthmus 2026a, Theorem 3.1). To
rebase at a point of an overlap, join the old mark to that point by a
piecewise analytic path inside the existing constructor's analytic
representative. Use the metric already derived by that local
constructor. A compact path is covered by finitely many relatively
compact point and spin charts inside larger analytic parents. The chiral
curvature, transported in parallel spin frames along sufficiently short
complex null geodesics, is holomorphic in the geodesic parameter on a
common positive disk for normalized spinor directions. Its \(r\)th
derivatives give diagonal evaluations of the free records, so Cauchy
estimates give \(M r!R^{-r}\). Polarization and the fixed spinor norms
cost only an exponential factor in \(r\), absorbed by one radius
decrease. Thus the complete records at the intermediate marks have
refreshed factorial bounds, uniformly on each retained compact
passive-parameter subsystem.

Subdivide the path so that successive marks lie in common short normal
parents. Each displacement and specified spin-frame change is then an
admitted \(R_{a_j,A_j}\) as in Theorems 3.3 and 8.3 of the local paper.
Their finitely many local source constructions and comparisons have
common smaller parents by the finite-word assertion. The exact
dictionaries return original source classes at every step, and local
inverse uniqueness identifies the resulting recentered germ with the
given analytic representative. This gives rebasing by a finite chain of
short instructions. The full spin labels are retained even when their
action on tensor records has a kernel.

For the coherent point-overlap diagrams used in atlas descent, work near
a common overlap point and refine to a local star with one reference
mark. Express its finitely many retained marks by admissible short
instructions, with its comparisons induced by the same initialized
identifications.
\hyperref[lem:compatible-source-amplitude-diagrams]{Lemma 5.3} gives
coherent physical families and analytic representatives on smaller
parents. It is not used to interpolate arbitrary independently specified
automorphism or loop relations to a flat family. Other admitted finite
diagrams of completed analytic objects use the conditioned
finite-diagram assertion of Theorem 9.10 of the local paper directly.
When the starting data already include a bounded coherent passive
family, use that family's comparison identities and apply the analytic
construction locally with those parameters retained. In each case it is
only necessary to reserve the buffers for the finitely many specified
maps, inverse maps and composite images.

The nonlinear execution in each point chart is the one prescribed in
§§5--6 of the local paper. Its correction at a weight \(m\ge2\) is
\begin{equation}\protect\phantomsection\label{eq:atlas-original-nonlinear-correction}{
\Delta c_{m+1}
=-L_m a_m+\iota_{m+1}Z_{m-2}.
}\end{equation} Here \(\Pi_{m+1}\) is the free-record map from the
closed contact correction classes at that weight to
\(W_{m-2}^+\oplus W_{m-2}^-\), \(\iota_{m+1}\) is its fixed right
inverse, and \(Z_{m-2}\) is the prescribed minus current free-record
component. The map \(o_m\) takes a closed contact correction class to
its optical obstruction class, and \(L_m=(o_m|_{\ker\Pi_{m+1}})^{-1}\)
on the compatible optical defects. Thus the first term removes the
defect without changing the free record, and the second supplies its
prescribed record correction, with \(\Pi_{m+1}\iota_{m+1}=1\).

The actual optical defect \(a_m\) and current record are read from that
chart's freshly recomputed incidence. The remaining flat scalar freedom
is \(a+b\cdot x+\gamma\eta(x,x)\). In the lifted lower Einstein scale
the already selected section is represented by \(1\); \(a,b,\gamma\)
describe the highest-stage correction (Isthmus 2026a, Theorem 5.5). The
value and four first derivatives fix its first five coefficients, and
the numerical coefficient fixes
\begin{equation}\protect\phantomsection\label{eq:atlas-sixth-scale-normalization}{
\gamma=(\lambda_{\mathrm{target}}-\lambda_*)/12.
}\end{equation} In this formula \(\lambda_{\mathrm{target}}\) and
\(\lambda_*\) are the prescribed and currently computed numerical-scalar
coefficients in the highest stage being corrected; the original Einstein
coefficient remains \(\Lambda\). These are the original six
normalizations, including their transverse scalar condition, rather than
only central sky data. The deterministic stored normalizer uses the
selected geometry and section, initialized normal and spin development,
its prescribed contact Moser maps, the same graph evaluations and fresh
optical recomputation.

Theorem 6.3 of the local paper identifies every finite prescribed stored
tuple with the corresponding internal-amplitude Taylor polynomial of one
analytic tuple. For each of finitely many selected vertices, write this
tuple as \(Y_i(s)\) and its degree-\(N\) polynomial as \(Y_{i,N}(s)\).
After the finite restrictions, choose \(T_*>1\) smaller than the
internal Cauchy radii and a common finite bound \(M\) for these finitely
many classical augmented tuples. Their analytic sums obey
\begin{equation}\protect\phantomsection\label{eq:finite-source-atlas-classical-tail}{
\max_i\|Y_i(1)-Y_{i,N}(1)\|
\le
\frac{M T_*^{-N-1}}{1-T_*^{-1}}.
}\end{equation} All required inverse components are included in these
tuples. The analytic operator and supported-module realizations use the
relative constructions of Appendices F--H of the local paper. In those
constructions the finite active operator presentations are compared
through normalized inner equivalences; the displayed classical tuple
estimate does not assert a literal equality of arbitrary raw quantum
representatives.

Physical comparisons preserve the same selected section and its
calibration by Proposition 8.1 and Theorem 9.12 of the local paper. Its
value and gradient are transported and recomputed in the chosen
auxiliary density frame. In particular, after an auxiliary scale change
they transform with the Cauchy-state matrix; their old coordinate
entries are not held fixed. The Einstein coefficient is preserved. Hence
the six normalizations agree geometrically under the comparison, even
when the stored coordinate tuples differ because their presentation
choices differ. The correction law and its obstruction are the same
local source law, including the preservation correction identified in
Proposition H.5 of the local paper.

A compact point subsystem lies in finitely many smaller chart cores with
larger parents. At each point of a retained overlap, the finite diagram
under consideration has representatives on a smaller neighborhood by the
preceding argument. Compactness yields a finite refinement covering the
required overlap subcores. Properness of the sky family makes its
inverse image over each compact core compact, so finitely many direction
charts and common finite buffer minima retain every whole sky. Repeating
this on a locally finite refinement gives source presentations on the
entire point manifold. Local identities are tested after common finite
refinement and agree as germs; no global minimum of all radii is taken.

Every constructed map is jointly holomorphic in the retained passive
parameters. Restriction to an actual finite Taylor neighborhood of a
coherent passive family preserves its comparison equations. This proves
the stated local analytic control and the compatibility used for
source-atlas descent. ∎

The proposition preserves the two original source operations and their
nonlinear correction as local executions joined by physical comparisons.
It does not identify the Hamiltonian representatives in distinct
reference charts as one global Hamiltonian. Nor does it give an exact
atlas cocycle for the polynomials obtained by truncating each vertex's
independent internal amplitude separately. Exact descent concerns the
analytic evaluations at internal amplitude one. Finite-jet descent
concerns a coherent passive or outer physical family, such as the common
family of \hyperref[lem:compatible-source-amplitude-diagrams]{Lemma
5.3}, whose comparison identities already hold before the jet
restriction.

\pandocstatementheading{\protect\phantomsection\label{lem:source-atlas-phase-completion}Lemma 5.5.}

An original-source atlas determines a proper holomorphic quadric bundle
\(p:F\to B\), its ordered rulings, a characteristic line \(K\), a
hyperplane distribution \(E\), and the point spin structure. These are
the point and phase data of
\hyperref[def:independent-monodromy-source-v092]{Definition 3.10}. The
source-complete null-path groupoid is canonically
\begin{equation}\protect\phantomsection\label{eq:source-atlas-monodromy-completion}{
\mathcal G=\operatorname{Mon}(K)\rightrightarrows F.
}\end{equation} Complete-sky evaluation and inverse null transport
determine its quadratic-affine representation and position. These
constructions use the source comparisons, are natural for source-atlas
morphisms, and commute with point restriction when the path groupoid is
recomputed.

\emph{Proof.} Apply the original analytic source constructor on each
retained chart. Its output has the full proper quadric incidence
\(p_i:F_i\to P_i\), line \(K_i\) and hyperplane \(E_i\). The comparisons
in Propositions 8.1 and 9.9 and Theorem 9.12 of the local paper induce
holomorphic incidence identifications over the specified point maps,
preserving the ordered rulings and these distributions. The source
cocycle (\ref{eq:original-source-atlas-cocycle}) makes their classical
incidence and full spin cocycles strict. Thus they glue the \(F_i\) over
\(B\) to a holomorphic quadric bundle, and glue \(K_i,E_i\) to \(K,E\).
Each local bundle is a bundle of the entire compact quadric, so \(p\) is
proper. The total space is Hausdorff and second countable by its local
bundle charts over \(B\).

The conditions \(\mathrm dp|_K\ne0\), \(K+\ker\mathrm dp\subset E\),
characteristic line \(K\) and nondegenerate transverse Levi form are
local conditions and hold in every constructed chart. The full spin
cocycle also glues the local spin bundles and their frame maps. No spin
sign is inferred from its action on a projective fiber.

Apply \hyperref[prop:holomorphic-null-monodromy]{Proposition 3.7} to the
resulting holomorphic characteristic line. Its monodromy integration has
the required source-fiber and anchor properties. It retains paths
through arbitrarily many point charts; the construction has not
identified those paths with a single short-ray overlap.

On each whole-sky chart, evaluation and inverse null transport
reconstruct \((\mathcal T_i,D_i,H_i,\mathsf I_i,\mathsf X_i)\). The
source comparisons preserve these objects. Their strict classical
cocycle glues them on \(B\). The proof of
\hyperref[lem:operator-atlas-quadratic-completion]{Lemma 4.5} now
applies: pull \(\mathcal T\) to \(F\), restrict its connection to \(K\),
and integrate this flat partial connection on \(\operatorname{Mon}(K)\).
It preserves the quadratic form, scale and affine lines, with position
retained as an incidence section. This determines the entire
representation from the source atlas.

A source-atlas morphism preserves the local incidence, evaluation and
spin maps. They glue to its phase and quadratic-affine maps, and mapping
paths by its phase map gives the groupoid functor. The local transport
equations prove equivariance and composition. For \(U\subset B\), the
restricted atlas glues the restricted phase \(F|_U\) and line
\(K|_{F|_U}\). Its completion is consequently
\(\operatorname{Mon}(K|_{F|_U})\), with the integrated inclusion into
the old groupoid. This is precisely the recomputed restriction of
\hyperref[prop:recomputed-monodromy-restriction]{Proposition 7.1}. ∎

\pandocstatementheading{\protect\phantomsection\label{thm:original-source-atlas-operator-equivalence}Theorem 5.6.}

Use the original local source results (Isthmus 2026a, Theorems 2.4, 6.3,
8.3 and 9.10--9.12) with their stated analytic-domain and physical-arrow
conventions, together with the contact-groupoid quantization proved in
\hyperref[thm:contact-groupoid-structured-quantization]{Theorem 4.2}.
The original source constructor and decoder extend to an equivalence
\begin{equation}\protect\phantomsection\label{eq:original-source-atlas-operator-equivalence}{
\mathsf{SrcAtl}^{\mathrm{hol}}_\Lambda
\simeq
\mathsf{OpMon}^{\mathrm{hol}}_\Lambda.
}\end{equation} It retains the chosen complex point manifold, full
compact point family, spin structure, selected Einstein scale and the
completed null-path representation. On each chart it is the original
two-input source correspondence. Both inverse comparisons preserve the
source presentations through their initialized physical comparisons, and
preserve the operator associators and framed local module comparisons.
The equivalence is fully faithful for the stated embedding morphisms.

\emph{Proof.} Start with a source atlas.
\hyperref[lem:source-atlas-phase-completion]{Lemma 5.5} constructs its
phase, groupoid and quadratic-affine data. On the \(i\)-chart apply the
source operator constructor \(\mathcal S_{\mathrm{op}}\) of Theorem 9.10
of the local paper to \(\mathbf a_i\). Its full symbols and first-unit
torsor are the original constructed contact-affine tuple, and its
locally standard compact modules are the original complete sky family.
Thus it supplies an admitted local operator chart on the retained point
domain.

The physical source comparisons \(c_{ji}\) cover \(e_j^{-1}e_i\) and
hence, under the retained identifications \(e_i\), the identity of
\(U_i\cap U_j\). Their operator and framed module lifts preserve the
completed incidence and quadratic-affine identifications by Proposition
9.9 and Theorem 9.10 of the local paper. On a triple overlap, the lift
of \(c_{kj}c_{ji}\) and the chosen lift of \(c_{ki}\) cover the same
physical map and spin label. The unique normalized comparison
isomorphism identifies these two operator representatives. Its module
component obeys the semilinear comparison equations
(\ref{eq:operator-source-comparison-descent}), with both associators. On
a quadruple overlap, the two composites of these comparison isomorphisms
have the same source, target and normalization, so they agree. This is
the weak operator cocycle required by
\hyperref[def:independent-operator-monodromy-source-v092]{Definition
4.4}.

Choose complete contact transversals of the global characteristic
foliation.
\hyperref[thm:contact-groupoid-structured-quantization]{Theorem 4.2}
supplies their structured groupoid algebroid. On each short-ray chart,
its pullback and \(\mathcal S_{\mathrm{op}}(\mathbf a_i)\) are
structured quantizations of the same contact symbols. Their structured
uniqueness supplies the comparison in clause 5 of
\hyperref[def:independent-operator-monodromy-source-v092]{Definition
4.4}, and the source affine torsor is retained under their first-unit
identification. Unique normalized comparisons make these identifications
compatible on intersections and with the physical operator lifts. All
clauses of the independent operator source are therefore satisfied. A
source-atlas morphism yields the corresponding phase and groupoid maps
by \hyperref[lem:source-atlas-phase-completion]{Lemma 5.5}, and its
local physical operator lifts glue by these same equations. This
constructs the functor \(\mathcal Q\) on physical Hom classes.

Conversely, let \(\mathcal O\) be an independent operator monodromy
source over \(B\). Near every point choose one of its admitted whole-sky
operator models, recentered at that point when necessary. Apply the
independent decoder \(\mathcal D_{\mathrm{op}}\) from Theorem 9.10 of
the local paper. It gives an original pair \(\mathbf a_i\) at the same
Einstein coefficient and with the same mark. The counit
\begin{equation}\protect\phantomsection\label{eq:source-atlas-local-operator-counit}{
\epsilon_i:\mathcal S_{\mathrm{op}}\mathcal D_{\mathrm{op}}(\mathcal O_i)
\longrightarrow\mathcal O_i
}\end{equation} identifies complete point families, so it gives an
identification \(e_i\) of a retained derived source point domain with a
neighborhood \(U_i\subset B\).

These neighborhoods can be chosen to cover all of \(B\). If a local
counit is initially represented only near its central sky, recenter the
supplied operator chart at each other point and use its local theorem
there. Then choose a point-cover refinement by such neighborhoods. This
refines the atlas and does not shrink the retained global point
manifold. Relatively compact subcharts inside larger parents supply the
conditioned local representatives when needed.

On a common point neighborhood, compose the supplied operator comparison
with the counits and their inverses. Since the supplied comparison
covers the identity of \(U_i\cap U_j\), this composite covers
\(e_j^{-1}e_i\) and respects the supplied global identifications. Local
full faithfulness in Isthmus (2026a, Theorem 9.10) gives a unique
original physical source comparison, with its specified full spin lift,
between the two decoded pairs after the necessary recentering. These
local comparisons agree as germs where two choices overlap. Indeed,
their operator images are the same physical Hom class, and the local
source-to-operator map is faithful on that class with its full spin
label.

The supplied operator cocycle commutes through its normalized comparison
isomorphisms. In physical Hom classes it is an equality. Faithfulness of
the original source correspondence therefore gives the strict identity,
inverse and triple laws of the source atlas in
(\ref{eq:original-source-atlas-cocycle}). The source comparisons are
holomorphic on their common parents by the analytic family assertion of
Theorem 9.10 of the local paper. Their contact-affine representatives
are those of Isthmus (2026a, Proposition 8.1), hence obey the retained
comparison equations. We obtain an independently admitted
original-source atlas without using a metric-image criterion. The same
local decoding applied to operator source morphisms defines a functor
\(\mathcal D\).

We verify the two inverse comparisons at the stated presentation level.
For a source atlas, the local equality
\begin{equation}\protect\phantomsection\label{eq:source-atlas-original-class-unit}{
\mathcal D_{\mathrm{op}}\mathcal S_{\mathrm{op}}(\mathbf a_i)
=\mathbf a_i
}\end{equation} is the equality of the two original input classes,
Einstein coefficient and mark in Theorem 9.11 of the local paper. The
fixed canonical projection gives the same canonical representatives. If
the comparison uses recentered charts or different retained parents,
Isthmus (2026a, Proposition 8.1) gives the initialized contact-affine
comparison on a common refinement. The full-spin physical-arrow
bijection identifies the decoded overlap instruction with the original
\(c_{ji}\). Consequently these local comparisons glue to the
source-atlas unit. This recovers the original source data and
comparisons.

For the operator composite, use the counits \(\epsilon_i\) from the
counit formula (\ref{eq:source-atlas-local-operator-counit}). On a point
overlap their two composites have the same physical comparison and full
spin lift. The unique normalized comparison identifies them with their
module components and both operator associators. On triples those
identifications are coherent for the same reason. Thus the local counits
define a comparison of the complete operator-chart descent data.

The phase identification supplied by these counits glues globally. The
reconstructed groupoid is the monodromy integration of that phase line,
while the supplied groupoid has the same anchor and source-completeness
conditions. The integration uniqueness of
\hyperref[thm:independent-monodromy-equivalence]{Theorem 3.11}
identifies the two groupoids. The whole-sky evaluation comparisons
identify their quadratic-affine bundles and partial connections;
initialized transport then identifies their full actions. Finally,
structured uniqueness in
\hyperref[sec:operator-contact-quantization-v08]{§4.1} identifies the
global contact-groupoid algebroids and makes that identification
compatible with the local counits. This is a counit in
\(\mathsf{OpMon}^{\mathrm{hol}}_\Lambda\), retaining its module and
associator data.

For fullness, take an operator source embedding between two objects
constructed from source atlases. On every common source and target point
chart the original local correspondence decodes its physical Hom class
to a source comparison with the given point map and spin lift. On
intersections these local decoded comparisons agree by local
faithfulness. They therefore define a source-atlas morphism. Its image
has exactly the original local operator Hom classes. The global phase
map is forced by the full incidence, the groupoid functor by its
differentiated map and source simple connectedness, and the
quadratic-affine map by evaluation and transport. The global algebroid
comparison is fixed up to its unique structured isomorphism. Hence this
image is the given operator source morphism.

For faithfulness, two source-atlas morphisms with the same operator
image have equal local physical arrows, including their full spin
labels, by the physical Hom bijection of Theorem 9.10 of the local
paper. They agree after common refinement and hence are the same
source-atlas morphism. Identity and composition are preserved by their
original physical laws. The local unit and counit naturality diagrams
commute by Theorem 9.10 of the local paper; their normalized operator
comparisons glue by uniqueness. This proves the equivalence and its
stated presentation-level naturality. ∎

\protect\phantomsection\label{def:conditioned-holomorphic-source-families}

A \emph{coherent holomorphic family of global sources} is specified over
a finite-dimensional complex manifold \(P\) by a holomorphic submersion
\(\pi:B_P\to P\) with four-dimensional point fibers and local
presentations of the relevant source type. Each point fiber carries an
admitted object of that source category. In an original-source
presentation, the chart inputs \(\mathbf a_i\), derived point-domain
identifications \(e_i\) and physical overlap comparisons \(c_{ji}\)
depend holomorphically on \(P\) in the conditioned sense of §2, and
their identity, inverse and cocycle equations hold before any parameter
truncation. Operator presentations retain holomorphic associators,
normalized comparison maps, roots and framed module maps satisfying
their stated equations. The geometric presentations retain the
corresponding holomorphic incidence and transport data. On each retained
compact subsystem and specified finite comparison diagram, the common
parents, inverse bounds and unit margins are those of Proposition 5.4.
The null monodromy groupoid is formed fiberwise; the parameter
assertions concern these local presentations and finite comparison
diagrams.

At \(p_0\in P\), write \(P_{p_0,N}\) for the finite Taylor neighborhood
with local algebra \(\mathcal O_{P,p_0}/\mathfrak m_{p_0}^{N+1}\).
Preservation of finite parameter jets means that the completed local
constructor and its comparisons commute with restriction to
\(P_{p_0,N}\) on every retained finite diagram. Thus for the
source-to-operator constructor \(\mathcal Q\),
\begin{equation}\protect\phantomsection\label{eq:source-family-finite-jet-base-change}{
\left.\mathcal Q(\mathfrak A)\right|_{P_{p_0,N}}
\simeq
\mathcal Q\!\left(\left.\mathfrak A\right|_{P_{p_0,N}}\right),
}\end{equation} where the right side denotes the Taylor restriction of
the prescribed analytic local construction and its descent equations.

\pandocstatementheading{\protect\phantomsection\label{prop:source-atlas-complete-diagram}Proposition 5.7.}

The equivalence of
\hyperref[thm:original-source-atlas-operator-equivalence]{Theorem 5.6}
retains the full local source-operation diagram in Theorem 9.12 of the
local paper on every atlas chart. In particular, it uses the same two
original classes, corrected nonlinear source constructor, six
normalization requirements and selected analytic representatives. These
local constructions and their physical comparisons determine one global
operator-incidence source. Its rank-four spin, optical and infinity data
of \hyperref[sec:spin-reality-v08]{§6}, and its calibrated incidence
module of \hyperref[sec:operator-parametrized-incidence-v08]{§4.3}, are
natural images of the same source atlas.

For a supplied coherent holomorphic family of source atlases on
conditioned common domains, these comparisons are holomorphic and
preserve its finite parameter jets on every retained finite diagram. The
assertion concerns the completed analytic local constructors and
coherent passive families. It does not identify independently truncated
internal construction-amplitude stages on different point charts.

\emph{Proof.} On each chart, use the source constructor of
\hyperref[thm:original-source-atlas-operator-equivalence]{Theorem 5.6}.
The comparison of Theorem 9.12 of the local paper identifies its
operations, correction cochains, fixed inverses, normalizations and
analytic representatives with those of the independent operator chart.
On an overlap, the comparison is the physical source comparison required
in \hyperref[def:independent-original-source-atlas]{Definition 5.1}.
Theorem 9.12 of the local paper identifies its contact, affine, optical,
operator and framed module components. The comparison and associator
equations used in
\hyperref[thm:original-source-atlas-operator-equivalence]{Theorem 5.6}
therefore glue these completed local outputs to one global source.

The selected scale and roots are transported by the physical
comparisons. Recomputing the six normalization components in a different
auxiliary chart represents this same selected calibration; it is not an
equality of six chart-independent raw numbers. The analytic
normalization and local finite-diagram comparison of
\hyperref[prop:source-atlas-local-analytic-execution]{Proposition 5.4}
give their common representatives.

The normal spin diagram is natural for these comparisons by
\hyperref[thm:global-projective-diagram-naturality]{Theorem 6.3}. The
same is true of the global incidence module by
\hyperref[thm:global-parametrized-point-module-family]{Theorem 4.8} and
\hyperref[cor:operator-incidence-enrichment-v08]{Corollary 4.9}. Hence
both are carried through the source-atlas unit and operator counit, with
their spin and root data.

For an admitted coherent parameter family, fix a supplied finite
conditioned comparison diagram.
\hyperref[prop:source-atlas-local-analytic-execution]{Proposition 5.4}
provides its local analytic representatives. A cocycle equation is
checked near a common overlap point, using the finite local star there
and the compatible construction of
\hyperref[lem:compatible-source-amplitude-diagrams]{Lemma 5.3}. All
composition and inverse equations are equalities of these holomorphic
maps or of their specified normalized comparison maps. Passing to finite
Taylor neighborhoods preserves these equations. A compact set can be
covered by finitely many such local parents, retaining their separate
bounds. Gluing uses these local representatives and their germs, with
the finite-star and local-bound conventions of
\hyperref[lem:compatible-source-amplitude-diagrams]{Lemma 5.3} and
\hyperref[prop:source-atlas-local-analytic-execution]{Proposition 5.4}.

Physical comparisons use each original constructor at its completed
analytic value, or an already supplied coherent passive family. The
distinction from independent internal-amplitude scaling is expressed by
(\ref{eq:source-amplitude-rebase-noncommutation}) in the discussion
following \hyperref[lem:compatible-source-amplitude-diagrams]{Lemma
5.3}. ∎

\protect\phantomsection\label{sec:spin-reality-v08}

\section{6. The projective spin diagram and Lorentz
reality}\label{the-projective-spin-diagram-and-lorentz-reality}

Let an independent monodromy source be given, and let \((B,g)\) be its
recognized holomorphic spin Einstein manifold. Write \(p:F\to B\) for
its full null phase, \(K\) for the characteristic line and
\(\mathcal G=\operatorname{Mon}(K)\) for its recognized groupoid. The
complete-sky evaluation data are
\((\mathcal T,D,H,\mathsf I,\mathsf X)\), and
\(\Theta(v)=D_v\mathsf X\). The Einstein coefficient is fixed by
\(\operatorname{Ric}(g)=\Lambda g\).

The local exterior-square, flag, optical and Lorentz comparisons are
those of Proposition I.3, Theorem J.2 and Propositions J.4 and J.7 of
the local paper, assembled in Theorem 9.12 of that paper. Their tensor
formulas are equivariant for the specified full spin maps. We use them
to construct the bundles on the point manifold and prove their
continuation along all null-path classes. The point spin structure is
retained before any quotient by rays.

\pandocstatementheading{\protect\phantomsection\label{prop:global-local-twistor-monodromy}Proposition 6.1.}

The point spin structure determines a rank-four holomorphic bundle
\(\mathcal T_{\mathrm{LT}}\), its volume-preserving normal spin
connection and a parallel identification
\begin{equation}\protect\phantomsection\label{eq:global-spin-exterior-square}{
\mathcal F_{\mathrm{spin}}:\mathcal T
\xrightarrow{\ \sim\ }\Lambda^2\mathcal T_{\mathrm{LT}}^*.
}\end{equation} The pullback \(\mathbb L=p^*\mathcal T_{\mathrm{LT}}\)
is a holomorphic representation of \(\mathcal G\) preserving its volume
and a flag \(\ell\subset h\subset\mathbb L\) of ranks one and three. For
the contact-line representation \(\mathcal L_F=TF/E\),
\begin{equation}\protect\phantomsection\label{eq:global-spin-flag-contact-line}{
\ell\otimes\operatorname{Ann}(h)\simeq\mathcal L_F^{-1}.
}\end{equation} On each point quadric \(Q_x\), \(\mathbb L\) is constant
with fiber \((\mathcal T_{\mathrm{LT}})_x\), while
\(\ell|_{Q_x}=\mathcal O(0,-1)\) and
\(\operatorname{Ann}(h)|_{Q_x}=\mathcal O(-1,0)\) in the specified order
of rulings.

The alternating form
\(I_{\mathrm{cov}}=\mathcal F_{\mathrm{spin}}(\mathsf I)\) and its
volume dual \(I_{\mathrm{contra}}\) are invariant. With \(c=\Lambda/6\),
their normalization is
\begin{equation}\protect\phantomsection\label{eq:global-infinity-pair-normalization}{
I_{\alpha\beta}I^{\beta\gamma}
=c\,\delta_\alpha{}^\gamma,
\qquad
2\operatorname{Pf}(I_{\mathrm{cov}})=-\Lambda/3.
}\end{equation} Both forms have rank four for \(\Lambda\ne0\) and rank
two for \(\Lambda=0\).

\emph{Proof.} In the selected Einstein scale, the spin bundle supplies
the two determinant-one spinor bundles and their rank-four sum. The
local normal spin connection is the Levi-Civita spin connection plus the
algebraic soldering term with coefficient \(c\). The exterior-square
comparison of Proposition I.3 of the local paper is spin-equivariant and
identifies its induced connection with \(D\). These formulas therefore
glue through the spin cocycle. Their trace is zero and their transition
determinants are one, giving the parallel volume and the displayed
global identification.

Pull the connection back to \(F\) and restrict it to \(K\). Its
curvature on this complex line distribution vanishes. Its transport
along a leafwise path consequently depends only on its endpoint-fixed
homotopy class. Denote it by \(\Pi_\gamma:\mathbb L_u\to\mathbb L_v\)
for \(\gamma:u\to v\). A representative path is covered by finitely many
holomorphic flow boxes. In each box, transport solves a holomorphic
linear matrix equation; finite composition gives holomorphic dependence
in a path-continuation chart. Homotopy invariance and ODE uniqueness
identify these expressions on chart intersections and give the identity,
inverse and composition laws. This constructs a representation on the
full, possibly non-Hausdorff, arrow manifold.

The parallel exterior-square comparison gives
\begin{equation}\protect\phantomsection\label{eq:global-spin-tractor-transport}{
\mathcal F_{\mathrm{spin},v}(P_\gamma Y)
=(\Pi_\gamma^{-1})^*
\mathcal F_{\mathrm{spin},u}(Y),
}\end{equation} where \(P_\gamma\) is the rank-six transport of the
geometric source. Transport around a nontrivial loop is retained in this
equation.

For a null tangent \(k^{AA'}=\kappa^A\pi^{A'}\), the local flag is
\(\ell=\mathbb C(0,\pi)\) and \(h=\ker(\kappa_A,0)\) in the spin
conventions of the local comparison. Rescaling either factor leaves
these subspaces unchanged. Along an affine null geodesic choose parallel
spin factors. The skew-spinor identities make both \((0,\pi)\) and
\((\kappa_A,0)\) parallel for the normal spin connection and its dual.
Thus every \(\Pi_\gamma\) preserves the flag. Their tensor product is
the tautological null-covector line, identified with
\(\mathcal L_F^{-1}\) by the characteristic pairing
\(g(k,\mathrm dp\,\cdot)\). Parallelism preserves this identification.
Restriction to \(Q_x\) gives the two stated tautological ruling lines.

Finally \(\mathsf I\) is parallel. In a compatible spin splitting, the
fixed volume and the selected section give \[
J_2=\begin{pmatrix}0&1\\-1&0\end{pmatrix},
\qquad
I_{\mathrm{cov}}=\operatorname{diag}(-cJ_2,J_2),
\qquad
I_{\mathrm{contra}}=\operatorname{diag}(J_2,-cJ_2).
\] The local exterior-square normalization is
\(2\operatorname{Pf}\mathcal F_{\mathrm{spin}}(Y)=H(Y,Y)\). The
displayed matrices therefore give both stated identities and ranks. They
glue by the same volume-preserving spin comparisons. This construction
keeps the selected Einstein section and the normalization
\(H(\mathsf I,\mathsf X)=1\); replacing its infinity form by an
arbitrary projective multiple would discard that scale information. No
inverse infinity form is used, so \(\Lambda=0\) is included. ∎

\pandocstatementheading{\protect\phantomsection\label{prop:global-optical-pencil-incidence}Proposition 6.2.}

There is a holomorphic totally null rank-two subrepresentation
\begin{equation}\protect\phantomsection\label{eq:global-optical-plane-representation}{
\mathbb V_u
=\operatorname{span}\{\mathsf X_x,\Theta_x(k)\}
\subset(p^*\mathcal T)_u,
\qquad u=(x,[k]).
}\end{equation} Its dual is the monodromy continuation of the intrinsic
optical solution bundle, with calibrated Wronskian
\(\Lambda^2\mathbb V^*\simeq\mathcal L_F\). The affine line is exactly
\begin{equation}\protect\phantomsection\label{eq:global-optical-affine-slice}{
(A_F)_u=\{v\in\mathbb V_u:H(\mathsf I,v)=1\}.
}\end{equation} Taking the kernel of the exterior-square form gives an
equivariant isomorphism from \(\mathbb P(\mathbb V_u)\) to the pencil of
two-planes \(L\) with \(\ell_u\subset L\subset h_u\). The position
selects
\begin{equation}\protect\phantomsection\label{eq:global-optical-point-plane}{
L_u^{\mathrm{point}}
=\ker\mathcal F_{\mathrm{spin}}(\mathsf X_x).
}\end{equation} This selected incidence section records the point's
position in the invariant pencil and need not itself be
groupoid-invariant.

\emph{Proof.} Rescaling \(k\) leaves the plane unchanged. The identities
\(H(\mathsf I,\mathsf X)=1\), \(H(\mathsf I,\Theta(k))=0\) and
\(\Theta(k)\ne0\) give independence of its generators. Their three
mutual \(H\)-pairings vanish. Along an affinely parametrized null
geodesic the canonical Einstein-scale connection satisfies \[
D_k\mathsf X=\Theta(k),
\qquad
D_k\Theta(k)=0.
\] Thus \(\mathbb V\) is parallel along \(K\), and its affine slice is
\(\mathsf X+\mathbb C\Theta(k)\).

On a simply connected short leaf segment, a parallel covector
\(\lambda\in\mathbb V^*\) determines \(y(t)=\lambda(\mathsf X(t))\). It
satisfies \(y''=0\), and its initial value and derivative determine
\(\lambda\). This is the local Cauchy identification with the intrinsic
optical equation. In an auxiliary conformal parameter it has the
projective optical equation and density transformation proved in
Propositions 5.3 and J.4 of the local paper. Since the identification
uses evaluation and parallel transport, it agrees on overlaps and
extends through every finite chain of leafwise continuation. It
therefore identifies representations, including their possible loop
actions.

The functional \(H(\mathsf I,\cdot)\) has value one on the affine slice
and kernel the translation line. Replacing \(\mathsf X\) by another
point in that slice leaves \(\mathsf X\wedge\Theta(k)\) unchanged. Hence
\(\det\mathbb V\) is canonically the translation line
\(\mathcal L_F^{-1}\), and its dual gives the calibrated value-slope
determinant on the optical solutions.

The kernel-pencil identity is the algebraic identity in Proposition J.4
of the local paper. For completeness, every nonzero form in the image of
the totally null plane \(\mathbb V_u\) has zero Pfaffian and rank two.
Its kernel contains \(\ell_u\) and lies in \(h_u\). A rank-two
alternating form on a four-dimensional space is determined up to scalar
by its kernel. Hence the kernel map is injective on
\(\mathbb P(\mathbb V_u)\); as a nonconstant holomorphic map between
compact projective lines it is surjective and therefore an isomorphism.
This calculation is intrinsic under the exterior-square comparison. The
transport identity (\ref{eq:global-spin-tractor-transport}) makes it
equivariant on every arrow. Evaluation at \(\mathsf X_x\) gives the
selected point plane. Since \(P_\gamma\mathsf X_x\) generally differs
from the position at the endpoint, invariance is asserted for the pencil
and not for that selected member. ∎

\pandocstatementheading{\protect\phantomsection\label{thm:global-projective-diagram-naturality}Theorem 6.3.}

The rank-four, flag, infinity, optical and point-incidence constructions
form a natural diagram on independent monodromy sources. On each
whole-sky chart it is the diagram in Theorem J.6 of the local paper. It
is natural for physical spin embedding morphisms, complete-transversal
changes and recomputed point-neighborhood restrictions. Consequently it
gives a diagram on the source germs along the entire real development
defined in \hyperref[sec:germ-correspondence-v08]{§7}, independent of
the representative complex extension under the specified real
comparison.

Normalized operator comparison isomorphisms over the same physical map
and spin label induce identity comparison on this finite-rank diagram.
The central negative spin element acts by \(-1\) on \(\mathbb L\).

\emph{Proof.} Complete-sky evaluation identifies the source data with
the canonical rank-six data on \(B\), while the spin comparisons
identify their normal spin lifts. The local comparisons of Theorem J.6
of the local paper therefore glue on the point manifold.

A holomorphic spin isometric embedding \(f:B\to B'\) induces a map
\(f_{\mathrm{LT}}\) intertwining the normal spin connections: it
preserves the Levi-Civita connection, soldering and numerical
coefficient \(c\). Initialized transport uniqueness on a finite flow-box
chain gives
\begin{equation}\protect\phantomsection\label{eq:global-projective-transport-naturality}{
f_{\mathrm{LT},v}\Pi_\gamma
=\Pi'_{F(f)\gamma}f_{\mathrm{LT},u}.
}\end{equation} This proves naturality on every path class even if the
induced groupoid functor identifies some classes. It preserves volume,
flag, infinity forms and optical planes. Evaluation and the kernel
construction also preserve the selected incidence maps. The comparisons
respect composition because the point-bundle maps and initialized
transport do.

Changing complete transversals restricts these same equivariant objects
to their common presentation. On restricting the point domain, the
groupoid is recomputed and maps to the original one by inclusion of
paths. Pullback of \(\Pi\) is exactly the transport of the restricted
connection. All the bundles and maps therefore restrict to the newly
constructed diagram, without an assertion that the restriction functor
is faithful or a Morita equivalence. The comparison of complexifications
along the entire real development and the germ correspondence proved in
\hyperref[sec:germ-correspondence-v08]{§7} then give the stated
independence.

A normalized operator comparison isomorphism preserves the same physical
symbols, torsor, root and spin label. Its finite-rank comparison is
identity by Proposition J.3 of the local paper. Transport uniqueness
propagates that identity along every arrow. This does not remove the
operator associators or their module actions. Finally the central
negative spin element acts by \(-1\) on both spinor summands. It is
invisible on their exterior square and projective flags but remains
\(-1\) on \(\mathbb L\). ∎

Assume now that \(\Lambda\) is real and that the source has compatible
Lorentz reality, with its real point spin lift. Write \(\bar u\) and
\(\bar\gamma\) for the conjugate phase point and path. When the
designated real locus is a globally hyperbolic development \(M\), its
real-path groupoid is \(\operatorname{Mon}(K_{\mathbb R})\). It consists
of paths in the real null phase, rather than all complex arrows between
real objects.

\pandocstatementheading{\protect\phantomsection\label{prop:global-projective-lorentz-duality}Proposition 6.4.}

The rank-four representation has an antilinear duality
\begin{equation}\protect\phantomsection\label{eq:global-local-twistor-real-duality}{
J_u:\mathbb L_u\longrightarrow\mathbb L_{\bar u}^*,
\qquad
J_v\Pi_\gamma=(\Pi_{\bar\gamma}^{-1})^*J_u.
}\end{equation} Its dual companion has the retained normalized square.
It exchanges \(\ell\) with \(\operatorname{Ann}(h)\), pairs a point
plane with the annihilator of its conjugate point plane, and exchanges
the covariant and contravariant infinity forms. These comparisons
commute with Theorem 6.3. On the designated real-path groupoid they
descend to a rank-four bundle with this duality over the real null-ray
space.

\emph{Proof.} The compatible point spin reality gives the local duality
in Proposition J.7 of the local paper. In normalized Lorentz spin frames
it is \(J(Z)=Z^\dagger H_{\mathrm{LT}}\), with \[
H_{\mathrm{LT}}=\begin{pmatrix}0&I_2\\I_2&0\end{pmatrix}.
\] This is a covector-valued antilinear map. The local normal spin
connection preserves it, and its dual companion has square one in these
frames, with the retained root and spin square globally. These tensor
identities glue through the spin cocycle and extend to paired complex
points. Integrating parallelism along a finite flow-box chain gives the
displayed transport equation. Consequently the local flag, point-plane
and infinity-form comparisons extend over every path class. The latter
use the volume dual of the infinity form and remain valid when
\(\Lambda=0\).

All the local identities are natural for the specified spin maps, and
initialized transport proves their compatibility with embedding and
restriction. On a globally hyperbolic real development, maximal real
null leaves are intervals. Real-path transport consequently has one
class between two points of the same real ray and descends through its
real ray submersion. This gives the final assertion for the real-path
groupoid of \hyperref[prop:real-path-groupoid-v08]{Proposition 3.9}. ∎

\pandocstatementheading{\protect\phantomsection\label{prop:operator-monodromy-lorentz-reality}Proposition 6.5.}

The operator equivalence in
\hyperref[thm:operator-monodromy-equivalence]{Theorem 4.6} retains
compatible Lorentz reality and the designated real-path groupoid. On
operator representatives, reality is a structured antilinear operation
with its specified square coherence. For a local object \(\mathcal O\)
and its conjugate \(c\mathcal O\),
\begin{equation}\protect\phantomsection\label{eq:operator-monodromy-real-square}{
\begin{aligned}
J_{c\mathcal O}\overline{J_{\mathcal O}}
&=\operatorname{Ad}(Z_{\mathcal O}),\\
J_{\mathcal O}(\overline{Z_{\mathcal O}})
&=Z_{c\mathcal O},\\
j_{c\mathcal O}\overline{j_{\mathcal O}}
&=Z_{\mathcal O}\cdot(-).
\end{aligned}
}\end{equation} Here \(Z_{\mathcal O}\) is the normalized square unit
and \(j_{\mathcal O}\) is the compatible framed module map. This
antilinear operation is distinct from the complex-linear formal adjoint.

\emph{Proof.} Apply the coupled Lorentz construction of Theorem G.5 of
the local paper on each operator chart, retaining the root square and
whole-family framing. It gives exactly the displayed equations. On a
point overlap, the conjugate comparison and the comparison of conjugates
cover the same classical Lorentz map and spin lift. The unique
normalized comparison identifies them.

Transporting a square relation across such a comparison gives two
normalized units implementing the same inner automorphism.
Normalized-inner injectivity makes them equal. Framed semilinear Hom
uniqueness gives the module equation and its coherence. The same
uniqueness on triple and quadruple overlaps makes these comparisons
compatible with both operator associators.

Choose contact transversals together with their conjugates. The
conjugate algebroid, pulled back by the antiholomorphic contact
involution, is another structured quantization of the same groupoid.
\hyperref[thm:contact-groupoid-structured-quantization]{Theorem 4.2}
gives its comparison with the original algebroid and the unique
comparison of its square with the identity. Its restrictions to
whole-sky charts agree with the preceding local construction through the
unique structured comparisons. Thus the antilinear action and its square
glue globally. The finite-rank square in Proposition 6.4 does not
strictify the chosen operator representatives.

Finally the original, conjugate and reconstructed whole-sky objects are
related by the local counits, which preserve these same reality data.
Their uniqueness makes both global inverse comparisons of Theorem 4.6
compatible with the antilinear action. All constructions retain the
designated real-path groupoid and use only its real paths for real
descent. ∎

The same equations give a weak Lorentz comparison on the incidence
module \(\mathscr M_F\): conjugation preserves the central point
parameters and sends each local incidence presentation to its conjugate
one; the framed maps \(j_{\mathcal O}\) glue by the module descent
equations of
\hyperref[thm:global-parametrized-point-module-family]{Theorem 4.8}.
Their square is the action of the retained units \(Z_{\mathcal O}\).
Central specialization sends the point object at \(x\) to that at
\(\bar x\), with its root and frame comparisons. Thus the full
parametrized point family, its finite-rank spin diagram and its operator
coherence have the same Lorentz continuation.

\protect\phantomsection\label{sec:germ-correspondence-v08}

\section{7. Germ correspondence and maximal Cauchy
development}\label{germ-correspondence-and-maximal-cauchy-development}

Fix the analytic Cauchy data of \hyperref[sec:local-input-v08]{§2},
including the orientations and prescribed initial spin structure. Their
development category is \hyperref[def:analytic-ghd-v2]{Definition A.1}.
We now remove the choice of complex point extension from the independent
source-atlas, operator and monodromy correspondences.

An \emph{extension object} has a connected point domain \((B,\sigma)\)
whose fixed locus is exactly a connected real four-manifold \(M\).
Require compatible antiholomorphic reality on its incidence, groupoid
and quadratic-affine data, and actual spin reality on its whole-sky
charts. The chart reality must have the Lorentz real form described in
§2 and in §8.4 of the local paper: its retained spin/tangent calibration
at each real mark has tangent form
\(\eta=\operatorname{diag}(1,-1,-1,-1)\) in the selected Einstein scale.
By \hyperref[lem:independent-ray-lorentz-recognition-v092]{Lemma 2.4},
the locally recognized metrics are real with this signature, the fixed
incidence \(F_{\mathbb R}\to M\) has real two-sphere fibers, and the
lines \(\mathrm dp(K_{\mathbb R})\) form their null cone. Supply a time
orientation of that cone, an orientation of \(M\), and an analytic
embedding \(\iota:\Sigma\to M\) that is spacelike and Cauchy for its
causal structure. In particular every inextendible real causal curve
must meet \(\iota(\Sigma)\) exactly once. Require that local recognition
along this embedding gives the prescribed forms \((h,K)\) and
orientation convention of \hyperref[def:analytic-ghd-v2]{Definition
A.1}, together with an actual identification of \(P_\Sigma\) with the
initial adapted-frame spin reduction. These are conditions on the
supplied source and its local models; no global Einstein metric is
supplied. Morphisms preserve all these markings. The designated
real-path groupoid is \(\operatorname{Mon}(K_{\mathbb R})\), with its
intrinsic functor into the complex monodromy groupoid, as in
\hyperref[prop:real-path-groupoid-v08]{Proposition 3.9}.

For an original-source atlas, an extension structure means the preceding
reality and initial markings on the phase constructed in
\hyperref[lem:source-atlas-phase-completion]{Lemma 5.5}, with the local
whole-sky models and their physical overlap comparisons compatible with
that reality. The Lorentz operation on the original pair is the
residue-transferred operation in §8.4 of the local paper; its two
classes retain their weights \(2\) and \(-6\). For an operator source,
use its supplied phase and the same markings, requiring the structured
local comparisons and their descent maps to preserve them. On operator
representatives the induced reality has the normalized inner square,
framed module square and overlap coherence of
\hyperref[prop:operator-monodromy-lorentz-reality]{Proposition 6.5}.
These are equations on structured representatives, considered through
their normalized comparisons. Morphisms in both extension categories are
their physical morphisms preserving this data, including the specified
real spin and initial identifications. Thus the extension structures are
defined on the independently constructed source phases before global
metric recognition.

An \emph{allowed point restriction} is a connected conjugation-invariant
open neighborhood of all of \(M\) in \(B\). Every such restriction
retains the entire real development, its initial hypersurface, and the
full complex quadric over every retained point.

\pandocstatementheading{\protect\phantomsection\label{prop:recomputed-monodromy-restriction}Proposition 7.1.}

Let \(\mathcal S\) be an extension monodromy source on \(B\), and let
\(B_0\subset B\) be an allowed point neighborhood. Its canonical
recomputed restriction \(\mathcal S^{\mathrm{mon}}_{B_0}\) has point
space \(B_0\), incidence \(F_0=p^{-1}(B_0)\) and groupoid
\begin{equation}\protect\phantomsection\label{eq:recomputed-monodromy-restriction}{
\mathcal G_0=\operatorname{Mon}(K|_{F_0}).
}\end{equation} The remaining data restrict to \(F_0\), with action
pulled back along the integrated inclusion
\(\mathcal G_0\to\mathcal G\). The result is an admitted extension
source with the same real development and markings, and inclusion gives
a source morphism
\(r_{B_0}:\mathcal S^{\mathrm{mon}}_{B_0}\to\mathcal S\). This operation
is canonically compatible with further restrictions and with source
morphisms whose point maps take the smaller source domain into the
smaller target domain.

\emph{Proof.} The supplied characteristic line restricts without using a
metric. Its monodromy integration is source connected and source simply
connected. The Lie algebroid inclusion integrates uniquely into the
supplied source groupoid, by the uniqueness of integration used in
\hyperref[prop:holomorphic-null-monodromy]{Proposition 3.7}. Pullback
along this functor gives the vector and affine actions. Whole-sky
conditions persist, since point restriction retains every full quadric;
properness of \(p\) is stable under this base restriction. The
short-path compatibility equations restrict and are preserved by
composition of the integrated inclusions. Reality and all initial
markings remain unchanged because every real point is retained.

For nested restrictions, the integrated functors send a path to the same
path in the larger domain. They agree by uniqueness of integration, as
do their representation comparisons. A source morphism sends paths and
their homotopies in a smaller source domain to paths and homotopies in
the smaller target domain whenever the stated image condition holds.
This proves functoriality. ∎

The groupoid \(\mathcal G_0\) is generally not the full subgroupoid of
\(\mathcal G\) on \(F_0\). Such a full subgroupoid includes paths
leaving \(F_0\), while homotopies in the larger leaf may identify
classes that remain distinct in the smaller leaf. The integrated
inclusion need not be faithful or a Morita equivalence.

Operator restriction simultaneously restricts symbols, affine torsors,
whole-family modules, roots, frames, associators and their comparison
maps, and pulls the groupoid action back along the same integrated
inclusion. All their defining equations persist. The structured
comparison of \hyperref[cor:null-monodromy-first-unit-v08]{Corollary
4.3} identifies this with the quantization of the new contact monodromy
presentation. Original-source restriction restricts every point chart
and physical comparison germ.
\hyperref[lem:source-atlas-phase-completion]{Lemma 5.5} identifies the
resulting completed groupoid with \(\operatorname{Mon}(K|_{F_0})\). The
fixed-domain source-atlas and operator equivalences therefore commute
with these restrictions through their inverse comparisons.

\pandocstatementheading{\protect\phantomsection\label{def:entire-real-source-germs-v08}Definition 7.2.}

The category \(\mathsf{MonRay}^{\mathrm{germ}}_\Lambda(\Sigma,h,K)\) has
the extension monodromy sources as objects. A morphism from
\(\mathcal S\) on \((B,M)\) to \(\mathcal S'\) on \((B',M')\) is
represented by an allowed neighborhood \(B_0\subset B\) and an enriched
source morphism \(f:\mathcal S^{\mathrm{mon}}_{B_0}\to\mathcal S'\). Its
real point map must preserve orientation, time orientation, the initial
embedding and the prescribed initial spin identification. Two
representatives define the same morphism when they agree after
recomputed restriction to a common allowed neighborhood of all of \(M\).
Equality includes the spin lift. Equivalently, a morphism is a roof
\begin{equation}\protect\phantomsection\label{eq:monodromy-germ-morphism-roof}{
\mathcal S
\xleftarrow{\ r_{B_0}\ }
\mathcal S^{\mathrm{mon}}_{B_0}
\xrightarrow{\ f\ }
\mathcal S',
}\end{equation} modulo common point restriction. The source-atlas and
operator categories
\(\mathsf{SrcAtl}^{\mathrm{germ}}_\Lambda(\Sigma,h,K)\) and
\(\mathsf{OpMon}^{\mathrm{germ}}_\Lambda(\Sigma,h,K)\) are defined by
the same construction with their respective physical morphisms and
restrictions. Atlas refinement and normalized operator comparison have
the conventions of
\hyperref[def:independent-original-source-atlas]{Definition 5.1} and
\hyperref[def:independent-operator-monodromy-source-v092]{Definition
4.4}. In particular operator equality is equality of the physical Hom
classes with their specified normalized comparisons, while the full spin
labels remain distinct.

For composition, if the second map is represented on \(B'_0\subset B'\),
restrict the first domain to the component containing \(M\) of
\begin{equation}\protect\phantomsection\label{eq:monodromy-germ-composition-domain}{
B_0\cap f^{-1}(B'_0).
}\end{equation} This component contains all of connected \(M\) and is
invariant under conjugation. Recompute the source groupoid there, use
the induced functor into the recomputed groupoid over \(B'_0\), and
compose. Common restriction proves independence of representatives and
associativity; identity source maps give identities. The restriction
arrow \(r_{B_0}\) is invertible in this category, with inverse
represented by the identity on the restricted object viewed as a germ
from the original object.

These categories are precisely the localizations by the specified
restriction arrows. Indeed, any functor \(L\) on extension objects that
inverts those arrows must send a roof to \(L(f)L(r_{B_0})^{-1}\). Common
restriction makes this expression well defined, and the composition
formula makes it functorial. Every germ arrow has a roof, so this
factorization is unique. This is a neighborhood localization of enriched
sources, defined before global metric recognition; it does not declare
the integrated restriction functors to be groupoid equivalences.

\pandocstatementheading{\protect\phantomsection\label{lem:whole-real-isometry-complexification}Lemma 7.3.}

Let \((B,g,\sigma)\) and \((B',g',\sigma')\) be compatible holomorphic
oriented spin extensions with fixed loci exactly \(M\) and \(M'\). Every
real analytic spin isometric embedding \(f:M\to M'\) has a unique germ
along all of \(M\) of a compatible holomorphic spin isometric embedding
into \(B'\). This allows noncompact \(M\) and nonproper \(f\), and
preserves composition and full spin lifts.

\emph{Proof.} Equal dimensions make \(f(M)\) open in \(M'\), with \(f\)
an analytic diffeomorphism onto its image. The fixed locus \(M'\) is
closed in the Hausdorff manifold \(B'\), and \(M'\setminus f(M)\) is
therefore closed in \(B'\). The invariant open complex manifold \[
\Omega=B'\setminus\bigl(M'\setminus f(M)\bigr)
\] has fixed locus exactly \(f(M)\). Via \(f\), the manifolds \(B\) and
\(\Omega\) are complexifications of the same real analytic manifold.
Germ uniqueness of complexification (Whitney and Bruhat 1959, sec. 1,
Proposition 1) gives a biholomorphism between neighborhoods of all of
\(M\) and all of \(f(M)\) extending \(f\). This gives injectivity on a
suitable source neighborhood, including when the real embedding is
nonproper.

The metric tensors agree on the maximal totally real slice and hence, by
holomorphic uniqueness, on a neighborhood of that entire slice.
Comparing the extension with its conjugate makes it equivariant after an
invariant restriction. The spin lift extends by the local holomorphic
double-cover construction of
\hyperref[prop:einstein-complexification-phase]{Proposition 3.1}. Its
sheet choices agree on real overlaps and therefore on smaller complex
overlaps, retaining the prescribed sign. Connected neighborhoods
containing all of \(M\) can be retained throughout.

Two extensions have equal holomorphic germs by uniqueness on the totally
real slice. Their spin lifts have the same sheet choices there, so their
germs also agree. Applying this uniqueness to composites proves the
composition assertion. The argument requires no uniform complex radius.
∎

\pandocstatementheading{\protect\phantomsection\label{thm:source-atlas-cauchy-equivalence-v08}Theorem 7.4.}

Under the local inputs of \hyperref[sec:local-input-v08]{§2}, real
recognition gives equivalences
\begin{equation}\protect\phantomsection\label{eq:source-atlas-cauchy-equivalence-v08}{\begin{aligned}
\mathsf{SrcAtl}^{\mathrm{germ}}_\Lambda(\Sigma,h,K)&\simeq
\mathsf{OpMon}^{\mathrm{germ}}_\Lambda(\Sigma,h,K) \\
&\simeq
\mathsf{MonRay}^{\mathrm{germ}}_\Lambda(\Sigma,h,K) \\
&\simeq
\mathsf{Ein}^{\omega}_\Lambda(\Sigma,h,K).
\end{aligned}}\end{equation} They retain the proper complete four-dimensional point
family, each full complex sky, the quadratic-affine representation, the
Einstein scale, Lorentz reality, the designated real-path groupoid and
the spin and initial identifications. On each whole-sky chart they agree
with the original local correspondence, including its two input classes,
source-operation comparisons and normalized operator and module data.
The projective spin diagram and parametrized incidence-module family are
carried by the same natural comparisons. Two compatible complex
extensions of one real development give canonically isomorphic source
germs under the specified real spin identification.

\emph{Proof.} First consider monodromy sources.
\hyperref[thm:independent-monodromy-equivalence]{Theorem 3.11}
recognizes an independent source as a holomorphic Einstein manifold and
identifies its source morphisms with holomorphic spin isometric
embeddings. By
\hyperref[lem:independent-ray-lorentz-recognition-v092]{Lemma 2.4},
local recognition intertwines the supplied conjugation and spin reality,
and \(\sigma^*\overline g=g\). Its fixed metric on \(M\) is real, and
the retained Lorentz tangent calibration gives signature \((+,-,-,-)\).
Its null directions are precisely \(\mathrm dp(K_{\mathbb R})\). The
admission conditions therefore make \(\iota(\Sigma)\) Cauchy for this
metric and give the prescribed \((h,K)\), orientations and initial spin
identification. Thus its real restriction is an analytic spin
development. Recomputed restriction changes none of its real data. Real
recognition defines the required functor to
\(\mathsf{Ein}^{\omega}_\Lambda(\Sigma,h,K)\).

Every analytic spin development has a compatible holomorphic extension
with precisely that real fixed locus, by
\hyperref[prop:einstein-complexification-phase]{Proposition 3.1}. The
canonical monodromy source on this extension is independently admitted
by \hyperref[thm:independent-monodromy-equivalence]{Theorem 3.11}, and
its recognition returns the original development. This proves essential
surjectivity.

For fullness, take a development morphism \(f:M\to M'\), with its spin
lift. \hyperref[lem:whole-real-isometry-complexification]{Lemma 7.3}
extends it to a holomorphic spin isometric embedding from a neighborhood
of all of \(M\) into the target extension. Fixed-domain full
faithfulness then supplies the corresponding enriched source morphism on
the recomputed source domain. Its germ recognizes to \(f\). The induced
groupoid functor sends each path class to the class of its image; no
injectivity on arrows is required.

For faithfulness, suppose two source-morphism germs recognize to the
same real spin embedding. The lemma makes their holomorphic point maps
and spin lifts agree on a common neighborhood of all of \(M\).
Fixed-domain full faithfulness identifies their enriched morphisms on
the recomputed source there. This comparison controls every path in that
smaller domain: each path is a finite composite of short paths, and the
source connected and source simply connected integration has a unique
functor inducing the fixed infinitesimal comparison. Thus the source
germs agree. The monodromy recognition functor is fully faithful and
essentially surjective.

Now use \hyperref[thm:operator-monodromy-equivalence]{Theorem 4.6} and
\hyperref[thm:original-source-atlas-operator-equivalence]{Theorem 5.6}
on the fixed point domains. Their local reality comparisons are those of
Proposition 8.4 and Theorems 9.10 and G.5 of the local paper. The source
cocycles and full spin lifts make these comparisons coherent on
overlaps. On operator representatives they preserve the normalized inner
square and the module square, with the supplied associators present; the
original class-level involution does not change these into strict
equations on arbitrary representatives. The designated real-path
groupoid is the integration of the real characteristic line throughout.

Both fixed-domain equivalences commute with recomputed restriction, as
established after
\hyperref[prop:recomputed-monodromy-restriction]{Proposition 7.1}. They
consequently induce functors on the roof categories. To verify fullness,
represent a target germ arrow after one allowed point restriction and
apply the fixed-domain full Hom correspondence on that domain. For
faithfulness, equality of two target germs is witnessed after a further
common point restriction; fixed-domain faithfulness then identifies
their lifted physical Hom classes there. For essential surjectivity,
construct or decode a representing object on its point domain by the
fixed-domain inverse. Transport the supplied phase reality and initial
markings through these inverse comparisons; their local reality
compatibility and the normalized operator square of Proposition 6.5 give
exactly the extension structures specified above. This proves both
remaining equivalences without any separate assumption about groupoid
equivalence of restriction arrows.

The inverse comparisons used here are the local evaluation,
source-operation and operator-module comparisons, extended over paths by
unique integration and glued with their specified associators. Their
naturality follows from the established Hom bijections.
\hyperref[prop:source-atlas-complete-diagram]{Proposition 5.7},
\hyperref[thm:global-projective-diagram-naturality]{Theorem 6.3}, and
\hyperref[thm:global-parametrized-point-module-family]{Theorem 4.8} show
that the indicated projective and parametrized-module data are carried
by these same comparisons. They commute with restriction and therefore
pass to the germ categories.

Finally, apply fullness to the identity real spin isometry between two
complex extensions, and faithfulness to the composites of the resulting
arrows. This gives their unique source-germ isomorphism under the
specified real identification. All restrictions retain every real point
and its whole complex sky, so the equivalence retains the entire
development. ∎

\pandocstatementheading{\protect\phantomsection\label{cor:maximal-source-atlas-germ-v08}Corollary 7.5.}

Each source germ category in
\hyperref[thm:source-atlas-cauchy-equivalence-v08]{Theorem 7.4} has a
terminal object whose recognized real point development is \(M_\omega\).
It is unique up to the unique initial-data-preserving physical spin
isomorphism. An analytic isomorphism of the initial data, with its
specified spatial spin lift, induces the unique corresponding
isomorphism of maximal source germs, compatibly with identities and
composition.

\emph{Proof.} By \hyperref[cor:analytic-cauchy-naturality-v08]{Corollary
A.8}, the prescribed analytic constraint data have the terminal analytic
spin development \(M_\omega\). Choose a compatible holomorphic extension
and a corresponding object in each source category. Any other source
germ recognizes to an analytic development \(N\). Its unique development
morphism \(N\to M_\omega\) lifts to a unique physical source-germ arrow
by the full faithfulness of
\hyperref[thm:source-atlas-cauchy-equivalence-v08]{Theorem 7.4}. This is
the terminal property, and it also proves independence of the extension
and source representatives.

For an initial-data isomorphism, transport the initial embedding and
spin identification as in
\hyperref[cor:analytic-cauchy-naturality-v08]{Corollary A.8}. Its unique
maximal real spin isometry lifts uniquely through the equivalences.
Uniqueness proves the identity and composition laws. These are the same
comparisons that preserve the original source operations and global
point-object diagram. In particular they preserve curvature tensors and
all their derivatives; no independently specified abstract Cauchy record
is required. ∎

The terminal object belongs to the analytic Cauchy germ category of
\hyperref[rem:analytic-maximality-scope]{Remark A.9}. Its source
presentation retains the full point family, spin lifts and parametrized
module under the restriction conventions of
\hyperref[prop:recomputed-monodromy-restriction]{Proposition 7.1} and
\hyperref[def:entire-real-source-germs-v08]{Definition 7.2}.

\protect\phantomsection\label{sec:nariai-obstructions-v08}

\section{8. Nariai obstructions}\label{nariai-obstructions}

The equal-radius Nariai development shows why the preceding
correspondence retains a monodromy presentation and parametrized
incidence. First, no Hausdorff holomorphic ray carrier can contain all
complex null directions over its entire real development. Second, in its
product complexification the exact ray image of a point is not the
support of an ordinary locally finitely generated global module, even
after restriction around all real points. The first assertion applies to
every complexification. The second holds for every real point on every
allowed product neighborhood and hence on a cofinal collection of
representatives of the germ of any compatible complexification.

Let \(a>0\) and set
\begin{equation}\protect\phantomsection\label{eq:nariai-product-complexification}{
\begin{aligned}
H_1&=\mathrm dx_0^2-\mathrm dx_1^2-\mathrm dx_2^2,
&H_2&=-\mathrm dy_1^2-\mathrm dy_2^2-\mathrm dy_3^2,\\
Q_i&=\{v:H_i(v,v)=-a^2\},
&B_0&=Q_1\times Q_2,\qquad
g_0=H_1|_{TQ_1}+H_2|_{TQ_2}.
\end{aligned}
}\end{equation} Write \(M\) for the conjugation-fixed locus. The product
geometry is standard (Beyer 2009, sec. 3.2.1); the obstructions are
proved below.

\pandocstatementheading{\protect\phantomsection\label{prop:nariai-maximal-einstein-development}Proposition 8.1.}

The real metric is
\begin{equation}\protect\phantomsection\label{eq:nariai-real-metric}{
M=\mathbb R\times S^1\times S^2,
\qquad
g=\mathrm dt^2-a^2\cosh^2(t/a)\,\mathrm d\theta^2-a^2g_{S^2}.
}\end{equation} It is analytic, geodesically complete and globally
hyperbolic, with \(\operatorname{Ric}(g)=\Lambda g\) for
\(\Lambda=-a^{-2}\) in our curvature convention, and with nonzero Weyl
tensor. Its compact Cauchy data at \(t=0\) are
\(h=a^2(\mathrm d\theta^2+g_{S^2})\) and \(K=0\). With either prescribed
spin structure on \(S^1\times S^2\), it is the maximal analytic spin
Cauchy development of these data.

\emph{Proof.} Parametrize the real factors by \[
\begin{gathered}
x_1=(a\sinh(t/a),a\cosh(t/a)\cos\theta,
a\cosh(t/a)\sin\theta), \\
x_2=a\omega,
\quad\omega\in S^2.
\end{gathered}
\] The induced metric is the displayed one. On each factor, \[
R_i(v,w)u=-a^{-2}\bigl(g_i(w,u)v-g_i(v,u)w\bigr).
\] Since both factors have dimension two, their Ricci tensors are
\(-a^{-2}g_i\), proving the Einstein equation. Mixed product curvature
vanishes, whereas for nonzero \(v\in TQ_1\) and non-null \(w\in TQ_2\),
\begin{equation}\protect\phantomsection\label{eq:nariai-nonzero-weyl}{
W(v,w)w=-\frac{\Lambda}{3}g(w,w)v\ne0.
}\end{equation} The function \(t\) is temporal. On every bounded time
interval the spatial metric is bounded below by a fixed positive metric
on the compact spatial manifold. A causal curve parametrized by \(t\)
thus has bounded spatial speed. A finite endpoint of its time range
would give a spacetime endpoint and an extension. Every inextendible
causal curve therefore meets every time slice exactly once. At zero the
derivative of the spatial metric vanishes, giving \(K=0\).

Each affine factor geodesic satisfies
\begin{equation}\protect\phantomsection\label{eq:nariai-factor-geodesic-equation}{
x_i''=(\varepsilon_i/a^2)x_i,
\qquad\varepsilon_i=g_i(x_i',x_i')\ \text{constant}.
}\end{equation} The linear equation has solutions for all real
parameters and preserves the quadric and tangency constraints. All
product geodesics are complete. Such a manifold has no proper connected
isometric open extension: choose a convex normal neighborhood of a
boundary point of a proposed image and a point of the image in that
neighborhood. The ambient geodesic joining them must be the image of the
complete source geodesic with the same initial data, placing the
boundary point in the image, a contradiction. Its embedding supplied by
the \hyperref[thm:maximal-analytic-development-v2]{maximal analytic
development theorem (Theorem A.7)} is consequently onto. Finally,
\(S^1\times S^2\) has two spin structures, parametrized by
\(H^1(S^1\times S^2;\mathbb Z_2)=\mathbb Z_2\), and each extends over
the product with \(\mathbb R\). The spin comparison in that theorem
gives the stated identification. ∎

\pandocstatementheading{\protect\phantomsection\label{thm:nariai-no-hausdorff-whole-sky-carrier}Theorem 8.2.}

Let \((B,g_B)\) be any holomorphic extension of this real Nariai metric,
and let \(F\to B\) be its full projective null bundle with line field
\(K\). There is no triple consisting of an open subset \(U\subset F\)
containing the entire quadric \(F_x\) for every \(x\in M\), a Hausdorff
complex manifold \(Z\), and a holomorphic map
\begin{equation}\protect\phantomsection\label{eq:nariai-forbidden-carrier-map}{
q:U\longrightarrow Z,
\qquad\ker\mathrm dq=K|_U.
}\end{equation} The obstruction is independent of projected-sky
embeddedness, affine enrichment and spin structure.

\emph{Proof.} At a real point, the tangent space of every
complexification identifies with the complexified real tangent space. We
can therefore use the following directions over real points without
choosing any nonreal extension. Fix \[
x_1=(0,a,0),\qquad x_2=ae_3,\qquad x=(x_1,x_2).
\] For a real parameter \(\eta\) near zero put
\begin{equation}\protect\phantomsection\label{eq:nariai-endpoint-family}{
\begin{aligned}
k_\eta&=(a\cosh\eta,0,a\sinh\eta),\\
y_{1,\eta}&=x_1\cosh\pi+k_\eta\sinh\pi,\\
k_\eta'&=x_1\sinh\pi+k_\eta\cosh\pi,
\qquad y_\eta=(y_{1,\eta},-x_2).
\end{aligned}
}\end{equation} The initial and final first-factor velocities have
squared length \(a^2\). For \(z\in\mathbb C^*\) let \[
u(z)=\frac a2\bigl((z+z^{-1})e_1-i(z-z^{-1})e_2\bigr).
\] This is tangent at \(x_2\), has squared length \(-a^2\), and is real
when \(|z|=1\). For these real directions,
\begin{equation}\protect\phantomsection\label{eq:nariai-antipodal-null-geodesics}{
\gamma_{\eta,z}(s)=
\bigl(x_1\cosh s+k_\eta\sinh s,
      x_2\cos s+u(z)\sin s\bigr),
\qquad 0\le s\le\pi,
}\end{equation} is an affine null geodesic from \(x\) to \(y_\eta\),
with initial and final velocities \((k_\eta,u(z))\) and
\((k_\eta',-u(z))\). Its phase lift lies in \(U\), so its endpoints have
the same image under \(q\).

These endpoint directions extend over the whole parameter curve
\(\mathbb P^1\). Put \[
v(S,T)=a(S^2+T^2)e_1-ia(S^2-T^2)e_2
\] and define
\begin{equation}\protect\phantomsection\label{eq:nariai-compact-direction-conics}{
\begin{aligned}
\alpha_\eta([S:T])&=
 \bigl(x,[2STk_\eta,v(S,T)]\bigr),\\
\beta_\eta([S:T])&=
 \bigl(y_\eta,[2STk_\eta',-v(S,T)]\bigr).
\end{aligned}
}\end{equation} The vectors never vanish and their two squared lengths
are \(4a^2S^2T^2\) and its negative. Thus these are holomorphic maps
into the full fibers retained in \(U\). For \(z=S/T\) on the unit
circle, the real geodesics give \(q\alpha_\eta=q\beta_\eta\). The
diagonal of the Hausdorff complex manifold \(Z\) is a closed analytic
subset. Hence the equalizer is a closed analytic subset of
\(\mathbb P^1\) containing the unit circle, and is all of
\(\mathbb P^1\).

At \([S:T]=[1:0]\), let \(v_+=a(e_1-ie_2)\). Projective sign
cancellation gives
\begin{equation}\protect\phantomsection\label{eq:nariai-forced-degenerate-ray-equality}{
q\bigl(x,[0,v_+]\bigr)
=q\bigl(y_\eta,[0,v_+]\bigr)
\qquad(\eta\text{ real near }0).
}\end{equation} Differentiate the right-hand side at zero, obtaining
\(w\in TF\) with \(\mathrm dq(w)=0\) and \[
\mathrm dp(w)=\bigl(a\sinh\pi\,e_2,0\bigr)\ne0,
\] where \(e_2\) in the first factor is its third ambient coordinate
vector. At \((y_0,[0,v_+])\), however, \(\mathrm dp(K)\) is the pure
second-factor line \(\mathbb C(0,v_+)\). The condition
\(\ker\mathrm dq=K\) is incompatible with this nonzero first-factor
derivative. ∎

In particular, no such extension has a closed embedded null-leaf
relation, by \hyperref[lem:ordinary-null-quotient-v08]{Lemma 3.12}. The
obstruction already precedes the descent of quadratic or operator data.
The real affine incidence and the monodromy construction of §3 remain
defined.

To analyze supports, now use the particular product extension \(B_0\),
its phase \(F\), and its maximal leaf relation \(R\). The following
explicit geometry controls every restriction of this extension around
the entire real development.

\pandocstatementheading{\protect\phantomsection\label{prop:nariai-proper-leaves-nonclosed-relation}Proposition 8.3.}

Every complex null leaf of \(B_0\) is a properly embedded copy of
\(\mathbb C\). Its partial tractor connection has trivial holonomy and
its affine position is injective. Nevertheless, \(R\) is not closed in
\(F\times F\).

\emph{Proof.} For initial null tangent \((k_1,k_2)\) at \((x_1,x_2)\)
write \(\varepsilon_i=H_i(k_i,k_i)\), so
\(\varepsilon_2=-\varepsilon_1\). If \(\varepsilon_1=0\), the affine
solution is the nonconstant proper embedded line
\((x_1+sk_1,x_2+sk_2)\). Otherwise choose \(r\) with
\(r^2=\varepsilon_1/a^2\). The entire solution is
\begin{equation}\protect\phantomsection\label{eq:nariai-complex-null-geodesic}{
\gamma(s)=
\bigl(x_1\cosh(rs)+r^{-1}k_1\sinh(rs),
      x_2\cos(rs)+r^{-1}k_2\sin(rs)\bigr).
}\end{equation} In each factor \(x_i,k_i\) are independent: they are
orthogonal and \(x_i\) has nonzero squared length. A return to the
starting point forces both \(rs\in2\pi i\mathbb Z\) and
\(rs\in2\pi\mathbb Z\), hence \(s=0\). Shifting the affine origin proves
injectivity. A bounded set of image points bounds
\(\cosh(rs),\sinh(rs),\cos(rs),\sin(rs)\), which bounds both real and
imaginary parts of \(rs\). This proves properness. The phase lift is
also proper and injective. Its tangent spans \(K\), and its image is
both open and closed in its connected leaf; it is therefore the entire
leaf. Simple connectivity gives trivial partial holonomy. In a parallel
tractor frame its affine position is \(\mathsf X_0+sV\) with \(V\ne0\),
proving injectivity.

For nonclosedness fix \(x_1=(0,a,0)\), \(x_2=ae_3\), \(u_2=ae_1\), and
for positive \(r\to0\) put
\begin{equation}\protect\phantomsection\label{eq:nariai-nonclosed-sequence}{
k_{1,r}=\frac a2(1+r^2,0,1-r^2),
\qquad k_{2,r}=ru_2,
\qquad s_r=2\pi i/r.
}\end{equation} The tangent is null with \(\varepsilon_1=a^2r^2\). At
time \(s_r\) the first factor and its tangent return to \(x_1,k_{1,r}\),
while the second factor and tangent are \[
y_2=x_2\cosh(2\pi)+iu_2\sinh(2\pi),
\qquad r\bigl(u_2\cosh(2\pi)-ix_2\sinh(2\pi)\bigr).
\] The related initial and final phase points converge to \[
\bigl((x_1,x_2),[k_{1,0},0]\bigr),\qquad
\bigl((x_1,y_2),[k_{1,0},0]\bigr),
\qquad k_{1,0}=\tfrac a2(1,0,1).
\] Their pure first-factor null leaves keep the second point constant.
Since \(y_2\ne x_2\), the limits are unrelated. ∎

\pandocstatementheading{\protect\phantomsection\label{prop:nariai-restricted-quotient-v08}Proposition 8.4.}

Let \(B'\subset B_0\) be any connected conjugation-invariant open
neighborhood of all of \(M\), and put \(F'=F|_{B'}\). Let \(R'\) relate
points in the same connected component of an original leaf's
intersection with \(F'\). Then \(R\) and \(R'\) are embedded, and
\begin{equation}\protect\phantomsection\label{eq:nariai-restricted-point-quotient}{
q':F'\longrightarrow Z'=F'/R'
}\end{equation} is an open holomorphic submersion onto a non-Hausdorff
contact fivefold. The quadratic and affine data descend to \(Z'\), and
affine position embeds \(F'\) as a Hausdorff open incidence domain in
the affine bundle. Its lifted full-quadric point family is proper and
locally complete in this Hausdorff incidence domain. For either
prescribed real spin structure, \(B'\) can be chosen so that the
prescribed spin structure extends.

\emph{Proof.} We first prove the following separation property, for
every \(x\in B_0\),
\begin{equation}\protect\phantomsection\label{eq:nariai-same-point-relation-closure}{
\overline R\cap(F_x\times F_x)=\Delta_{F_x}.
}\end{equation} Let \((u_n,v_n)\in R\) converge to
\((u,v)\in F_x\times F_x\). Normalize representatives \(k_n\) of the
initial directions by a fixed linear functional nonzero at the limiting
direction. Thus \(k_n\to k\ne0\), and the initial and final base points
\(x_n,y_n\) tend to \(x\). Proposition 8.3 gives a unique affine
parameter \(s_n\) from \(u_n\) to \(v_n\).

For nonzero first-factor energy, set \(r_n^2=H_1(k_{1,n},k_{1,n})/a^2\)
and \(z_n=r_ns_n\). Taking the ambient inner products of the endpoint
equations with the initial points gives \[
\cosh z_n=-a^{-2}H_1(x_{1,n},y_{1,n})\longrightarrow1,
\qquad
\cos z_n=-a^{-2}H_2(x_{2,n},y_{2,n})\longrightarrow1.
\] These bounds bound \(\operatorname{Re}z_n\) and
\(\operatorname{Im}z_n\), respectively. Every limit is in
\(2\pi i\mathbb Z\cap2\pi\mathbb Z=\{0\}\), so \(z_n\to0\). Choose a
factor with \(k_i\ne0\). Its endpoint equation is \[
\begin{gathered}
y_{i,n}-x_{i,n}f_i(z_n)=k_{i,n}s_nh_i(z_n), \\
(f_1,h_1)=(\cosh,\sinh(z)/z),\quad
(f_2,h_2)=(\cos,\sin(z)/z).
\end{gathered}
\] Both removable values \(h_i(0)\) equal one. Apply a fixed linear
functional nonzero on \(k_i\) to see \(s_n\to0\). For zero energy the
same follows directly from \(y_n-x_n=s_nk_n\). Short-flow continuity now
gives \(v=u\), proving the separation property.

In particular, related pairs approaching a diagonal pair have joining
parameter tending to zero. In a flow box their relation is therefore
exactly the embedded local leaf relation, with no additional branches.
Finite compositions of local leafwise flows in one factor transfer this
to every related pair, so \(R\) is embedded.

The relation \(R'\) is open in \(R\cap(F'\times F')\). A path connecting
a related pair inside \(F'\) is compact and has a finite flow-box cover
there; its endpoint comparison persists for nearby related pairs. Hence
\(R'\) is embedded too. Lemma 3.12 gives the quotient charts, open
submersion and contact structure. This quotient cannot be Hausdorff by
Theorem 8.2.

Every restricted leaf is an open subset of an original affine-parameter
line. Its partial connection extends over that entire line, so all
restricted loop transports are trivial, and its affine position stays
injective. Lemma 3.12 therefore descends the bundles and embeds \(F'\)
into the affine bundle. Properness of the lifted point family is that of
\(F'\to B'\). In this incidence domain its fiber has normal bundle
\(T_xB'\otimes\mathcal O_Q\). To see local completeness there directly,
a small compact deformation stays over a coordinate neighborhood of
\(x\); each coordinate of its point projection is holomorphic on a
compact connected manifold and is constant. Equal dimensions force the
deformation to be the entire fiber over that constant point. Evaluation
on a local section shows that this point depends holomorphically in a
holomorphic family.

Finally, Proposition 3.1 extends either specified real spin structure to
a holomorphic metric extension. Germ uniqueness identifies a smaller
conjugation-invariant neighborhood of all of \(M\) with one in \(B_0\),
preserving metric and reality. The selected spin structure thus extends
to a suitable \(B'\). ∎

For \(x\in B'\) let \(S_x=q'(F_x)\). An embedding into a possibly
non-Hausdorff manifold is an injective holomorphic immersion that is a
homeomorphism onto its image; the image need not be closed.

\pandocstatementheading{\protect\phantomsection\label{lem:nariai-individual-sky-embedding}Lemma 8.5.}

The map \(q'|_{F_x}\) is a holomorphic embedding and \(S_x\) is
Hausdorff in its subspace topology.

\emph{Proof.} By the separation property
(\ref{eq:nariai-same-point-relation-closure}), distinct \(u,v\in F_x\)
have neighborhoods \(U,V\subset F'\) with
\((U\times V)\cap R'=\varnothing\). The open sets \(q'(U),q'(V)\) are
disjoint, making \(S_x\) Hausdorff. Proposition 8.3 shows that no
original null geodesic returns to its base point, so \(q'|_{F_x}\) is
injective. It is a continuous bijection from a compact quadric to its
Hausdorff image, hence a homeomorphism. Its differential is injective
because \(K\cap\ker\mathrm dp=0\). ∎

\pandocstatementheading{\protect\phantomsection\label{thm:nariai-all-real-nonclosed-point-skies-v09}Theorem 8.6.}

For every neighborhood \(B'\) in Proposition 8.4 and every real point
\(x\in M\), the compact embedded sky \(S_x\) is not closed in \(Z'\).
Equivalently,
\begin{equation}\protect\phantomsection\label{eq:nariai-point-saturation}{
\operatorname{Sat}_{R'}(F_x)
=\{u\in F':\text{the restricted leaf through }u\text{ meets }F_x\}
=q'^{-1}(S_x)
}\end{equation} is not closed in the Hausdorff phase manifold \(F'\).

\emph{Proof.} First take \(x=((0,a,0),ae_3)\). Set \(k_0=(a,0,0)\) and
use the real endpoint and velocity \[
y_1=x_1\cosh\pi+k_0\sinh\pi,
\qquad k'=x_1\sinh\pi+k_0\cosh\pi,
\qquad y=(y_1,-x_2).
\] The direction conics from Theorem 8.2 at \(\eta=0\) are
\begin{equation}\protect\phantomsection\label{eq:nariai-point-support-direction-conics}{
\alpha([S:T])=(x,[2STk_0,v(S,T)]),
\qquad
\beta([S:T])=(y,[2STk',-v(S,T)]).
}\end{equation} They lie in \(F'\) because their entire fibers are over
real points. For \(|S/T|=1\) the real geodesic between their endpoints
stays in \(M\), hence in \(B'\), so \(q'\beta=q'\alpha\in S_x\) there.

If \(S_x\) were closed, Lemma 8.5 would make it a closed analytic subset
of \(Z'\): use submanifold coordinates at its image points and its open
complement elsewhere. The inverse image \((q'\beta)^{-1}(S_x)\) would be
a closed analytic subset of \(\mathbb P^1\) containing the unit circle,
and therefore all of \(\mathbb P^1\).

But at \([1:0]\) the final phase is \((y,[0,v_+])\), with
\(v_+=a(e_1-ie_2)\). Its original geodesic is pure second-factor null
and keeps its first base point equal to \(y_1\ne x_1\). It cannot meet
\(F_x\) even in the full product extension, and neither can its
restricted leaf. Thus \(q'\beta([1:0])\notin S_x\), a contradiction.

The group \(SO_0(1,2)\times SO(3)\) acts transitively on \(M\). Each of
its elements extends by the same real matrices to a holomorphic isometry
of \(B_0\) commuting with conjugation. For an arbitrary real point
\(x\), choose such an isometry \(A\) carrying the distinguished point
\(x_*=((0,a,0),ae_3)\) to \(x\). The neighborhood \(A^{-1}(B')\) is
again allowed. The induced phase isomorphism preserves the restricted
null-leaf relation and identifies its contact quotient with \(Z'\),
carrying the sky at \(x_*\) to \(S_x\). The distinguished-point argument
applied to \(A^{-1}(B')\) therefore proves nonclosedness for \(S_x\).
Finally, \(q'\) is a quotient map, so closedness of a subset of \(Z'\)
is equivalent to closedness of its inverse image. ∎

\pandocstatementheading{\protect\phantomsection\label{lem:finite-module-closed-support}Lemma 8.7.}

The nonzero-stalk support of a locally finitely generated module over a
sheaf of rings is closed. The same holds chartwise for an algebroid
module and after inverse image to a presentation atlas.

\emph{Proof.} At a point where the stalk vanishes, choose finitely many
local generators. Each vanishes on a smaller neighborhood, and on their
finite intersection the module is zero. Its zero-stalk locus is open.
Algebroid comparisons are equivalences and preserve zero modules, so
supports agree on overlaps. Inverse image preserves finite local
generation and has the original stalk at each point above a base point.
∎

\pandocstatementheading{\protect\phantomsection\label{cor:nariai-cofinal-point-module-obstruction-v09}Corollary 8.8.}

For every neighborhood \(B'\) in Proposition 8.4 and every real point
\(x\in M\), no ordinary locally finitely generated module over an
operator algebroid on \(Z'\) has support exactly \(S_x\). On the
corresponding null monodromy contact-groupoid presentation, no such
module has phase-atlas support exactly \(\operatorname{Sat}_{R'}(F_x)\).
In particular, there is no ordinary coherent simple point module with
that exact incidence support.

Every compatible holomorphic Einstein extension of the entire real
Nariai development has an allowed point neighborhood \(U\) such that
these nonexistence assertions hold on \(U\) and on every further allowed
point restriction, for every real point. These representatives are
cofinal among its neighborhoods of the entire real development. The
conclusion holds for either prescribed real spin structure.

\emph{Proof.} The first assertion contradicts Lemma 8.7 and Theorem 8.6.
For the groupoid version, pull the module and algebroid to the Hausdorff
phase atlas \(F'\). This inverse image of operator germs preserves local
finite generation. Its support must therefore be closed, whereas the
required saturation is not.

Let \((\widetilde B,\widetilde g,\widetilde\sigma)\) be any compatible
extension of the entire real Nariai development. By
\hyperref[lem:whole-real-isometry-complexification]{Lemma 7.3}, applied
to the identity real isometry, there is a biholomorphism from an allowed
neighborhood \(U\subset\widetilde B\) onto an allowed neighborhood in
\(B_0\), preserving the metric, conjugation and the specified real
identification. If a real spin structure is prescribed,
\hyperref[prop:nariai-restricted-quotient-v08]{Proposition 8.4} first
supplies a product neighborhood carrying its extension, and the lemma
retains the specified spin identification after a further restriction.
Every allowed neighborhood inside \(U\) is identified in the same way
with a product neighborhood. The induced phase maps identify the
recomputed monodromy presentations and their point-sky saturations, so
the first part applies to every real point on every such neighborhood.
Given any other allowed neighborhood \(W\subset\widetilde B\), the
component containing all of connected \(M\) of \(U\cap W\) is an allowed
neighborhood contained in both. This proves cofinality. ∎

For nested allowed product neighborhoods \(B''\subset B'\), the
recomputed saturations satisfy
\begin{equation}\protect\phantomsection\label{eq:nariai-recomputed-saturation-inclusion-v09}{
\operatorname{Sat}_{R''}(F_x)
\subseteq
\operatorname{Sat}_{R'}(F_x)\cap F'',
\qquad F''=F|_{B''},
}\end{equation} and the inclusion need not be equality: a joining null
path in \(B'\) may leave \(B''\). Ordinary restriction of a module on a
larger extension therefore need not have the recomputed exact support.
The corollary concerns the stated cofinal representatives; it neither
excludes isolated modules on arbitrary larger complexifications nor
asserts a neutralization obstruction for the operator algebroid. Local
whole-sky simple modules remain available, and their parametrized
realization on the Hausdorff incidence is the construction of
\hyperref[sec:operator-parametrized-incidence-v08]{§4.3}.

\protect\phantomsection\label{sec:analytic-cauchy-appendix-v08}

\section{Appendix A. The analytic Cauchy
category}\label{appendix-a.-the-analytic-cauchy-category}

Fix a connected compact oriented real analytic three-manifold \(\Sigma\)
without boundary, an analytic positive definite metric \(h\), an
analytic symmetric tensor \(K\), and a specified spatial spin structure
\(P_\Sigma\). Let \(\Lambda\in\mathbb R\) be fixed. The data satisfy the
Einstein constraints (\ref{eq:ordinary-analytic-constraints-v08}),
namely
\(\operatorname{Scal}(h)+(\operatorname{tr}_hK)^2-|K|_h^2=-2\Lambda\)
and \(\operatorname{div}_hK-\mathrm d\operatorname{tr}_hK=0\). We use
signature \((+---)\), \(\operatorname{Ric}(g)=\Lambda g\), and
\(K(v,w)=-g(\nabla_v\nu,w)\) for the future unit normal \(\nu\).

The required Cauchy input concerns these analytic constraint data. We
give the analytic existence and maximality argument explicitly,
including the relative spin identification used by the source
correspondence.

\pandocstatementheading{\protect\phantomsection\label{def:analytic-ghd-v2}Definition A.1.}

A smooth globally hyperbolic development of \((\Sigma,h,K)\) is a
connected oriented and time-oriented smooth Lorentzian four-manifold
\((M,g)\) satisfying \(\operatorname{Ric}(g)=\Lambda g\), with a
spacelike embedding \(\iota:\Sigma\to M\) inducing \((h,K)\) and having
Cauchy image. Here every inextendible causal curve must meet the image
exactly once. The spacetime orientation is the one for which
\((\nu_\iota,\mathrm d\iota(v_1),\mathrm d\iota(v_2),\mathrm d\iota(v_3))\)
is positive whenever \((v_1,v_2,v_3)\) is positive on \(\Sigma\). An
analytic development has an analytic manifold structure, metric and
initial embedding. A development morphism is an orientation- and
time-orientation-preserving isometric embedding \(f\) satisfying
\(f\circ\iota=\iota'\); it is analytic in the analytic category.

A spin development also has a Lorentz spin structure with a specified
identification of \(P_\Sigma\) with its reduction along the initial
hypersurface to frames adapted to \(\nu_\iota\). This identification
covers
\begin{equation}\protect\phantomsection\label{eq:initial-adapted-frame-identification}{
(e_1,e_2,e_3)\longmapsto
(\nu_\iota,\mathrm d\iota(e_1),\mathrm d\iota(e_2),\mathrm d\iota(e_3)).
}\end{equation} Spin morphisms include a spin lift preserving that
identification. The resulting analytic spin category is denoted by
\(\mathsf{Ein}^{\omega}_\Lambda(\Sigma,h,K)\), with the orientation and
\(P_\Sigma\) understood. A maximal development means a terminal object:
it receives a unique morphism from every development in the category.

\pandocstatementheading{\protect\phantomsection\label{prop:ordinary-analytic-cauchy-collar-v08}Proposition A.2.}

The specified data admit an analytic spin globally hyperbolic
development on \((-\varepsilon,\varepsilon)\times\Sigma\), with initial
embedding \(x\mapsto(0,x)\) and metric \(g=\mathrm dt^2-h_t\).

\emph{Proof.} Put \(\lambda=-\Lambda\) and solve the Gaussian evolution
equations
\begin{equation}\protect\phantomsection\label{eq:ordinary-analytic-gaussian-evolution-v08}{\begin{gathered}
\partial_t h_t=2K_t, \\
\partial_tK_t=-\operatorname{Ric}(h_t)
 -(\operatorname{tr}_{h_t}K_t)K_t
 +2K_th_t^{-1}K_t+\lambda h_t, \\
(h_0,K_0)=(h,K).
\end{gathered}}\end{equation} Eliminating \(K_t\) gives a system solved for
\(\partial_t^2h_t\). Its right-hand side is analytic in the invertible
metric, its first time derivative, and tangential derivatives of order
at most two; the coefficient of the highest normal derivative is the
identity. The analytic Cauchy theorem in normal form therefore gives a
unique local analytic solution in every spatial chart. This is the same
normal form and constraint-propagation calculation used in Appendix A.1
of the local paper, now applied directly to the prescribed constraint
data.

For completeness, set \(\widehat g=-\mathrm dt^2+h_t\) and
\(\tau=\operatorname{tr}_{h_t}K_t\). Its time-containing Christoffel
symbols are \(\widehat\Gamma^0_{ij}=(K_t)_{ij}\) and
\(\widehat\Gamma^i_{0j}=(K_t)^i{}_j\). Consequently \[
\begin{aligned}
\operatorname{Ric}(\widehat g)_{ij}
 &=\operatorname{Ric}(h_t)_{ij}+\partial_t(K_t)_{ij}
   +\tau(K_t)_{ij}-2(K_t)_i{}^q(K_t)_{qj},\\
\operatorname{Ric}(\widehat g)_{00}
 &=-\partial_t\tau-|K_t|_{h_t}^2,
\qquad
\operatorname{Ric}(\widehat g)_{0i}
 =\nabla^j(K_t)_{ij}-\partial_i\tau.
\end{aligned}
\] Thus the spatial components of
\(E=\operatorname{Ric}(\widehat g)-\lambda\widehat g\) vanish.
Contracted Bianchi, with \(H=E_{00}\) and \(M_i=E_{0i}\), gives
\begin{equation}\protect\phantomsection\label{eq:ordinary-analytic-constraint-propagation-v08}{
\partial_tH=2\nabla^iM_i-2\tau H,
\qquad
\partial_tM_i=\tfrac12\partial_iH-\tau M_i.
}\end{equation} Tracing the evolution equations shows that initially
\(H=\operatorname{Scal}(h)+(\operatorname{tr}_hK)^2-|K|_h^2-2\lambda=0\);
the momentum constraint gives \(M=0\). Uniqueness for this homogeneous
analytic Cauchy system gives \(E=0\). Constant reversal of metric sign
leaves the Levi-Civita connection and covariant Ricci tensor unchanged,
so \(g=-\widehat g\) satisfies \(\operatorname{Ric}(g)=\Lambda g\).

Time-independent analytic spatial coordinate changes preserve the
evolution equations and initial data. Analytic uniqueness therefore
identifies the local solutions on overlaps near the initial
hypersurface. Choose a finite cover by smaller spatial charts with
closures in their construction charts. Compactness gives a common
positive time interval on which all solutions and overlap
identifications exist. They glue on the fixed product with \(\Sigma\);
uniqueness under complex conjugation gives real coefficients. After
decreasing the interval, \(h_t\) is positive definite and extends to its
closed time interval with uniform upper and lower bounds relative to
\(h\).

The function \(t\) is temporal. A future-directed causal curve,
parametrized by \(t\), satisfies
\(h_t(\mathrm dx/\mathrm dt,\mathrm dx/\mathrm dt)\le1\) almost
everywhere. Uniform comparison with \(h\) bounds its spatial speed. If
its time range had a finite endpoint inside the interval, its spatial
component would be Cauchy as that endpoint is approached, and
completeness of the compact Riemannian manifold \((\Sigma,h)\) would
give a limit. The curve could then be continued through that spacetime
point by a vertical timelike segment. Therefore an inextendible causal
curve has the entire time interval as its time range and meets every
slice once. In particular the initial slice is Cauchy.

The prescribed orientations extend over the product.
\hyperref[lem:relative-spin-extension-v2]{Lemma A.6} extends
\(P_\Sigma\) to this development with its required initial
identification and analytic covering structure. ∎

\pandocstatementheading{\protect\phantomsection\label{thm:smooth-mghd-input}Theorem A.3 (Smooth Cauchy theorem).}

Smooth Einstein constraint data for fixed \(\Lambda\) admit a smooth
globally hyperbolic development
\((M_{\mathrm{sm}},g_{\mathrm{sm}},\iota_{\mathrm{sm}})\) receiving a
unique initial-data- and time-orientation-preserving isometric embedding
from every smooth globally hyperbolic development of those data.

For \(\Lambda=0\) this is the maximal-development result of Sbierski
(2016, Theorems 2.6 and 2.8, Corollary 3.2); the embedding and
common-development arguments are his Lemma 3.3 and Theorem 3.4. The
hypotheses are connected smooth developments and smooth constraint data.
Before Definition 2.1 he notes that the same proof applies with a
cosmological constant once local existence and uniqueness hold. In the
harmonic reduction, \(\Lambda g\) is a smooth zeroth-order term, so the
reduced equations have the same quasilinear wave principal part and the
same local existence and uniqueness argument. The contracted Bianchi
identity gives homogeneous gauge and constraint propagation because
\(\nabla g=0\) and \(\mathrm d\Lambda=0\); the constraint calculation in
Proposition A.2 is its Gaussian form. These observations give the local
input needed for the same maximal-development construction. Signature
reversal changes the Einstein coefficient and preserves the argument.
The required orientation extends from the initial hypersurface by the
product decomposition proved below. An initial-data- and
time-orientation-preserving embedding preserves that orientation on the
initial hypersurface, hence everywhere on its connected domain.

\pandocstatementheading{\protect\phantomsection\label{lem:analytic-isometries}Lemma A.4.}

A smooth local isometry between real analytic pseudo-Riemannian
manifolds of the same dimension is analytic. A fixed smooth metric
consequently has at most one compatible maximal analytic atlas in which
it is analytic.

\emph{Proof.} Analytic dependence for the geodesic equation makes the
exponential maps analytic near their zero vectors, with analytic local
inverses. Preservation of the Levi-Civita connection gives
\begin{equation}\protect\phantomsection\label{eq:analytic-isometry-exponential}{
f\circ\exp_p^{g_N}
 =\exp_{f(p)}^{g_M}\circ\mathrm df_p.
}\end{equation} This formula proves analyticity near each \(p\). Apply
it to the identity between two proposed analytic structures, in both
directions, to obtain compatibility of their maximal atlases. ∎

\pandocstatementheading{\protect\phantomsection\label{lem:cauchy-union}Lemma A.5.}

Let \((M,g)\) be globally hyperbolic with connected Cauchy hypersurface
\(S\). If a nonempty set of open subsets \(U_\alpha\subset M\) all
contain \(S\) and have \(S\) as a Cauchy hypersurface, then \(S\) is
Cauchy in \(U=\bigcup_\alpha U_\alpha\). The union is a connected
globally hyperbolic open subspacetime.

\emph{Proof.} For an inextendible causal curve \(\gamma:I\to U\), choose
a parameter value whose image belongs to some \(U_\alpha\), and take the
component \(J\) of \(\gamma^{-1}(U_\alpha)\) containing it. The
restriction \(\gamma|_J\) is inextendible in \(U_\alpha\): an endpoint
of \(J\) inside \(I\) has image outside \(U_\alpha\), whereas an
extension at an endpoint of \(I\) would extend \(\gamma\) in \(U\). Thus
\(\gamma|_J\) meets \(S\). Two intersections of \(\gamma\) with \(S\)
would give an ambient causal segment that extends to an ambient
inextendible causal curve meeting \(S\) twice, contrary to the ambient
Cauchy property. Each component of each \(U_\alpha\) meets \(S\), by
applying the same argument to an inextendible timelike curve.
Connectedness of \(S\) implies connectedness of the union. The existence
of its Cauchy hypersurface gives global hyperbolicity. ∎

This is a union inside a fixed ambient development, as in the causal
argument in Sbierski (2016, Theorem 3.4). It requires no uniform
analytic radius.

\pandocstatementheading{\protect\phantomsection\label{lem:relative-spin-extension-v2}Lemma A.6.}

The prescribed \(P_\Sigma\) extends to every smooth development,
uniquely up to a unique isomorphism with prescribed initial
identification. On an analytic development the extension has a unique
analytic covering structure. Every analytic development morphism has a
unique analytic spin lift extending a prescribed compatible initial
lift.

\emph{Proof.} Choose a future-directed timelike vector field and rescale
it to have bounded norm in a complete auxiliary Riemannian metric. Its
flow is complete, and each of its inextendible timelike orbits meets the
initial Cauchy hypersurface exactly once. The resulting map
\(\mathbb R\times\Sigma\to M\) is bijective and a local diffeomorphism,
since the field is transverse to the initial hypersurface. It is
therefore a diffeomorphism and exhibits that hypersurface as a
deformation retract.

Choose a smooth future unit timelike field \(T\) agreeing with the
initial normal. On \(E=T^\perp\), endowed with \(-g|_E\), parallel
transport for a metric connection along the product lines identifies
\(E\) with the pullback of \((T\Sigma,h)\), relative to the identity on
the initial slice. Pull back \(P_\Sigma\) and extend its structure group
to \(\operatorname{Spin}^{+}(1,3)\). This constructs an extension
without requiring an analytic product decomposition.

The two spin structures in a comparison differ by a principal
\(\mathbb Z_2\)-bundle on \(M\). The prescribed initial identification
is a section of that difference bundle on the deformation retract.
Covering-space lifting extends the section over \(M\); two extensions
differ by a locally constant sign fixed on the initial hypersurface, so
they coincide on connected \(M\). This proves relative uniqueness.

For an analytic metric the orthonormal frame bundle is analytic. Its
spin cover has the unique analytic atlas making the covering projection
a local analytic diffeomorphism. The spin action and initial
identification are analytic by local lifting. For a morphism, pull back
the target spin structure and apply relative uniqueness on the source.
The frame map is analytic, hence so is its covering lift. Uniqueness
proves the composition law. Initial lifts differing by the central sign
have distinct extensions. ∎

\pandocstatementheading{\protect\phantomsection\label{thm:maximal-analytic-development-v2}Theorem A.7 (Maximal analytic development).}

Assume that the analytic initial data admit an analytic globally
hyperbolic development. There is an analytic spin development
\((M_\omega,g_\omega,\iota_\omega,P_\omega)\) receiving a unique
analytic spin morphism from every analytic spin development of the same
oriented data and specified initial spin structure. It is unique up to
the unique analytic spin isometry preserving the initial identification.
Its underlying smooth spacetime is an open subspacetime of the smooth
maximal development.

\emph{Proof.} Fix the smooth maximal development and put
\(S=\iota_{\mathrm{sm}}(\Sigma)\). Let \(\mathscr U_\omega\) consist of
the open subsets \(U\subset M_{\mathrm{sm}}\) containing \(S\), with
\(S\) Cauchy in \(U\), that admit a compatible analytic structure making
the restricted metric and initial embedding analytic. This is a set of
subsets of a fixed manifold; \hyperref[lem:analytic-isometries]{Lemma
A.4} gives at most one compatible maximal analytic atlas on each member.
It is nonempty: the assumed analytic development embeds smoothly in
\(M_{\mathrm{sm}}\), and its analytic atlas pushes forward to its open
image.

Set
\begin{equation}\protect\phantomsection\label{eq:maximal-analytic-union}{
M_\omega=\bigcup_{U\in\mathscr U_\omega}U,
\qquad g_\omega=g_{\mathrm{sm}}|_{M_\omega},
\qquad \iota_\omega=\iota_{\mathrm{sm}}:\Sigma\longrightarrow M_\omega.
}\end{equation} This open subspace is Hausdorff and second countable. On
the intersection of two members the identity is a smooth local isometry,
so \hyperref[lem:analytic-isometries]{Lemma A.4} makes their analytic
atlases compatible. The union atlas, with a countable subatlas if
desired, makes the metric and initial embedding analytic.
\hyperref[lem:cauchy-union]{Lemma A.5} makes \(S\) Cauchy in
\(M_\omega\).

Every analytic development \(N\) has a unique smooth
initial-data-preserving embedding into \(M_{\mathrm{sm}}\). Its open
image, with its pushed-forward analytic atlas, belongs to
\(\mathscr U_\omega\); the embedding therefore factors analytically
through \(M_\omega\). Any second factor would give a second such
embedding into \(M_{\mathrm{sm}}\), contrary to the smooth theorem.
\hyperref[lem:relative-spin-extension-v2]{Lemma A.6} extends the
prescribed spin structure to \(M_\omega\) and uniquely lifts each factor
with its initial identification. Two objects with this universal
property admit unique morphisms in both directions, whose composites are
identities. This also removes dependence on the chosen smooth maximal
development. ∎

\pandocstatementheading{\protect\phantomsection\label{cor:analytic-cauchy-naturality-v08}Corollary A.8.}

The analytic constraint data fixed above have a maximal analytic spin
development, and their analytic Cauchy collar has a unique spin
embedding into it. Every analytic isomorphism of such initial data with
a specified spatial spin lift extends uniquely to an analytic spin
isometry between their maximal analytic developments. These extensions
preserve identities and composition.

\emph{Proof.}
\hyperref[prop:ordinary-analytic-cauchy-collar-v08]{Proposition A.2}
supplies the existence hypothesis of
\hyperref[thm:maximal-analytic-development-v2]{Theorem A.7}. For an
initial-data isomorphism, reparametrize the target initial embedding and
transport its initial spin identification by the specified lift. The
universal property gives a morphism for these identified data. Apply it
also to the inverse data isomorphism. The two composites are identities
by uniqueness; the same argument proves the identity and composition
laws. ∎

\pandocstatementheading{\protect\phantomsection\label{rem:analytic-maximality-scope}Remark A.9.}

Maximality here is within the analytic globally hyperbolic category of
the fixed Cauchy data. The construction proves
\(M_\omega\subset M_{\mathrm{sm}}\) as an open subspacetime. Equality
would require a separate analytic-regularity propagation argument.
Neither geodesic completeness nor maximality among extensions in which
the initial hypersurface ceases to be Cauchy is asserted.

\protect\phantomsection\label{sec:references}

\section{References}\label{references}

\protect\phantomsection\label{ref:AdamoAranedaSeetSharma2026}

Adamo, Tim, Bernardo Araneda, Sean Seet, and Atul Sharma. 2026.
\emph{Schwarzschild black holes from twistor space}.
\href{https://arxiv.org/abs/2607.06236}{arXiv:2607.06236}, July 7.

\protect\phantomsection\label{ref:BastonMason1987}

Baston, Robert J., and Lionel J. Mason. 1987. ``Conformal gravity, the
Einstein equations and spaces of complex null geodesics.''
\emph{Classical and Quantum Gravity} \textbf{4} (4): 815--826.
\href{https://doi.org/10.1088/0264-9381/4/4/018}{doi:10.1088/0264-9381/4/4/018}.

\protect\phantomsection\label{ref:BautistaIbortLafuente2015}

Bautista, A., A. Ibort, and J. Lafuente. 2015. ``Causality and skies: is
non-refocussing necessary?'' \emph{Classical and Quantum Gravity}
\textbf{32} (10): 105002.
\href{https://doi.org/10.1088/0264-9381/32/10/105002}{doi:10.1088/0264-9381/32/10/105002}.

\protect\phantomsection\label{ref:Beyer2009}

Beyer, Florian. 2009. ``Non-genericity of the Nariai solutions: I.
Asymptotics and spatially homogeneous perturbations.'' \emph{Classical
and Quantum Gravity} \textbf{26} (23): 235015.
\href{https://doi.org/10.1088/0264-9381/26/23/235015}{doi:10.1088/0264-9381/26/23/235015}.
Section 3.2.1 of
\href{https://arxiv.org/abs/0902.2531v2}{arXiv:0902.2531v2} records the
real Nariai metric, compact Cauchy hypersurfaces and completeness; the
metric and curvature signs are fixed in §2 and verified in §8.

\protect\phantomsection\label{ref:Chernov2018}

Chernov, Vladimir. 2018. \emph{Causality and Legendrian linking for
higher dimensional spacetimes}.
\href{https://arxiv.org/abs/1803.04590v2}{arXiv:1803.04590v2}, June 8.
The classical real contact-ray identification is recalled in §1.

\protect\phantomsection\label{ref:CrainicFernandes2003}

Crainic, Marius, and Rui Loja Fernandes. 2003. ``Integrability of Lie
brackets.'' \emph{Annals of Mathematics} \textbf{157} (2): 575--620.
\href{https://doi.org/10.4007/annals.2003.157.575}{\nolinkurl{doi:10.4007/annals.2003.157.575}}.
The source-simply-connected integration and Lie functor statements are
recalled on pp.~577--579; Proposition 1.6 gives the homotopy invariance
of representation transport.

\protect\phantomsection\label{ref:IsthmusTwistor}

Isthmus, Lior. 2026a. \emph{A projective twistor correspondence for
local analytic Einstein germs}. Hexagon preprint \textbf{2611.00123v1},
November 5. \href{https://hexagonmath.org/2611.00123v1}{Posted version
v1}.

\protect\phantomsection\label{ref:IsthmusGlobalComputations}

Isthmus, Lior. 2026b. \emph{Computational supplement for Global twistor
sources for analytic Einstein spacetimes}. Version 1.0.0. Computational
files accompanying version 1.0 of this paper, in the
\texttt{computations/} directory of
\href{https://github.com/lioristhmus/global-twistor-sources/releases/tag/v1.0}{GitHub
release v1.0}.

\protect\phantomsection\label{ref:Kashiwara1996}

Kashiwara, Masaki. 1996. ``Quantization of Contact Manifolds.''
\emph{Publications of the Research Institute for Mathematical Sciences}
\textbf{32} (1): 1--7.
\href{https://doi.org/10.2977/PRIMS/1195163179}{doi:10.2977/PRIMS/1195163179}.
Theorem 2 gives the canonical stack of microdifferential modules; the
paragraph following it states the invertible-object criterion for a
global algebra presentation. The contact-groupoid descent in §4 is
proved here from the structured local quantization input.

\protect\phantomsection\label{ref:KashiwaraSchapira2012}

Kashiwara, Masaki, and Pierre Schapira. 2012. \emph{Deformation
quantization modules}. Astérisque \textbf{345}. Société Mathématique de
France.
\href{https://www.numdam.org/item/AST_2012__345__R1_0/}{Published
volume}. Section 2.1 gives algebroid descent and its module
interpretation; Proposition 2.1.13 gives the bimodule gluing
construction without a paracompactness hypothesis. The numbering cited
here is that of
\href{https://arxiv.org/abs/1003.3304v3}{arXiv:1003.3304v3}.

\protect\phantomsection\label{ref:LaurentGengouxStienonXu2009}

Laurent-Gengoux, Camille, Mathieu Stiénon, and Ping Xu. 2009.
``Integration of Holomorphic Lie Algebroids.'' \emph{Mathematische
Annalen} \textbf{345} (4): 895--923.
\href{https://doi.org/10.1007/s00208-009-0388-7}{doi:10.1007/s00208-009-0388-7}.
The holomorphic integration statement used in §3 is Theorem 3.17 of
\href{https://arxiv.org/abs/0803.2031v2}{arXiv:0803.2031v2}.

\protect\phantomsection\label{ref:LeBrun1983}

LeBrun, Claude. 1983. ``Spaces of complex null geodesics in
complex-Riemannian geometry.'' \emph{Transactions of the American
Mathematical Society} \textbf{278} (1): 209--231.
\href{https://doi.org/10.1090/S0002-9947-1983-0697071-9}{doi:10.1090/S0002-9947-1983-0697071-9}.

\protect\phantomsection\label{ref:LeBrun1985}

LeBrun, Claude. 1985. ``Ambi-twistors and Einstein's equations.''
\emph{Classical and Quantum Gravity} \textbf{2} (4): 555--563.
\href{https://doi.org/10.1088/0264-9381/2/4/020}{doi:10.1088/0264-9381/2/4/020}.

\protect\phantomsection\label{ref:LeBrun1991}

LeBrun, Claude. 1991. ``Thickenings and conformal gravity.''
\emph{Communications in Mathematical Physics} \textbf{139} (1): 1--43.
\href{https://doi.org/10.1007/BF02102727}{doi:10.1007/BF02102727}.

\protect\phantomsection\label{ref:LeBrunMason2007}

LeBrun, Claude, and L. J. Mason. 2007. ``Nonlinear gravitons, null
geodesics, and holomorphic disks.'' \emph{Duke Mathematical Journal}
\textbf{136} (2): 205--273.
\href{https://doi.org/10.1215/S0012-7094-07-13621-4}{doi:10.1215/S0012-7094-07-13621-4}.

\protect\phantomsection\label{ref:MarcolliPenrose2021}

Marcolli, Matilde, and Roger Penrose. 2021. ``Gluing Non-commutative
Twistor Spaces.'' \emph{The Quarterly Journal of Mathematics}
\textbf{72} (1--2): 417--454.
\href{https://doi.org/10.1093/qmath/haab024}{doi:10.1093/qmath/haab024}.

\protect\phantomsection\label{ref:Penrose1976}

Penrose, Roger. 1976. ``Nonlinear gravitons and curved twistor theory.''
\emph{General Relativity and Gravitation} \textbf{7}: 31--52.
\href{https://doi.org/10.1007/BF00762011}{doi:10.1007/BF00762011}.

\protect\phantomsection\label{ref:Penrose2015}

Penrose, Roger. 2015. ``Palatial twistor theory and the twistor googly
problem.'' \emph{Philosophical Transactions of the Royal Society A}
\textbf{373} (2047): 20140237.
\href{https://doi.org/10.1098/rsta.2014.0237}{doi:10.1098/rsta.2014.0237}.

\protect\phantomsection\label{ref:Penrose2026}

Penrose, Roger. 2026. ``Bi-twistors: An extension of twistor theory for
general spacetimes.'' \emph{International Journal of Modern Physics D}
\textbf{35} (6): 2540015. Published online May 6, 2025.
\href{https://doi.org/10.1142/S0218271825400152}{doi:10.1142/S0218271825400152}.

\protect\phantomsection\label{ref:Polesello2005}

Polesello, Pietro. 2005. \emph{Uniqueness of quantization of complex
contact manifolds}.
\href{https://arxiv.org/abs/math/0512501}{arXiv:math/0512501}, December
21. The distinction between a global algebra and its algebroid is
recalled in the Introduction; Theorem 3.3 is the uniqueness theorem.

\protect\phantomsection\label{ref:Sbierski2016}

Sbierski, Jan.~2016. ``On the existence of a maximal Cauchy development
for the Einstein equations: a dezornification.'' \emph{Annales Henri
Poincaré} \textbf{17} (2): 301--329.
\href{https://doi.org/10.1007/s00023-015-0401-5}{doi:10.1007/s00023-015-0401-5}.
The theorem numbering cited in the text is that of
\href{https://arxiv.org/abs/1309.7591v3}{arXiv:1309.7591v3}.

\protect\phantomsection\label{ref:Sharma2021}

Sharma, Atul. 2021. \emph{Twistor action for general relativity}.
\href{https://arxiv.org/abs/2104.07031v2}{\nolinkurl{arXiv:2104.07031v2}}, July 28.

\protect\phantomsection\label{ref:StacksProject}

The Stacks project authors. 2026. \emph{The Stacks Project}.
\href{https://stacks.math.columbia.edu/tag/02ZC}{Tag 02ZC, ``Descent
data in fibred categories''}, and
\href{https://stacks.math.columbia.edu/tag/0268}{Tag 0268, ``Stacks''}.
These sections provide the descent-data, refinement and effectivity
terminology; the source-atlas construction and its analytic comparisons
are proved in §5.

\protect\phantomsection\label{ref:WhitneyBruhat1959}

Whitney, Hassler, and François Bruhat. 1959. ``Quelques propriétés
fondamentales des ensembles analytiques-réels.'' \emph{Commentarii
Mathematici Helvetici} \textbf{33}: 132--160.
\href{https://doi.org/10.1007/BF02565913}{doi:10.1007/BF02565913}. The
complexification and germ uniqueness theorem is §1, Proposition 1; the
compatible involution with exactly the given real fixed locus is the
following remark, p.~136.

\protect\phantomsection\label{ref:WoodhouseMason1988}

Woodhouse, N. M. J., and L. J. Mason. 1988. ``The Geroch group and
non-Hausdorff twistor spaces.'' \emph{Nonlinearity} \textbf{1} (1):
73--114.
\href{https://doi.org/10.1088/0951-7715/1/1/004}{doi:10.1088/0951-7715/1/1/004}.

\end{document}
